\documentclass[12pt,a4paper]{article}

\usepackage[utf8]{inputenc}
\usepackage[T1]{fontenc}
\usepackage{lmodern}
\usepackage{amsmath,amssymb,amsthm}
\usepackage{enumitem}
\usepackage{xcolor}
\usepackage{geometry}
\geometry{a4paper, margin=2.5cm}
\usepackage{hyperref}

\hypersetup{
    colorlinks=true,
    linkcolor=blue!70!black,
    citecolor=green!50!black,
    urlcolor=blue!70!black
}

\theoremstyle{definition}
\newtheorem{definition}{Definition}[section]
\theoremstyle{plain}
\newtheorem{theorem}[definition]{Satz}
\newtheorem{lemma}[definition]{Lemma}
\newtheorem{proposition}[definition]{Proposition}
\newtheorem{korollar}[definition]{Korollar}
\newtheorem{axiom}[definition]{Axiom}
\theoremstyle{remark}
\newtheorem{remark}[definition]{Bemerkung}

\renewcommand{\proofname}{Beweis}

\title{\textbf{A Thermodynamic Obstruction to Finite-Time Blow-Up\\ in the 3D Navier-Stokes Equations (Conditional Framework)}}
\author{Lukas Geiger\\
\small Unabhaengiger Forscher, Bernau im Schwarzwald, Deutschland\\
\small ORCID: \href{https://orcid.org/0009-0005-7296-1534}{0009-0005-7296-1534}%
\thanks{KI-Werkzeuge (Claude von Anthropic) wurden f\"ur mathematische Formalisierung, Literaturverifikation und redaktionelle Verfeinerung eingesetzt. Alle mathematischen Inhalte wurden vom Autor entwickelt, verifiziert und freigegeben.}}
\date{Maerz 2026}

\begin{document}
\maketitle

\begin{abstract}
Wir schlagen ein doppelt bedingtes Regularitaetskriterium fuer die inkompressiblen 3D-Navier-Stokes-Gleichungen vor, das auf einem thermodynamischen Abstandsfunktional basiert. Anstatt den Geschwindigkeitsgradienten global zu beschraenken, konstruieren wir ein Attraktor-Abstandsfunktional $H(u \mid \mathcal{A})$, das den $L^2$-Abstand der aktiven Stroemung vom globalen Leray--Hopf-Attraktor misst. Unter zwei Bedingungen -- (\textbf{Assumption~G}) dass die zeitabhaengige minimierende Projektion auf den Attraktor hinreichend regulaer ist, und (\textbf{Condition~D}) dass die kumulative viskose Dissipation die Gronwall-Kosten nahe einer hypothetischen Blow-up-Zeit dominiert -- zeigen wir, dass ein Blow-up in endlicher Zeit $H$ unter Null zwingen wuerde, ein Widerspruch zur Nichtnegativitaet. Das Argument reduziert die Regularitaetsfrage auf zwei strukturelle Eigenschaften: die Geometrie des Attraktors (Assumption~G) und das Verhaeltnis von Dissipation zu nichtlinearem Wachstum (Condition~D). Wir identifizieren hinreichende geometrische Bedingungen fuer~G (positiver Reach, log-Lipschitz-Projektionen, BV-Selektionen) und praesentieren einen $H^1$-Level-Lift, bei dem die Enstrophie-Dissipation -- die im Gegensatz zur Energie-Dissipation keine a-priori-Oberschranke unter Blow-up besitzt -- als divergenter Term dienen kann. Wir diskutieren den Status dieser Annahmen, ihre Verbindung zur Theorie inertialer Mannigfaltigkeiten und die Luecken, die das vorliegende Rahmenwerk von einem unbedingten Beweis trennen.
Diese Arbeit etabliert eine \emph{bedingte Reduktion}: Globale Regularitaet ist aequivalent zur Existenz einer BV-regulaeren Selektion auf dem Attraktor (Assumption~G), die in der Begleitarbeit zur Log-Distanz-Integrierbarkeit \cite{Geiger2026-LDI} behandelt wird.
\end{abstract}

\section{Einleitung}
Die klassische Herausforderung bezueglich der inkompressiblen 3D-Navier-Stokes-Gleichungen besteht darin, die Entstehung von Singularitaeten in endlicher Zeit (Blow-ups) im Geschwindigkeitsfeld $u(x,t) \in \mathbb{R}^3$ auszuschliessen -- ein Problem, das seit Lerays grundlegender Arbeit \cite{Leray1934} offen geblieben ist. Die Stroemung wird beschrieben durch:
\begin{equation}
    \partial_t u + (u \cdot \nabla) u = -\nabla p + \nu \Delta u, \quad \nabla \cdot u = 0, \quad \text{auf } \Omega \subset \mathbb{R}^3.
\end{equation}
Das bekannte Beale--Kato--Majda-Kriterium (BKM) \cite{BealeKatoMajda1984} besagt, dass eine Singularitaet zur Zeit $T^*$ genau dann auftritt, wenn die Vortizitaet divergiert:
\begin{equation}
    \int_0^{T^*} \|\nabla \times u(t)\|_{L^\infty} \, dt = \infty.
\end{equation}
In dieser Arbeit konstruieren wir ein thermodynamisches Ljapunow-Funktional, das den Abstand der aktiven Stroemung vom globalen Leray--Hopf-Attraktor misst. Wir zeigen, dass -- \emph{bedingt durch eine Regularitaetsannahme ueber die Attraktorprojektion} (Assumption~G, Abschnitt~3) und eine \emph{Dissipations-Dominanz-Bedingung}~(D) -- jeder hypothetische Blow-up zu einem Widerspruch fuehrt: Die Enstrophie divergiert punktweise und treibt das nichtnegative Abstandsfunktional unter Null, sobald die kumulative viskose Dissipation die Gronwall-Kosten uebersteigt. Das resultierende doppelt bedingte Regularitaetskriterium reduziert das vollstaendige Regularitaetsproblem auf zwei strukturelle Fragen: die Geometrie des Navier-Stokes-Attraktors (Assumption~G) und die relativen Blow-up-Raten von Dissipation gegenueber Gronwall-Faktoren (Condition~D).

\section{Das Energie-Dissipations-Divergenzfunktional}
Wir definieren ein verallgemeinertes Divergenzfunktional $H(u \mid \mathcal{A})$, das den Abstand des aktiven Geschwindigkeitsfeldes $u$ vom globalen Attraktor $\mathcal{A}$ darstellt.

\begin{definition}[Divergenz vom globalen Attraktor]
Sei $\Omega = \mathbb{T}^3$ der flache Torus. Wir definieren die Referenzmannigfaltigkeit $\mathcal{A}$ als den globalen Leray--Hopf-Attraktor des Navier-Stokes-Semiflusses (vgl.\ \cite{FoiasManleyRosaTemam2001,Temam2001}). Fuer jedes aktive Stroemungsfeld $u \in L^2(\Omega)$ ist das Abstandsfunktional:
\begin{equation}
    H(u \mid \mathcal{A}) := \inf_{v \in \mathcal{A}} E_{v}(u) = \inf_{v \in \mathcal{A}} \frac{1}{2} \int_\Omega |u - v|^2 \, dx.
\end{equation}
Da $\mathcal{A}$ eine kompakte Teilmenge von $L^2(\Omega)$ ist, wird das Infimum angenommen. Wir bezeichnen mit $v(t) \in \mathcal{A}$ eine messbare Auswahl, die diesen Abstand zur Zeit $t$ minimiert, sodass $H(u \mid \mathcal{A}) = \frac{1}{2} \int_\Omega |w|^2 \, dx$ mit $w := u - v$.
\end{definition}

\begin{proposition}[Regularitaet des globalen Attraktors]
\label{prop:attractor-regularity}
Sei $\Omega = \mathbb{T}^3$ und sei $f \in L^\infty_t H^1_x(\Omega)$ eine glatte
Antriebskraft. Der globale Leray--Hopf-Attraktor $\mathcal{A}$ erfuellt:
\begin{enumerate}
    \item[\emph{(i)}]   $\mathcal{A}$ ist kompakt in $H^1(\Omega)$.
    \item[\emph{(ii)}]  Jedes $v \in \mathcal{A}$ ist glatt: $v \in C^\infty(\Omega)$.
    \item[\emph{(iii)}] $\sup_{v \in \mathcal{A}} \|v\|_{W^{1,\infty}} < \infty$.
\end{enumerate}
\end{proposition}

\begin{proof}
\noindent\textbf{(i) Kompaktheit in $H^1$.}
Der Navier--Stokes-Semifluss $S(t)$ auf $\mathbb{T}^3$ mit Antrieb
$f \in L^\infty_t H^1_x$ besitzt eine absorbierende Kugel
$\mathcal{B}_1 \subset H^1(\Omega)$: Es existiert $R_1 > 0$ (abhaengig von
$\nu$ und $\|f\|_{L^\infty H^1}$) und $t_1 > 0$, sodass
$S(t)\bigl(\mathcal{B}_0\bigr) \subset \mathcal{B}_1 := \{u : \|u\|_{H^1} \le R_1\}$
fuer alle $t \ge t_1$ und jede beschraenkte Menge $\mathcal{B}_0 \subset L^2(\Omega)$.
Dies folgt aus der Standard-Energieabschaetzung: Testen der Navier--Stokes-Gleichung
mit $-\Delta u$ liefert eine Differentialungleichung fuer $\|\nabla u\|_{L^2}^2$,
die, kombiniert mit der Poincar\'{e}-Ungleichung, gleichmaessige $H^1$-Schranken
nach einer Transiente erzeugt (vgl.\ \cite{Temam2001}, Kapitel~III, \S3;
\cite{FoiasManleyRosaTemam2001}, \S9.2).

Der globale Attraktor ist definiert als
$\mathcal{A} = \bigcap_{t\ge 0} S(t)\bigl(\mathcal{B}_1\bigr)$, und er ist die
maximal kompakte invariante Menge des Semiflusses. Auf $\mathbb{T}^3$ ist der Semifluss
$S(t)$ \emph{kompakt} als Abbildung von $H^1$ nach $H^1$ (er bildet beschraenkte
Mengen in $L^2$ in beschraenkte Mengen in $H^2$ ab fuer $t > 0$, und die Einbettung
$H^2(\mathbb{T}^3) \hookrightarrow H^1(\mathbb{T}^3)$ ist kompakt nach dem
Rellich--Kondrachov-Satz). Nach dem abstrakten Attraktor-Existenzsatz
(\cite{Temam2001}, Theorem~I.1.1; \cite{FoiasManleyRosaTemam2001}, Theorem~9.6)
ist $\mathcal{A}$ kompakt in $H^1(\Omega)$.

\medskip
\noindent\textbf{(ii) Glattheit der Attraktorelemente.}
Sei $v \in \mathcal{A}$. Da $\mathcal{A}$ invariant ist, existiert eine
vollstaendige beschraenkte Trajektorie $\{v(t)\}_{t\in\mathbb{R}} \subset \mathcal{A}$
durch $v = v(0)$. Jedes $v(t)$ erfuellt die Navier--Stokes-Gleichung
\[
    \partial_t v + (v\cdot\nabla)v + \nabla p = \nu\Delta v + f.
\]
Da $v(t) \in H^1(\mathbb{T}^3)$ gleichmaessig und $f \in H^1$, bootstrappen wir
die Regularitaet wie folgt.

\emph{Erster Schritt ($H^1 \to H^2$):} Der nichtlineare Term
$(v\cdot\nabla)v \in L^2(\mathbb{T}^3)$ folgt aus der Sobolev-Einbettung
$H^1(\mathbb{T}^3) \hookrightarrow L^6(\mathbb{T}^3)$ und der H\"older-Ungleichung.
Umschreiben der Gleichung als stationaeres Stokes-System
$-\nu\Delta v + \nabla p = f - (v\cdot\nabla)v - \partial_t v$,
elliptische Regularitaet fuer den Stokes-Operator auf $\mathbb{T}^3$
(\cite{ConstantinFoias1988}, \S2.4; \cite{Temam2001}, \S II.3) liefert
$v(t) \in H^2(\mathbb{T}^3)$.

\emph{Induktionsschritt ($H^k \to H^{k+1}$):} Sei $v \in H^k$ fuer ein
$k \ge 2$. Da $H^k(\mathbb{T}^3)$ eine Banach-Algebra fuer $k \ge 2$ ist
(wegen der Sobolev-Einbettung $H^k \hookrightarrow L^\infty$ fuer $k > 3/2$),
gilt $(v\cdot\nabla)v \in H^{k-1}$. Elliptische Regularitaet fuer
den Stokes-Operator liefert dann $v \in H^{k+1}$. Per Induktion folgt
$v \in \bigcap_{k\ge 1} H^k(\mathbb{T}^3) = C^\infty(\mathbb{T}^3)$.

\medskip
\noindent\textbf{(iii) Gleichmaessige $W^{1,\infty}$-Schranke.}
Aus Teil~(ii) und dem Bootstrap-Argument folgt $\mathcal{A} \subset H^k(\mathbb{T}^3)$
fuer jedes $k \ge 1$. Zudem erstreckt sich das Argument der absorbierenden Kugel auf
hoehere Sobolev-Normen: Fuer jedes $k$ existiert $R_k > 0$ mit
$\sup_{v\in\mathcal{A}}\|v\|_{H^k} \le R_k$ (vgl.\ \cite{FoiasManleyRosaTemam2001},
Theorem~12.1; \cite{Temam2001}, \S III.5).
Mit der Wahl $k > 5/2$ (z.B.\ $k = 3$) liefert der Sobolev-Einbettungssatz auf
$\mathbb{T}^3$ die stetige Einbettung $H^k(\mathbb{T}^3) \hookrightarrow W^{1,\infty}(\mathbb{T}^3)$.
Daher
\[
    \sup_{v\in\mathcal{A}} \|v\|_{W^{1,\infty}}
    \;\le\; C_{\mathrm{Sob}}\,\sup_{v\in\mathcal{A}} \|v\|_{H^k}
    \;\le\; C_{\mathrm{Sob}}\,R_k
    \;<\; \infty,
\]
wobei $C_{\mathrm{Sob}}$ die Einbettungskonstante ist.
\end{proof}

\begin{remark}[Zum Attraktor ohne unbedingten Regularitaetsbeweis]
\label{rem:circularity}
Die Existenz eines kompakten globalen Attraktors fuer den 3D-Navier--Stokes-Semifluss
auf~$\mathbb{T}^3$ ist Standard, wenn der Semifluss ein stetiges dynamisches System
auf~$H^1$ erzeugt
(vgl.~\cite{Temam2001,FoiasManleyRosaTemam2001}). Ohne einen unbedingten
Regularitaetsbeweis ist der Semifluss nur durch Leray--Hopf-schwache Loesungen
definiert, fuer die Eindeutigkeit offen bleibt.
Gleichwohl existiert der \emph{schwache Attraktor} (die $\omega$-Limesmenge der
absorbierenden Kugel unter dem Leray--Hopf-Semifluss; siehe
\cite{FoiasManleyRosaTemam2001}, \S12) und ist kompakt
in~$L^2$, und jedes Element auf dem schwachen Attraktor ist glatt nach dem
Standard-Bootstrap-Argument (Invarianz liefert eine vollstaendige beschraenkte
Trajektorie, die ueber den Stokes-Operator regularisiert). Die vorliegende
Arbeit bewegt sich vollstaendig in diesem Rahmen: Der Attraktor
$\mathcal{A}$ ist der schwache Leray--Hopf-Attraktor, und die Glattheit
seiner Elemente (Proposition~\ref{prop:attractor-regularity}(ii)--(iii))
setzt globale Regularitaet beliebiger Loesungen \emph{nicht} voraus.
\end{remark}

Durch Verfolgen der Entwicklung dieses Funktionals leiten wir eine exakte Dissipationsschranke her.

\begin{lemma}[Energie-Dissipations-Divergenzschranke]
\label{lem:entropy-bound}
Seien $u, v \in L^\infty_t H^1_x(\Omega)$ divergenzfreie Loesungen der Navier-Stokes-Gleichungen mit identischem glattem Antrieb $f$. Die Zeitentwicklung des spezifischen Energieabstands $E_{v}(t) = \frac{1}{2}\| u - v \|_{L^2}^2$ ist beschraenkt durch:
\begin{equation}
    \frac{d}{dt} E_{v}(t) \le - \frac{\nu}{2} \|\nabla u\|_{L^2}^2 + C_1(t) E_{v}(t) + C_2(t),
\end{equation}
wobei $C_1(t) = 2\|\nabla v\|_{L^\infty} + 1$ und $C_2(t) = \frac{1}{2}\|(v \cdot \nabla)v\|_{L^2}^2 + \frac{1}{2}\|f - \partial_t v\|_{L^2}^2 + \frac{\nu}{2}\|\nabla v\|_{L^2}^2$.
\end{lemma}

\begin{proof}
Wir subtrahieren die Gleichungen fuer $v$ von denen fuer $u$ und erhalten die Entwicklungsgleichung fuer $w = u - v$:
\begin{equation}
    \frac{d}{dt} E_{v}(t) = \int_\Omega w \cdot \partial_t w \, dx = \int_\Omega w \cdot \left( -(u \cdot \nabla)u - \nabla p + \nu \Delta u + f - \partial_t v \right) dx.
\end{equation}
Wegen $\nabla \cdot w = 0$ und der Periodizitaet von $\Omega$ verschwindet der Druckterm: $\int_\Omega w \cdot \nabla p \, dx = 0$.
Partielle Integration des viskosen Terms ergibt:
\begin{equation}
    \nu \int_\Omega w \cdot \Delta u \, dx = - \nu \int_\Omega \nabla w : \nabla u \, dx = - \nu \|\nabla u\|_{L^2}^2 + \nu \int_\Omega \nabla v : \nabla u \, dx.
\end{equation}
Anwenden der Cauchy-Schwarz- und Young-Ungleichung liefert:
\begin{equation}
    \nu \int_\Omega w \cdot \Delta u \, dx \le - \frac{\nu}{2} \|\nabla u\|_{L^2}^2 + \frac{\nu}{2} \|\nabla v\|_{L^2}^2.
\end{equation}
Fuer das kritische konvektive Integral $I := \int_\Omega w \cdot (u \cdot \nabla)u \, dx$ zerlegen wir $u = v + w$. Unter Nutzung der fundamentalen Aufhebungseigenschaft $\int_\Omega w \cdot (y \cdot \nabla)w \, dx = 0$ fuer jedes divergenzfreie $y$ isolieren wir die verbleibenden Terme:
\begin{equation}
    - I = - \int_\Omega w \cdot (v \cdot \nabla)v \, dx - \int_\Omega w \cdot (w \cdot \nabla)v \, dx.
\end{equation}
Diese Terme sind beschraenkt durch:
\begin{equation}
    |-I| \le \|w\|_{L^2} \|(v \cdot \nabla)v\|_{L^2} + \|\nabla v\|_{L^\infty} \|w\|_{L^2}^2.
\end{equation}
Einsetzen dieser Schranken in die Zeitableitungsgleichung und Anwenden der Young-Ungleichung auf die restlichen linearen Terme $\|w\|_{L^2} \|f - \partial_t v\|_{L^2}$ und $\|w\|_{L^2} \|(v \cdot \nabla)v\|_{L^2}$ liefert das behauptete Resultat.
\end{proof}

Da $\mathcal{A}$ der globale Leray-Hopf-Attraktor auf $\mathbb{T}^3$ ist, sichern Abschaetzungen der absorbierenden Kugel die gleichmaessigen Schranken $\sup_{v \in \mathcal{A}} C_1(t) < \infty$ und $\sup_{v \in \mathcal{A}} C_2(t) < \infty$ global fuer alle $t$. Infimumbildung ueber $\mathcal{A}$ ueberfuehrt die individuelle Energieschranke in eine absolute Stabilitaetsgrenze fuer die Stroemungsdivergenz $H(u \mid \mathcal{A})$.

\begin{remark}[Der Minimierer ist keine einzelne NS-Trajektorie]
\label{rem:minimiser-trajectory}
Im Beweis von Satz~\ref{thm:main} ist die minimierende Auswahl
$v(t) \in \mathcal{A}$ das naechste Attraktorelement zur Zeit~$t$,
\emph{nicht} eine einzelne NS-Loesungstrajektorie. Waehrend $u(t)$ evolviert,
kann der Minimierer zwischen verschiedenen Punkten von~$\mathcal{A}$ springen.
Lemma~\ref{lem:entropy-bound} ist fuer eine \emph{feste} NS-Loesung~$v$
formuliert; bei Anwendung auf den zeitvariablen Minimierer muss der Term
$\partial_t v$ in $C_2$ als die gesamte Aenderungsrate der Auswahl
interpretiert werden (einschliesslich Spruenge), nicht lediglich als die
NS-Evolution bei~$v(t)$. Genau deshalb ist Assumption~G (zur Kontrolle
von $\partial_t v$) oder die BV-Modifikation (Stieltjes-Integral,
Gleichung~\eqref{eq:stieltjes-bound}) erforderlich.
\end{remark}

\section{Hauptresultat: Bedingter Ausschluss von Singularitaeten in endlicher Zeit}

Wir formulieren nun das zentrale Regularitaetsresultat. Zwei Hilfslemmas werden
vorab bewiesen, um das Hauptargument in sich geschlossen zu halten.

% ------------------------------------------------------------------
\subsection*{Lemma A (Biot--Savart logarithmische Abschaetzung)}
% ------------------------------------------------------------------
\begin{lemma}[Biot--Savart-Abschaetzung auf $\mathbb{T}^3$]
\label{lem:BS}
Sei $u$ ein periodisches, divergenzfreies Vektorfeld auf $\mathbb{T}^3$, und sei
$\omega = \nabla\times u$ seine Vortizitaet. Fuer jedes $s > 3/2$ existiert eine
Konstante $C_{\mathrm{BS}} = C_{\mathrm{BS}}(s,\Omega) > 0$ (die nur von
$s$ und der Geometrie von $\Omega = \mathbb{T}^3$ abhaengt, \emph{nicht} aber von $u$), sodass
\begin{equation}\label{eq:BS}
    \|\nabla u\|_{L^\infty}
    \;\le\;
    C_{\mathrm{BS}}\,\|\omega\|_{L^\infty}
    \!\left(1 + \log\frac{\|\omega\|_{H^s}}{\|\omega\|_{L^\infty}}\right),
    \qquad \text{falls } \|\omega\|_{L^\infty} > 0.
\end{equation}
\end{lemma}

\begin{proof}
Wir rekonstruieren $\nabla u$ aus $\omega$ mittels des periodischen Biot--Savart-Gesetzes:
$u = K * \omega$, wobei $K$ der matrixwertige Faltungskern auf
$\mathbb{T}^3$ ist, der aus der Inversion des Rotor-Divergenz-Systems entsteht.
Insbesondere gilt $\nabla u = \nabla K * \omega$, und $\nabla K$ ist ein
Calder\'{o}n--Zygmund-Singulaerintegral-Kern (vgl.\ \cite{BertozziMajda2002},
Kapitel~2).

\medskip
\noindent\textbf{Schritt 1 (Littlewood--Paley-Zerlegung).}
Seien $\{\Delta_j\}_{j\ge 0}$ die Standard-Littlewood--Paley-Projektionsoperatoren
auf $\mathbb{T}^3$, assoziiert mit einer glatten dyadischen Zerlegung der Eins
$\{\varphi_j\}$ im Frequenzraum, sodass
$\operatorname{supp}\hat\varphi_j \subset \{|\xi| \sim 2^j\}$ und
$\sum_{j\ge 0}\varphi_j(\xi) = 1$ fuer alle $\xi \neq 0$.
Wir spalten $\nabla u$ bei einem freien ganzzahligen Parameter $N \ge 1$ auf:
\[
    \nabla u \;=\; \underbrace{\sum_{j=0}^{N} \Delta_j(\nabla u)}_{=:\,S_{\le N}}
                  \;+\; \underbrace{\sum_{j=N+1}^{\infty} \Delta_j(\nabla u)}_{=:\,S_{>N}}.
\]

\medskip
\noindent\textbf{Schritt 2 (Niedrige Frequenzen $j \le N$).}
Da $\nabla K$ ein Calder\'{o}n--Zygmund-Operator nullter Ordnung auf
$\mathbb{T}^3$ ist, bildet er $L^\infty$ nach $\mathrm{BMO}$ ab, und auf jedem Littlewood--Paley-Block $\Delta_j$ liefert die Bernstein-Einbettung
$\mathrm{BMO} \hookrightarrow L^\infty$ (bei fester Frequenzskala)
\[
    \|\Delta_j(\nabla u)\|_{L^\infty}
    \;\le\; C_{\mathrm{CZ}}\,\|\Delta_j \omega\|_{L^\infty}
    \;\le\; C_{\mathrm{CZ}}\,\|\omega\|_{L^\infty},
\]
wobei $C_{\mathrm{CZ}}$ nur von der Operatornorm von $\nabla K$ und
der Geometrie von $\mathbb{T}^3$ abhaengt. Summation ueber $j = 0, \ldots, N$:
\begin{equation}\label{eq:BS-low}
    \|S_{\le N}\|_{L^\infty}
    \;\le\; C_{\mathrm{CZ}}\,(N+1)\,\|\omega\|_{L^\infty}.
\end{equation}

\medskip
\noindent\textbf{Schritt 3 (Hohe Frequenzen $j > N$).}
Fuer jeden dyadischen Block mit Fourier-Traeger in $\{|\xi| \sim 2^j\}$ liefert die
Bernstein-Ungleichung auf $\mathbb{T}^3$ (mit $p=2$, $q=\infty$):
\[
    \|\Delta_j(\nabla u)\|_{L^\infty}
    \;\le\; C_B\,2^{3j/2}\,\|\Delta_j(\nabla u)\|_{L^2}.
\]
Da $u = K * \omega$ und $\widehat{\nabla K}(\xi)$ ein Symbol nullter Ordnung ist
(oben und unten beschraenkt auf jeder Frequenzschale), gilt
$\|\Delta_j(\nabla u)\|_{L^2} \le C\,\|\Delta_j \omega\|_{L^2}$.
Ausdruecken von $\|\Delta_j \omega\|_{L^2}$ mittels der $H^s$-Norm
durch $\|\Delta_j\omega\|_{L^2} \le 2^{-js}\,\bigl(2^{js}\|\Delta_j\omega\|_{L^2}\bigr)
\le 2^{-js}\,\|\omega\|_{H^s}$ ergibt
\[
    \|\Delta_j(\nabla u)\|_{L^\infty}
    \;\le\; C'\,2^{j(3/2-s)}\,\|\omega\|_{H^s}.
\]
Da $s > 3/2$, ist der Exponent $\alpha := s - 3/2 > 0$ strikt positiv,
und die Restsumme ist eine konvergente geometrische Reihe:
\begin{equation}\label{eq:BS-high}
    \|S_{>N}\|_{L^\infty}
    \;\le\; \sum_{j>N} C'\,2^{-j\alpha}\,\|\omega\|_{H^s}
    \;=\; \frac{C'\,2^{-(N+1)\alpha}}{1 - 2^{-\alpha}}\,\|\omega\|_{H^s}
    \;\le\; C''\,2^{-N\alpha}\,\|\omega\|_{H^s}.
\end{equation}

\medskip
\noindent\textbf{Schritt 4 (Optimierung ueber $N$).}
Kombination von \eqref{eq:BS-low} und \eqref{eq:BS-high}:
\[
    \|\nabla u\|_{L^\infty}
    \;\le\; C_{\mathrm{CZ}}\,(N+1)\,\|\omega\|_{L^\infty}
          + C''\,2^{-N\alpha}\,\|\omega\|_{H^s}.
\]
Wir waehlen $N \in \mathbb{N}$, sodass die beiden Terme ungefaehr ausgeglichen sind.
Setzen von
\[
    2^{N\alpha} \;=\; \frac{\|\omega\|_{H^s}}{\|\omega\|_{L^\infty}},
    \qquad\text{d.h.}\qquad
    N \;=\; \frac{1}{\alpha\,\ln 2}\,
            \ln\!\frac{\|\omega\|_{H^s}}{\|\omega\|_{L^\infty}},
\]
liefert fuer den Hochfrequenzbeitrag
$C''\,2^{-N\alpha}\,\|\omega\|_{H^s} = C''\,\|\omega\|_{L^\infty}$.
Fuer den Niederfrequenzanteil gilt
\[
    C_{\mathrm{CZ}}\,(N+1)\,\|\omega\|_{L^\infty}
    \;\le\; \tilde C\,\|\omega\|_{L^\infty}\,
            \Bigl(1 + \ln\!\frac{\|\omega\|_{H^s}}{\|\omega\|_{L^\infty}}\Bigr).
\]
Addition beider Abschaetzungen und Absorption der universellen Konstanten in
$C_{\mathrm{BS}} = C_{\mathrm{BS}}(s,\Omega)$ ergibt
\[
    \|\nabla u\|_{L^\infty}
    \;\le\;
    C_{\mathrm{BS}}\,\|\omega\|_{L^\infty}\,
    \Bigl(1 + \log\frac{\|\omega\|_{H^s}}{\|\omega\|_{L^\infty}}\Bigr),
\]
was die behauptete Ungleichung \eqref{eq:BS} ist.
\end{proof}

\begin{remark}
Die Konstante $C_{\mathrm{BS}}$ kann explizit gemacht werden: Sorgfaeltiges Nachverfolgen der
Calder\'{o}n--Zygmund-Norm von $K$ und der Bernstein-Konstante ergibt
$C_{\mathrm{BS}} \lesssim C(s)(1 + \|K\|_{L^{3/2,\infty}})$. Insbesondere
ist $C_{\mathrm{BS}}$ \emph{gleichmaessig} ueber $\mathcal{A}$, da alle Attraktorelemente
denselben umgebenden Torus teilen.
\end{remark}

% ------------------------------------------------------------------
\subsection*{Lemma B (Enstrophie-Bilanz und punktweiser Blow-up)}
% ------------------------------------------------------------------
\begin{lemma}[Enstrophie-Bilanz und punktweiser Blow-up]
\label{lem:enstrophy}
Sei $u$ eine Leray--Hopf-Loesung der Navier--Stokes-Gleichungen mit glattem
Antrieb $f \in L^\infty_t H^1_x$ auf $\mathbb{T}^3$. Die Enstrophie erfuellt (im
schwachen Sinne fuer f.u.\ $t$)
\begin{equation}\label{eq:enstrophy}
    \frac{1}{2}\frac{d}{dt}\|\omega(t)\|_{L^2}^2
    = -\nu\|\nabla\omega\|_{L^2}^2
      + \int_\Omega (\omega\cdot\nabla)u\cdot\omega\,dx
      + \int_\Omega (\nabla\times f)\cdot\omega\,dx.
\end{equation}
Wenn $\int_0^{T^*}\|\omega(t)\|_{L^\infty}\,dt = \infty$ (BKM-Blow-up), dann
divergiert die Enstrophie punktweise:
\begin{equation}\label{eq:diss-blowup}
    \limsup_{t \nearrow T^*}\|\nabla u(t)\|_{L^2}^2 = \infty.
\end{equation}
\end{lemma}

\begin{remark}[Was divergiert und was nicht bei Blow-up]
\label{rem:divergence-landscape}
Unter einem hypothetischen BKM-Blow-up bei $T^* < \infty$:
\begin{itemize}
  \item \textbf{Divergiert (punktweise):}\;
    $\|\omega(t)\|_{L^\infty} \to \infty$,\;
    $\|\nabla u(t)\|_{L^2} \to \infty$ fuer $t \nearrow T^*$.
  \item \textbf{Bleibt endlich (Integral):}\;
    $\int_0^{T^*}\|\nabla u\|_{L^2}^2\,dt < \infty$.
    Dies folgt aus der Leray-Energieidentitaet:
    $\nu\int_0^{T^*}\|\nabla u\|^2
    = \tfrac{1}{2}\|u_0\|^2 - \tfrac{1}{2}\|u(T^*)\|^2
    + \int_0^{T^*}\langle f,u\rangle\,dt$,
    deren rechte Seite beschraenkt ist, da $\|u\|_{L^2}$ durch die
    Energieungleichung kontrolliert wird, $f$ glatt ist und $T^* < \infty$.
\end{itemize}
Die punktweise Divergenz von $\|\nabla u(t)\|_{L^2}$ ist vereinbar mit
endlicher $L^2$-in-Zeit-Integrierbarkeit: zum Beispiel divergiert
$(T^*-t)^{-\alpha}$ punktweise, ist aber integrierbar fuer $\alpha < 1$.
Diese Unterscheidung ist zentral fuer die bedingte Struktur von
Satz~\ref{thm:main}; siehe Bemerkung~\ref{rem:diss-gap} unten.
\end{remark}

\begin{proof}
Der Beweis besteht aus zwei Teilen: der Herleitung der Enstrophie-Identitaet
\eqref{eq:enstrophy} und der Implikation \eqref{eq:diss-blowup}.

\medskip
\noindent\textbf{Teil 1: Herleitung der Enstrophie-Identitaet.}
Wir wenden den Rotationsoperator $\nabla\times$ auf die Navier--Stokes-Gleichung
$\partial_t u + (u\cdot\nabla)u = -\nabla p + \nu\Delta u + f$ an.
Da $\nabla\times(\nabla p) = 0$ identisch gilt, faellt der Druck heraus.
Die Standard-Vektoridentitaet fuer den nichtlinearen Term liefert die Vortizitaetsgleichung:
\begin{equation}\label{eq:vort-eq}
    \partial_t \omega + (u\cdot\nabla)\omega
    \;=\; (\omega\cdot\nabla)u + \nu\Delta\omega + \nabla\times f.
\end{equation}
Wir testen nun \eqref{eq:vort-eq} gegen $\omega$ in $L^2(\Omega)$, d.h.\
wir bilden das Skalarprodukt $\int_\Omega (\cdot)\cdot\omega\,dx$:
\begin{align}
    \int_\Omega (\partial_t\omega)\cdot\omega\,dx
    &\;=\; \frac{1}{2}\frac{d}{dt}\|\omega\|_{L^2}^2, \label{eq:enst-lhs}\\[4pt]
    \int_\Omega \bigl[(u\cdot\nabla)\omega\bigr]\cdot\omega\,dx
    &\;=\; 0,\label{eq:enst-transport}
\end{align}
wobei \eqref{eq:enst-transport} gilt, weil $\nabla\cdot u = 0$ und
$\Omega = \mathbb{T}^3$ periodisch ist: partielle Integration ergibt
$\int_\Omega (u\cdot\nabla)\tfrac{|\omega|^2}{2}\,dx
= -\int_\Omega \tfrac{|\omega|^2}{2}\,(\nabla\cdot u)\,dx = 0$.
Fuer den viskosen Term liefert partielle Integration (zweimal, unter Nutzung der Periodizitaet)
\[
    \nu\int_\Omega (\Delta\omega)\cdot\omega\,dx
    \;=\; -\nu\|\nabla\omega\|_{L^2}^2.
\]
Zusammenfassen aller Terme ergibt die Identitaet \eqref{eq:enstrophy}.

\medskip
\noindent\textbf{Teil 2: Punktweise Divergenz der Enstrophie.}
Angenommen $\int_0^{T^*}\|\omega(t)\|_{L^\infty}\,dt = \infty$. Wir zeigen,
dass $\limsup_{t \nearrow T^*}\|\nabla u(t)\|_{L^2}^2 = \infty$.

\smallskip
\noindent\emph{Schritt 2a (Wirbelstreckungsschranke).}
Das Wirbelstreckungsintegral in \eqref{eq:enstrophy} erfuellt
\begin{equation}\label{eq:stretch-bound}
    \Bigl|\int_\Omega (\omega\cdot\nabla)u\cdot\omega\,dx\Bigr|
    \;\le\; \|\nabla u\|_{L^\infty}\,\|\omega\|_{L^2}^2.
\end{equation}
Dies folgt aus der punktweisen Abschaetzung
$|(\omega\cdot\nabla)u\cdot\omega| \le |\nabla u|\,|\omega|^2$ und
der H\"older-Ungleichung.

\smallskip
\noindent\emph{Schritt 2b (Einsetzen der Biot--Savart-Abschaetzung).}
Nach Lemma~\ref{lem:BS}, fuer ein festes $s > 3/2$:
\[
    \|\nabla u\|_{L^\infty}
    \;\le\; C_{\mathrm{BS}}\,\|\omega\|_{L^\infty}\,
            \bigl(1 + \log^+\!\tfrac{\|\omega\|_{H^s}}{\|\omega\|_{L^\infty}}\bigr).
\]
Auf $\mathbb{T}^3$ liefert die Gagliardo--Nirenberg-Interpolationsungleichung
$\|\omega\|_{H^s} \le C\,\|\omega\|_{L^2}^{1-s/k}\,\|\omega\|_{H^k}^{s/k}$
fuer jedes ganzzahlige $k > s$; Wahl von $k = \lceil s \rceil + 1$ und
$\|\omega\|_{H^k} \le C'\,\|\nabla^k \omega\|_{L^2} + C''\,\|\omega\|_{L^2}$
(Poincar\'{e} auf $\mathbb{T}^3$; vgl.\ \cite{ConstantinFoias1988}).
Insbesondere waechst $\log^+\!\tfrac{\|\omega\|_{H^s}}{\|\omega\|_{L^\infty}}$
hoechstens logarithmisch in den hoeheren Sobolev-Normen von~$\omega$.

Einsetzen in \eqref{eq:stretch-bound}:
\begin{equation}\label{eq:stretch-BS}
    \Bigl|\int_\Omega (\omega\cdot\nabla)u\cdot\omega\,dx\Bigr|
    \;\le\; C_{\mathrm{BS}}\,\|\omega\|_{L^\infty}\,
            \bigl(1 + \log^+\!\tfrac{\|\omega\|_{H^s}}{\|\omega\|_{L^\infty}}\bigr)\,
            \|\omega\|_{L^2}^2.
\end{equation}

\smallskip
\noindent\emph{Schritt 2c (Gronwall-Kette).}
Einsetzen von \eqref{eq:stretch-BS} und der Antriebsschranke
$|\int_\Omega(\nabla\times f)\cdot\omega\,dx|
\le \|\nabla\times f\|_{L^2}\|\omega\|_{L^2}
\le C_f + \tfrac{1}{4}\|\omega\|_{L^2}^2$ in \eqref{eq:enstrophy} ergibt
\[
    \frac{1}{2}\frac{d}{dt}\|\omega\|_{L^2}^2
    \;\le\;
    -\nu\|\nabla\omega\|_{L^2}^2
    + C_{\mathrm{BS}}\,\|\omega\|_{L^\infty}\,
      \bigl(1 + \log^+\!\tfrac{\|\omega\|_{H^s}}{\|\omega\|_{L^\infty}}\bigr)\,
      \|\omega\|_{L^2}^2
    + C_f'.
\]
Definiere $\Phi(t) := \|\omega(t)\|_{L^2}^2$. Da der logarithmische Faktor
hoechstens wie $C + C\log\|\nabla\omega\|_{L^2}$ waechst, und nach der Young-Ungleichung
$a\log b \le \varepsilon b^2 + C_\varepsilon a(\log a + 1)$ fuer $a,b > 0$,
kann der logarithmisch-streckende Term in die viskose Dissipation
$\nu\|\nabla\omega\|_{L^2}^2$ absorbiert werden, auf Kosten eines multiplikativen Gronwall-Faktors
mit $\|\omega\|_{L^\infty}$. Genauer existieren Konstanten
$c_1, c_2 > 0$ (abhaengig von $\nu, s, C_{\mathrm{BS}}$), sodass
\begin{equation}\label{eq:Gronwall-enst}
    \frac{d}{dt}\Phi(t)
    \;\le\;
    -\nu\|\nabla\omega\|_{L^2}^2
    + c_1\,\|\omega(t)\|_{L^\infty}\,(1 + \log^+\Phi(t))\,\Phi(t)
    + c_2.
\end{equation}
Falls $\int_0^{T^*}\|\omega\|_{L^\infty}\,dt = \infty$ (BKM), divergiert der
Gronwall-Faktor $\exp\!\bigl(c_1\int_0^t\|\omega\|_{L^\infty}(1+\log^+\Phi)\,d\tau\bigr)$,
was $\Phi(t) = \|\omega(t)\|_{L^2}^2 \to \infty$ fuer
$t \nearrow T^*$ erzwingt.

\smallskip
\noindent\emph{Schritt 2d (Hodge-Identitaet und punktweiser Blow-up).}
Auf $\mathbb{T}^3$ mit $\nabla\cdot u = 0$ und Mittelwert Null gilt die Hodge-Identitaet
$\|\omega\|_{L^2} = \|\nabla u\|_{L^2}$.
Da $\Phi(t) = \|\omega(t)\|_{L^2}^2 = \|\nabla u(t)\|_{L^2}^2 \to \infty$,
erhalten wir den \emph{punktweisen} Blow-up der Enstrophie
\eqref{eq:diss-blowup}.

\medskip
\noindent\textbf{Wichtige Unterscheidung:}\;
Die Leray-Energieidentitaet impliziert, dass das \emph{Zeitintegral}
$\int_0^{T^*}\|\nabla u\|_{L^2}^2\,dt$ selbst unter
Blow-up endlich bleibt (siehe Bemerkung~\ref{rem:divergence-landscape}). Was divergiert,
ist der \emph{punktweise} Wert $\|\nabla u(t)\|_{L^2}^2$ fuer
$t \nearrow T^*$, nicht sein Zeitintegral. Diese Unterscheidung ist
entscheidend fuer die bedingte Struktur von Satz~\ref{thm:main}.
\end{proof}

% ------------------------------------------------------------------
\subsection*{Assumption G (Gronwall-Regularitaet der Attraktorprojektion)}
\label{sec:assumption-G}
% ------------------------------------------------------------------

Die folgende Annahme erfasst die zentrale strukturelle Eigenschaft, die unser
Argument benoetigt, aber nicht aus ersten Prinzipien herleitet.

\begin{definition}[Assumption G]\label{ass:G}
Sei $v(t) \in \mathcal{A}$ die messbare Auswahl, die
$H(u(t)\mid\mathcal{A}) = \tfrac{1}{2}\|u(t)-v(t)\|_{L^2}^2$ minimiert. Wir sagen,
dass \textbf{Assumption~G} gilt, wenn der zeitabhaengige Minimierer $v(t)$ erfuellt:
\begin{enumerate}
    \item[\emph{(G1)}] Die Abbildung $t \mapsto v(t)$ ist absolut stetig in
          $L^2(\Omega)$, sodass $\partial_t v(t)$ fuer f.u.\ $t$ existiert.
    \item[\emph{(G2)}] Die Gronwall-Koeffizienten $C_1(t)$ und $C_2(t)$ aus
          Lemma~\ref{lem:entropy-bound}, ausgewertet an der minimierenden Auswahl
          $v(t)$, erfuellen $C_1, C_2 \in L^1_{\mathrm{loc}}([0,\infty))$.
\end{enumerate}
\end{definition}

\begin{definition}[Assumption G$'$ (abgeschwaecht)]\label{ass:Gprime}
Wir sagen, dass \textbf{Assumption~G$'$} gilt, wenn:
\begin{enumerate}
    \item[\emph{(G1$'$)}] Die Abbildung $t \mapsto v(t)$ hat \emph{beschraenkte Variation}
          in $L^2(\Omega)$: $\mathrm{Var}_{L^2}(v; [0,T]) :=
          \sup \sum_{i} \|v(t_{i+1}) - v(t_i)\|_{L^2} < \infty$ fuer jedes
          $T > 0$.
    \item[\emph{(G2$'$)}] Die Gronwall-Koeffizienten erfuellen
          $C_1, C_2 \in L^1_{\mathrm{loc}}$ bezueglich des Stieltjes-Masses
          $\mathrm{d}\|v\|(t)$ (dem Totalvariationsmass von $v$).
\end{enumerate}
\end{definition}

\begin{proposition}[Abgeschwaechte Annahme genuegt]\label{prop:Gprime}
Satz~\ref{thm:main} bleibt gueltig, wenn Assumption~G durch die
schwaechere Assumption~G$'$ ersetzt wird.
\end{proposition}

\begin{proof}[Beweisskizze]
Die Gronwall-Ungleichung in Schritt~2 des Beweises von
Satz~\ref{thm:main} verwendet $\int_0^T C_1(t)\,\mathrm{d} t$ als kumulative
Schranke. Unter (G1$'$) kann der Minimierer $v(t)$ springen, aber die Totalvariation $\mathrm{Var}_{L^2}(v;[0,T])$ ist endlich. Der Schluesselterm
$\langle w, \partial_t v \rangle$ in der Zeitableitung von $H$ wird
durch das Stieltjes-Integral $\int_0^T \langle w(t),
\mathrm{d} v(t)\rangle$ ersetzt, das beschraenkt ist durch $\|w\|_{L^\infty_t L^2}\cdot
\mathrm{Var}_{L^2}(v;[0,T])$. Alle resultierenden Terme auf der rechten Seite
der integrierten Ungleichung sind endlich (siehe die BV-Modifikation
im Beweis von Satz~\ref{thm:main} fuer die detaillierte
Schranke~\eqref{eq:H-integral-BV}). Der Widerspruch $H < 0$ folgt
dann unter derselben Dissipations-Dominanz-Bedingung~\eqref{eq:diss-dom}.
\end{proof}

Wir identifizieren nun hinreichende geometrische Bedingungen, unter denen G$'$ ein
\emph{Satz} ist. Die Struktur spiegelt die DFC$\Rightarrow$NL$'$-Reduktion
in der Turbulenzsituation~\cite{FST-TU2026} wider: eine opake PDE-Annahme wird
durch verifizierbare geometrische Eigenschaften des Attraktors ersetzt.

\begin{definition}[Attractor Geometry Condition]\label{def:AGC}
Der Attraktor $\mathcal{A}$ erfuellt die \emph{Attractor Geometry
Condition} (AGC), wenn:
\begin{enumerate}
  \item[\textup{(AGC1)}] \textbf{Positiver Reach.}\;
    Es existiert $\delta > 0$, sodass jeder Punkt $x$ mit
    $\mathrm{dist}(x, \mathcal{A}) < \delta$ einen \emph{eindeutigen}
    naechsten Punkt $\pi(x) \in \mathcal{A}$ besitzt, und die Naechste-Punkt-Abbildung
    $\pi$ Lipschitz-stetig auf
    $\mathcal{A}^\delta := \{x : \mathrm{dist}(x,\mathcal{A}) < \delta\}$
    mit Lipschitz-Konstante~$1$ ist.
  \item[\textup{(AGC2)}] \textbf{Exponentielles Tracking.}\;
    Es existieren $C_0, \lambda > 0$, sodass fuer jede Leray--Hopf-Loesung
    $u(t)$ mit $u_0$ in der absorbierenden Kugel:
    $\mathrm{dist}(u(t), \mathcal{A}) \le C_0\,e^{-\lambda t}$ fuer alle
    $t \ge 0$.
\end{enumerate}
\end{definition}

\begin{lemma}[AGC impliziert G]\label{lem:AGC-implies-G}
Falls \textup{(AGC1)--(AGC2)} gelten, dann gilt Assumption~G
\textup{(}nicht bloss G$'$\textup{)}.
\end{lemma}

\begin{proof}
Nach (AGC2) existiert $t_0 > 0$, sodass
$\mathrm{dist}(u(t), \mathcal{A}) < \delta$ fuer alle $t \ge t_0$.
Fuer $t \ge t_0$ ist der Minimierer $v(t) = \pi(u(t))$ wohldefiniert, und
die Lipschitz-Eigenschaft von $\pi$ (AGC1) liefert:
\[
  \|v(t_2) - v(t_1)\|_{L^2}
  \;\le\;
  \|u(t_2) - u(t_1)\|_{L^2}
  \qquad \forall\, t_0 \le t_1 < t_2.
\]
Da $u(t)$ absolut stetig in $L^2$ ist (Leray--Hopf-Regularitaet),
ist $v(t)$ absolut stetig mit $\|\partial_t v\| \le \|\partial_t u\|$
f.u. Dies ergibt (G1). Fuer (G2) haengen die Gronwall-Koeffizienten $C_1(t),
C_2(t)$ von $\|v(t)\|_{W^{1,\infty}}$ und $\|\partial_t v(t)\|$ ab;
beide sind beschraenkt durch
Proposition~\ref{prop:attractor-regularity}(iii) und die Lipschitz-Abschaetzung.
\end{proof}

\begin{remark}[Status von AGC und die Reach-Luecke]
\label{rem:AGC-status}
\textup{(AGC2)} ist eine Standardkonsequenz der Eigenschaft der absorbierenden Kugel
und ist bewiesen fuer 3D NS auf $\mathbb{T}^3$ mit glattem
Antrieb~\cite{Temam2001}.

\textup{(AGC1)} ist der schwierige Teil.
Positiver Reach von $\mathcal{A}$ wuerde aus der Existenz einer
\emph{Inertialen Mannigfaltigkeit} folgen -- einem Lipschitz-Graph-Attraktor mit glattem
Normalenbuendel. Inertiale Mannigfaltigkeiten sind bewiesen fuer hyperviskose NS
($(-\Delta)^\beta$, $\beta \ge 5/4$) und fuer 2D
NS~\cite{MalletParetSell1988,EdenFoiasNicolaenko1994}, bleiben aber offen fuer 3D NS mit $\beta=1$.
Fuer eine moderne Uebersicht zu Attraktordimension, Squeezing und exponentiellen
Attraktoren im Kontext dissipativer PDEs siehe Zelik~\cite{Zelik2022}.
Das Hindernis ist die unzureichende Spektralluecke des 3D-Stokes-Operators:
Eigenwertluecken $\mu_{N+1} - \mu_N$ wachsen wie $N^{2/3}$, waehrend
die Konstruktion inertialer Mannigfaltigkeiten Wachstum $\gtrsim N$ erfordert.

\medskip\noindent
\textbf{Hinreichende Bedingungen fuer G$'$ (schwaecher als AGC1).}\;
Fuer die BV-Version G$'$ genuegen drei progressiv schwaechere geometrische
Bedingungen jeweils. Wir formulieren sie als Propositionen, um die
Landschaft der Angriffsvektoren zu klaeren.

\begin{proposition}[Trajektorien-Reach]\label{prop:traj-reach}
\textup{(AGC1$'$)}\;
Angenommen es existiert $\delta > 0$, sodass die Naechste-Punkt-Projektion
$\pi$ wohldefiniert und Lipschitz auf der tubularen Menge
$\{x : \mathrm{dist}(x, \mathcal{A}) < \delta\}
\cap \{u(t) : t \ge t_0\}$
ist \textup{(}d.h., nur entlang der Trajektorie\textup{)}. Dann gilt G
auf $[t_0, \infty)$.
\end{proposition}

\begin{proof}
Identisch zu Lemma~\ref{lem:AGC-implies-G}, eingeschraenkt auf die
Trajektorie.
\end{proof}

\begin{proposition}[Lipschitz-Graph ueber niedrigen Moden]\label{prop:graph}
\textup{(AGC1$''$)}\;
Angenommen es existiert eine Galerkin-Trunkierung $N$ und eine Lipschitz-Abbildung
$\Phi : P_N L^2 \to Q_N L^2$ (wobei $P_N$ der Stokes-Projektor
auf die ersten $N$ Eigenmoden ist, $Q_N = I - P_N$), sodass
$\mathcal{A} \subset \mathrm{graph}(\Phi) =
\{x + \Phi(x) : x \in P_N L^2\}$.
Dann ist die Projektion $\pi(u) = x^* + \Phi(x^*)$ mit
$x^* = \arg\min_{x \in P_N\mathcal{A}} \|P_N u - x\|$ Lipschitz-stetig,
und G gilt.
\end{proposition}

\begin{proof}
Die Graph-Eigenschaft stellt sicher, dass $\pi$ ueber die
endlichdimensionale Projektion $P_N$ gefolgt von der Lipschitz-Abbildung
$\Phi$ faktorisiert. Da $P_N u(t)$ absolut stetig in $\mathbb{R}^N$ ist (Leray--Hopf-Loesungen sind absolut stetig in jeder endlichen Galerkin-Projektion),
erbt $v(t) = \pi(u(t))$ die absolute Stetigkeit.
\end{proof}

\begin{proposition}[BV-Auswahl aus Minimierer-Multivaluation]
\label{prop:bv-selection}
\textup{(AGC1$'''$)}\;
Angenommen die mengenwertige Abbildung $t \mapsto \arg\min_{v \in \mathcal{A}}
\|u(t)-v\|$ besitzt eine messbare Auswahl $v(t)$ mit
$\mathrm{Var}_{L^2}(v; [0,T]) < \infty$ fuer jedes $T > 0$.
Dann gilt G$'$ nach Definition.
\end{proposition}

\begin{proof}
Tautologisch: (AGC1$'''$) ist genau die Aussage (G1$'$).
\end{proof}

\noindent
\textbf{Hierarchie:}\;
$\textup{AGC1 (globaler Reach)}
\;\Longrightarrow\;
\textup{AGC1}' \textup{ (Traj.-Reach)}
\;\Longrightarrow\;
\textup{G}
\;\Longrightarrow\;
\textup{G}'
\;\Longleftarrow\;
\textup{AGC1}'''$.
Die Graph-Bedingung AGC1$''$ impliziert AGC1 in der Normalumgebung des Graphen
und damit~G.
Keine der Bedingungen AGC1--AGC1$'''$ ist bewiesen fuer den 3D-Navier--Stokes-Attraktor
mit klassischer Viskositaet $\beta = 1$. Jede stellt einen anderen
geometrischen Angriffsvektor dar, und ein Beweis einer \emph{beliebigen} davon wuerde die
Regularitaetsluecke schliessen.

\medskip\noindent
\textbf{Der offene Kern (Auswahlregularitaetsproblem).}\;
Die gesamte Regularitaetsfrage fuer die inkompressiblen 3D-Navier--Stokes-Gleichungen
reduziert sich \emph{innerhalb des Rahmenwerks dieser Arbeit} auf die
folgende einzelne Aussage:

\begin{center}
\fbox{\parbox{0.85\textwidth}{%
\textbf{Offenes Problem (Auswahlregularitaet).}\;
Sei $u(t)$ eine Leray--Hopf-Loesung auf $\mathbb{T}^3$ und
$\mathcal{A}$ der globale Attraktor. Existiert eine messbare
Auswahl $v(t) \in \mathcal{A}$ mit
$\|u(t)-v(t)\|_{L^2} \le \delta(t)$ (Tracking) und
$\mathrm{Var}_{L^2}(v;\,[0,T]) < \infty$ fuer jedes $T > 0$?
}}
\end{center}

\noindent
Eine bejahende Antwort impliziert G$'$, was globale Regularitaet impliziert
(Satz~\ref{thm:main}). Vier geometrische hinreichende Bedingungen sind
oben identifiziert (AGC1--AGC1$'''$); die staerkste (AGC1, positiver Reach)
ist aequivalent zur Existenz inertialer Mannigfaltigkeiten. Die schwaechste
(AGC1$'''$, BV-Auswahl) ist genau die obige Aussage.
Dies ist eine \emph{geometrische} Frage ueber die Struktur von
$\mathcal{A}$ als Teilmenge von $L^2$, unabhaengig von der spezifischen
Trajektorie~$u(t)$.

\medskip\noindent
\textbf{Bekannte Resultate, die beinahe genuegen.}\;
Mehrere klassische Resultate kommen der Loesung des Auswahlregularitaetsproblems nahe,
scheitern aber an einem praezisen Punkt:
\begin{itemize}
  \item \textbf{Federer-Reach / Prox-Regularitaet}
    (Federer 1959; Poliquin--Rockafellar--Thibault 2000):
    Mengen mit positivem Reach $r > 0$ haben Lipschitz-1 metrische Projektion
    auf der $r$-tubularen Umgebung. Dies wuerde trivial~G liefern.
    Jedoch erzwingt positiver Reach lokale $C^{1,1}$-Mannigfaltigkeitsstruktur
    und ist \emph{inkompatibel mit fraktaler Geometrie} -- Attraktoren
    mit nichtganzzahliger fraktaler Dimension haben generisch Reach~$0$.
  \item \textbf{Chistyakov BV-Auswahl} (2004):
    Falls eine kompaktwertige Mengenabbildung $F(t)$ beschraenkte Hausdorff-Variation hat,
    dann existiert eine BV-Auswahl mit
    $\mathrm{Var}(f) \le \mathrm{Var}_{\mathrm{Haus}}(F)$.
    Angewandt auf $F(t) = \arg\min_{v \in \mathcal{A}} \|u(t)-v\|$
    reduziert dies G$'$ auf: \emph{hat die Minimierermenge beschraenkte
    Hausdorff-Variation entlang NS-Trajektorien?} Dies ist offen.
  \item \textbf{Sweeping-Prozesse}
    (Moreau 1977; Edmond--Thibault 2006; Colombo--Monteiro Marques 2003):
    BV-Loesungen des Sweeping-Prozesses $-\dot{v} \in N_{C(t)}(v)$
    existieren fuer prox-regulaeres $C(t)$ mit beschraenkter Retraktion.
    Fuer nicht prox-regulaere Mengen \emph{existiert keine allgemeine BV-Theorie}.
  \item \textbf{Man\'e-Projektoren}
    (Man\'e 1981; Zelik 2022):
    Fuer Attraktoren mit endlicher fraktaler Dimension~$d$ existiert eine
    generische lineare Projektion $P : H \to \mathbb{R}^{2d+1}$, sodass
    $P|_{\mathcal{A}}$ injektiv ist und die Inverse
    $(P|_{\mathcal{A}})^{-1}$ H\"older-stetig ist.
    Bi-Lipschitz Man\'e-Projektoren (die~G via der Graph-Eigenschaft
    AGC1$''$ liefern wuerden) existieren nur unter zusaetzlichen dynamischen Bedingungen
    und sind fuer 3D NS \emph{nicht} verfuegbar.
\end{itemize}

\noindent
\textbf{Die Literaturluecke.}\;
Die Luecke zwischen bekannten Resultaten und dem Auswahlregularitaetsproblem ist
praezise charakterisiert: Alle BV-Projektionssaetze erfordern entweder
Prox-Regularitaet (Federer/Clarke/Thibault) oder beschraenkte Hausdorff-Variation
(Chistyakov), und keine davon folgt allein aus endlicher fraktaler Dimension.
Die vielversprechendste Bruecke waere:
\begin{enumerate}
  \item Eine ``log-Lipschitz BV''-Theorie, die Chistyakovs Satz auf
    Mengenabbildungen mit H\"older-kontrollierter (nicht Hausdorff-beschraenkter)
    Variation erweitert -- \emph{unten formalisiert} in
    den Definitionen~\ref{def:TLL}--\ref{def:LDI} und
    Lemma~\ref{lem:LL-BV}.
  \item Ein Beweis, dass der 3D-NS-Attraktor ``approximativ
    prox-regulaer'' ist (prox-regulaer nach Entfernung einer mageren singulaeren
    Menge) -- kombiniert mit der Sweeping-Prozess-Theorie.
  \item Erweiterung von Zeliks bi-Lipschitz Man\'e-Projektor-Resultaten
    von Ginzburg--Landau auf 3D NS -- dies wuerde AGC1$''$
    direkt liefern.
\end{enumerate}

\medskip\noindent
\textbf{Abhaengigkeitstabelle.}
\medskip

\begin{center}
\begin{tabular}{llll}
\hline
\textbf{Label} & \textbf{Aussage} & \textbf{Typ} & \textbf{Quelle} \\
\hline
AGC1 & Positiver Reach $\delta > 0$ & \textbf{Offen} &
  Folgt aus Inertialer Mannigfaltigkeit \\
AGC2 & Exponentielles Tracking & Satz &
  \cite{Temam2001} \\
Lem.~\ref{lem:AGC-implies-G} & AGC $\Rightarrow$ G &
  \textbf{Satz} & Diese Arbeit \\
Prop.~\ref{prop:Gprime} & G$'$ $\Rightarrow$
  Satz~\ref{thm:main} & \textbf{Satz} & Diese Arbeit \\
\hline
\end{tabular}
\end{center}

\medskip\noindent
Die logische Kette ist:
$\textup{AGC1} + \textup{AGC2}
\;\Longrightarrow\;
\textup{G}
\;\Longrightarrow\;
\textup{G}'
\;\Longrightarrow\;
\textup{Satz~\ref{thm:main}}$.
Der einzige unbewiesene Input ist \textup{(AGC1)}, eine geometrische
Eigenschaft von $\mathcal{A}$, unabhaengig von der spezifischen Trajektorie
$u(t)$. Dies reduziert die Regularitaetsfrage von einem PDE-Problem
(``bleibt $u$ glatt?'') auf ein geometrisches Problem (``hat
$\mathcal{A}$ positiven Reach?'').
\end{remark}

\begin{remark}[Bezug zur Vanishing-Viscosity-Selektion]
\label{rem:Drivas-VV-selection}
Das obige Auswahlregularitaetsproblem betrifft die Regularitaet der
Naechste-Punkt-Projektion auf den NS-Attraktor $\mathcal{A}$ -- eine
geometrische Frage ueber $\mathcal{A} \subset L^2$.

Ein konzeptuell verwandtes, aber technisch verschiedenes Selektionsproblem
entsteht im Vanishing-Viscosity-Programm: Drivas und
Nguyen~\cite{DrivasNguyen2019} bewiesen, dass -- falls die lokalen
Strukturfunktions-Exponenten zweiter Ordnung gleichmaessig in der
Viskositaet positiv bleiben -- Raum-Zeit $L^2$-schwache Limiten von
Leray--Hopf-Loesungen schwache Euler-Loesungen sind, diese Limiten
aber \emph{nicht-eindeutig und dissipativ} sein koennen.
Die physikalische Selektion unter solchen Limiten erfordert
zusaetzliche Kriterien (lokale Energie-Ungleichung, exaktes
Dritte-Ordnung-Gesetz, Skalen-Lokalitaet des Flusses).

De~Rosa, Drivas und Inversi~\cite{DeRosaDrivasInversi2024} bewiesen, dass jede
beschraenkte Euler-Loesung, die als Null-Viskositaets-Limes entsteht, anomale
Dissipation nicht auf einer Menge mit Hausdorff-Dimension kleiner als~$d$
tragen kann, wobei das Duchon--Robert-Framework fuer Randdissipation
modifiziert wurde.
De~Rosa, Drivas, Inversi und Isett~\cite{DeRosaDrivasInversiIsett2025}
zeigten ferner, dass die Duchon--Robert-Distribution Besov-regulaerer
schwacher Euler-Loesungen verbesserte Regularitaet in einem negativen
Besov-Raum besitzt und, falls ein Radon-Mass, absolutstetig bzgl.\ eines
geeigneten Hausdorff-Masses ist -- mit quantitativen Einschraenkungen an
die fraktale Dimension der dissipativen Menge.

Diese Resultate adressieren die \emph{Euler-seitige} Selektion (welche
schwachen Limiten physikalisch zulaessig sind), nicht die
\emph{NS-Attraktor-Geometrie} (ob $\mathcal{A}$ eine BV-Naechste-Punkt-Auswahl
zulaesst). Sie bestaetigen jedoch, dass Selektionsprobleme sowohl fuer das
Regularitaets- als auch fuer das Turbulenz-Closure-Programm zentral sind,
und dass das Duchon--Robert-Defektmass die natuerliche Bruecke zwischen
beiden darstellt.
\end{remark}

% ==================================================================
% Bruecke I: Log-Lipschitz-Projektion und BV-Kettenregel
% ==================================================================

\subsection*{Bruecke~I: Log-Lipschitz-Projektion und BV-Auswahl}

Positiver Reach (AGC1) ist hinreichend fuer~G, aber wahrscheinlich zu stark fuer
fraktale Attraktoren. \emph{Log-Lipschitz}-Regularitaet -- die kritische
Klasse zwischen H\"older (zu schwach fuer BV-Komposition) und Lipschitz
(zu stark fuer fraktale Geometrie) -- tritt natuerlich in dissipativen
PDE-Situationen auf und bietet einen zugaenglicheren Weg zu~G$'$.

Die Schluesselerkenntnis ist, dass eine \emph{globale} log-Lipschitz-Schranke
$\|P(x)-P(y)\| \le C\|x-y\|(1+\log(R/\|x-y\|))$
\emph{nicht} direkt BV von $P\circ u$ fuer AC-Kurven~$u$ impliziert:
die Partitionssummen $\sum_i \Delta_i(1+\log(R/\Delta_i))$ divergieren
logarithmisch bei Verfeinerung. Die korrekte Formulierung ersetzt das
Inkrement-Skalen-Log-Gewicht durch das
\emph{Abstand-zum-Attraktor}-Gewicht $\log(R/d(t))$, was eine
trajektorienweise Bedingung liefert, die diese Divergenz vermeidet.

\begin{definition}[Trajektorien-log-Lipschitz-Projektion]
\label{def:TLL}
Sei $u \in \mathrm{AC}([0,T];\,H)$ mit $u(t) \in \mathcal{A}^R
:= \{x \in H : \mathrm{dist}_H(x,\mathcal{A}) < R\}$ fuer alle~$t$.
Setze $d(t) := \mathrm{dist}_H(u(t),\mathcal{A})$ und
$v(t) := P(u(t))$ fuer eine Auswahl
$P: \mathcal{A}^R \to \mathcal{A}$ mit $P|_{\mathcal{A}} = \mathrm{id}$.
Wir sagen, dass $P$ die \emph{Trajektorien-log-Lipschitz-Bedingung} (TLL) entlang~$u$ erfuellt, wenn $C_{\mathrm{TLL}} > 0$ und
$h_0 > 0$ existieren, sodass fuer alle $t \in [0,T)$ mit $d(t) > 0$ und
alle $0 < h < \min(h_0,\, T-t)$:
\begin{equation}\label{eq:TLL}
  \|P(u(t\!+\!h)) - P(u(t))\|_H
  \;\le\;
  C_{\mathrm{TLL}}\,\|u(t\!+\!h) - u(t)\|_H\,
  \bigl(1 + \log\tfrac{R}{d(t)}\bigr).
  \tag{TLL}
\end{equation}
Auf der Menge $\{t : d(t)=0\}$ (wo $u(t) \in \mathcal{A}$) gilt
$P(u(t)) = u(t)$ durch die Idempotenz von~$P$, sodass keine separate Schranke
noetig ist.
\end{definition}

\begin{definition}[Log-Abstand-Integrierbarkeit]
\label{def:LDI}
In der Situation von Definition~\ref{def:TLL} sagen wir, dass
\emph{Log-Abstand-Integrierbarkeit} gilt, wenn
\begin{equation}\label{eq:LDI}
  \int_0^T \|u'(t)\|_H\,
  \bigl(1 + \log\tfrac{R}{d(t)}\bigr)\,\mathrm{d}t
  \;<\; \infty.
  \tag{LDI}
\end{equation}
\end{definition}

\begin{lemma}[BV-Kettenregel fuer log-Lipschitz-Projektionen]
\label{lem:LL-BV}
Sei $u \in \mathrm{AC}([0,T];\,H)$ mit $u(t) \in \mathcal{A}^R$
fuer alle~$t$, und sei $P$ eine Projektion, die \eqref{eq:TLL} erfuellt.
Falls \eqref{eq:LDI} gilt, dann ist $v := P\circ u \in \mathrm{BV}([0,T];\,H)$
mit der expliziten Schranke
\begin{equation}\label{eq:BV-bound}
  \mathrm{Var}_H(v;\,[0,T])
  \;\le\;
  C_{\mathrm{TLL}}
  \int_0^T \|u'(t)\|_H\,
  \bigl(1 + \log\tfrac{R}{d(t)}\bigr)\,\mathrm{d}t.
\end{equation}
Insbesondere gilt Assumption~G$'$
\textup{(}Definition~\ref{ass:Gprime}\textup{)}, und somit ist
Satz~\ref{thm:main} anwendbar.
\end{lemma}

\begin{proof}
\textbf{Schritt 1 (Metrische Ableitungsschranke).}\;
Fuer f.u.\ $t$ mit $d(t)>0$ ist die metrische Ableitung von $v$ bei $t$
beschraenkt durch
\[
  |v'|(t)
  \;:=\;
  \limsup_{h\to 0}\frac{\|v(t+h)-v(t)\|_H}{|h|}
  \;\le\;
  C_{\mathrm{TLL}}\,|u'|(t)\,
  \bigl(1 + \log\tfrac{R}{d(t)}\bigr),
\]
wobei $|u'|(t) = \lim_{h\to 0}\|u(t+h)-u(t)\|_H/|h|$ f.u.\ existiert,
da $u$ absolut stetig ist. Dies folgt direkt aus Division von \eqref{eq:TLL}
durch $|h|$ und Grenzuebergang $h\to 0$.

Auf $\{t : d(t)=0\}$: da $P|_{\mathcal{A}} = \mathrm{id}$,
gilt dort $v(t) = u(t)$, und $|v'|(t) = |u'|(t) \le |u'|(t)\,
(1+\log(R/d(t)))$ mit der Konvention $0 \cdot \infty = 0$
(das Integral in \eqref{eq:LDI} absorbiert diese Menge automatisch).

\medskip
\noindent
\textbf{Schritt 2 (Stetigkeit von $v$).}\;
Die Naechste-Punkt-Projektion auf eine nichtkonvexe kompakte Menge muss
im Allgemeinen nicht stetig sein. Die TLL-Bedingung~\eqref{eq:TLL}
sichert jedoch die Stetigkeit von $v = P\circ u$ \emph{entlang der Trajektorie}:
Fuer $t$ mit $d(t) > 0$ gilt
$\|v(t+h)-v(t)\| \le C_{\mathrm{TLL}}\,\|u(t+h)-u(t)\|\,
(1+\log(R/d(t))) \to 0$ fuer $h \to 0$
(da $u$ absolut stetig ist). Auf $\{d=0\}$ ist $v(t)=u(t)$
automatisch stetig. Also ist $v: [0,T] \to H$ stetig.

\medskip
\noindent
\textbf{Schritt 3 (BV-Kriterium).}\;
Ein Standardresultat der BV-Theorie (vgl.\ Ambrosio, Fusco und
Pallara~\cite{AmbrosioFuscoPallara2000}, \S3.2) besagt: Eine stetige
Funktion $v: [0,T] \to H$ mit
$|v'|(t) \le g(t)$ f.u.\ fuer ein $g \in L^1(0,T)$ hat beschraenkte
Variation mit $\mathrm{Var}(v) \le \int_0^T g(t)\,\mathrm{d}t$.
Anwendung mit $g(t) = C_{\mathrm{TLL}}\,|u'|(t)\,
(1+\log(R/d(t)))$ (integrierbar nach \eqref{eq:LDI}) ergibt
\eqref{eq:BV-bound}.

\medskip
\noindent
\textbf{Schritt 4 (Konsequenz).}\;
Da $\mathrm{Var}_H(v;\,[0,T]) < \infty$, ist Assumption~G$'$
(Definition~\ref{ass:Gprime}) erfuellt.
Proposition~\ref{prop:Gprime} liefert dann die Aussage von
Satz~\ref{thm:main}: kein Blow-up in endlicher Zeit.
\end{proof}

\begin{remark}[Zusammenhang mit der globalen log-Lipschitz-Bedingung]
\label{rem:global-LL}
Eine \emph{globale} log-Lipschitz-Auswahl
$P : \mathcal{A}^R \to \mathcal{A}$ mit
$\|P(x)-P(y)\| \le C\|x-y\|(1+\log(R/\|x-y\|))$
fuer alle $x,y \in \mathcal{A}^R$ impliziert \emph{nicht} direkt TLL
\eqref{eq:TLL}. Der Grund ist strukturell: fuer feine Partitionen
koennen die Inkremente $\|u(t_{i+1})-u(t_i)\|$ viel kleiner als
$d(t_i)$ sein, sodass $\log(R/\|u(t_{i+1})-u(t_i)\|) \gg \log(R/d(t_i))$.
Die resultierenden Partitionssummen
$\sum_i \Delta_i(1+\log(R/\Delta_i))$ divergieren logarithmisch,
selbst fuer Lipschitz-Kurven~$u$.

TLL ersetzt das \emph{Inkrement-Skalen}-Gewicht
$\log(R/\|x-y\|)$ durch das \emph{Abstand-Skalen}-Gewicht
$\log(R/d(t))$. Dies ist die geometrisch natuerliche Skala: sie misst,
wie ``rau'' der Attraktorrand auf der fuer die Trajektorie relevanten Skala ist,
nicht auf der Skala des Partitionsgitters.
Die globale LL-Bedingung sollte als Motivation fuer TLL betrachtet werden,
nicht als Ersatz.
\end{remark}

Der gesamte neue mathematische Inhalt konzentriert sich auf die
Integrierbarkeitsbedingung~\eqref{eq:LDI}. Wir identifizieren zwei
hinreichende Bedingungen, unter denen \eqref{eq:LDI} aus der
Struktur der Navier--Stokes-Gleichungen folgt.

\begin{proposition}[Variante~A: Log-Abstand hat beschraenkte Variation]
\label{prop:LDI-A}
Falls $\log(R/d(t)) \in \mathrm{BV}_{\mathrm{loc}}(0,T)$, dann
gilt \eqref{eq:LDI}.
\end{proposition}

\begin{proof}
Das Produkt einer $\mathrm{AC}$-Funktion ($\|u'\|$) und einer $\mathrm{BV}$-Funktion ($1 + \log(R/d)$) ist integrierbar nach der Standard-Formel
fuer partielle Integration fuer $\mathrm{AC}\times\mathrm{BV}$-Paare.
\end{proof}

\begin{proposition}[Variante~B: Subniveau-Zeitmass-Kontrolle]
\label{prop:LDI-B}
Angenommen es existieren $C_*, \alpha > 0$, sodass
\begin{equation}\label{eq:STC}
  \bigl|\{t \in [0,T] : d(t) \le \varepsilon\}\bigr|
  \;\le\;
  C_*\,\varepsilon^\alpha
  \qquad \text{fuer alle } \varepsilon \in (0, R].
  \tag{STC}
\end{equation}
Dann gilt \eqref{eq:LDI}.
\end{proposition}

\begin{proof}
Mittels $\log(R/d(t)) = \int_{d(t)}^R s^{-1}\,\mathrm{d}s$ und Tonelli:
$\int_0^T \|u'\|\,\log(R/d)\,\mathrm{d}t
= \int_0^R s^{-1}(\int_{\{d \le s\}} \|u'\|\,\mathrm{d}t)\,\mathrm{d}s$.
Nach H\"older und \eqref{eq:STC} ist das innere Integral beschraenkt durch
$\|u'\|_{L^2}\,C_*^{1/2}\,s^{\alpha/2}$. Einsetzen:
$\int_0^R s^{\alpha/2-1}\,\mathrm{d}s = (2/\alpha)\,R^{\alpha/2} < \infty$,
da $\alpha > 0$.
Die Erweiterung auf allgemeine $L^p$-Schranken folgt durch dieselbe
Layer-Cake-Zerlegung mit $\|u'\|_{L^p}$ anstelle von $\|u'\|_{L^2}$.
\end{proof}

\begin{remark}[Physikalische Plausibilitaet von TLL und LDI in der
Navier--Stokes-Situation]
\label{rem:LDI-plausibility}
Die dissipative Struktur der Navier--Stokes-Gleichungen liefert
zwei unabhaengige Kontrollen, die sowohl TLL als auch LDI stuetzen:
\begin{enumerate}[label=(\roman*)]
  \item \textbf{Exponentielles Tracking (AGC2):}\;
    $d(t) \le C_0\,e^{-\lambda t}$ fuer $t \ge t_0$
    (Temam~\cite{Temam2001}), also
    $\log(R/d(t)) \le \lambda t + C$ auf $[t_0, T]$ -- hoechstens
    linear wachsend, daher $\log(R/d) \in L^\infty_{\mathrm{loc}}$
    im absorbierenden Regime.
  \item \textbf{Energiedissipation:}\;
    Fuer glatte Loesungen auf $[0,T]$ (was das Regime im
    Widerspruchsbeweis ist) gilt
    $\int_0^T \|u'(t)\|_{L^2}\,\mathrm{d}t < \infty$.
\end{enumerate}
Kombiniert ergibt dies
$\int_0^T \|u'\|\,\log(R/d)\,\mathrm{d}t
\le \|u'\|_{L^1}\,\|\log(R/d)\|_{L^\infty}$,
was auf jedem kompakten Intervall $[0,T]$ mit
$T < T^*$ endlich ist. Somit ist \eqref{eq:LDI} automatisch erfuellt
im Vor-Blow-up-Regime.

Die Subtilitaet entsteht an der hypothetischen Blow-up-Zeit $T^*$:
$\|u'(t)\| \to \infty$ fuer $t \nearrow T^*$, und die Frage ist,
ob das Log-Gewicht dieses Wachstum verstaerkt oder absorbiert.
Unter Blow-up bewegt sich die Loesung \emph{weg} vom Attraktor
in $H^1$ (da $\|u\|_{H^1} \to \infty$, waehrend Attraktorelemente
gleichmaessig beschraenkt sind), kann aber in $L^2$ nahe bleiben -- die $L^2$-Energie bleibt beschraenkt nach der Leray--Hopf-Ungleichung. Das Zusammenspiel
zwischen der divergierenden Geschwindigkeit $\|u'\|$ und dem beschraenkten (oder langsam
wachsenden) Log-Abstand-Gewicht bestimmt, ob \eqref{eq:LDI}
im Limes $T \nearrow T^*$ ueberlebt.

Die Bedingung TLL selbst ist plausibel, weil sie eine
``geometrische Regularitaet auf der Trajektorienskala'' erfasst: sie erlaubt dem
Attraktor, fraktal zu sein (was Lipschitz-Projektion ausschliesst), fordert aber,
dass die Projektionsinstabilitaet hoechstens logarithmisch waechst,
wenn sich die Trajektorie dem Attraktor naehert.
Dies ist die natuerliche Kritikalitaetsschwelle fuer dissipative PDEs,
wo log-Lipschitz-Abschaetzungen generisch aus Quetschungseigenschaften des Semiflusses
entstehen (vgl.\ Barbu--Cannone~\cite{BarbuCannone2016}).
\end{remark}

\medskip\noindent
\textbf{Erweiterte Abhaengigkeitskette.}\;
Der log-Lipschitz-Pfad bietet eine \emph{parallele} Route zu~G$'$,
unabhaengig vom AGC-Pfad:
\[
  \underbrace{\textup{TLL} + \textup{LDI}}_{\text{offen}}
  \;\overset{\text{Lem.~\ref{lem:LL-BV}}}{\Longrightarrow}\;
  \textup{BV}
  \;\Longrightarrow\;
  \textup{G}'
  \;\overset{\text{Prop.~\ref{prop:Gprime}}}{\Longrightarrow}\;
  \textup{Satz~\ref{thm:main}}.
\]
Der Vorteil gegenueber dem AGC-Pfad ist, dass TLL und LDI
\emph{quantitative Trajektorien-Ebene}-Bedingungen sind, statt
\emph{globaler geometrischer} Eigenschaften von~$\mathcal{A}$.
Ein Beweis von TLL wuerde keine Existenz inertialer Mannigfaltigkeiten erfordern;
es wuerde genuegen, log-Lipschitz-Stabilitaet der Naechste-Punkt-Projektion
in Trajektorienrichtung zu etablieren.

\begin{proposition}[Squeezing impliziert TLL]\label{prop:squeezing-tll}
Sei $\mathcal{A}$ der globale Attraktor einer dissipativen PDE auf einem Hilbertraum $H$ mit kompakter Einbettung $V \hookrightarrow H$. Angenommen, das System erf\"ullt die \emph{Squeezing-Eigenschaft}: es existieren $N \in \mathbb{N}$, $\theta \in (0,1)$ und $T_0 > 0$, sodass f\"ur je zwei Trajektorien $u_1(t), u_2(t)$ auf $\mathcal{A}$ gilt:
\begin{equation}\label{eq:squeezing}
  \|Q_N(u_1(T_0) - u_2(T_0))\|_H \leq \theta \|u_1(0) - u_2(0)\|_H
\end{equation}
wobei $Q_N = I - P_N$ die Projektion auf die hohen Moden ist. Dann erf\"ullt die N\"achste-Punkt-Projektion $P_\mathcal{A}: \mathcal{A}^R \to \mathcal{A}$ die TLL-Bedingung (Definition~\ref{def:tll}) mit $C_{\mathrm{TLL}} = C(N, \theta)$, das nur von $N$ und $\theta$ abh\"angt.
\end{proposition}

\begin{proof}[Beweisskizze]
Die Squeezing-Eigenschaft impliziert, dass $\mathcal{A}$ in einem Lipschitz-Graphen \"uber den ersten $N$ Fourier-Moden enthalten ist (Foias--Temam, 1984). Auf einem solchen Graphen erf\"ullt die N\"achste-Punkt-Projektion
$\|P_\mathcal{A}(x) - P_\mathcal{A}(y)\| \leq (1 + L_N) \|P_N(x - y)\| + L_N \|Q_N(x - y)\|$,
wobei $L_N$ die Lipschitz-Konstante des Graphen ist. F\"ur Punkte $x$ im Abstand $d$ von $\mathcal{A}$ ergibt sich der logarithmische Faktor aus der Kegelbedingung: der Hochmoden-Fehler $\|Q_N(x - P_\mathcal{A}(x))\|$ wird durch $d \cdot (1 + \log(R/d))$ kontrolliert, wenn $\mathcal{A}$ endliche fraktale Dimension hat (Man\'e-Projektionssatz). Zusammengenommen ergibt dies TLL mit $C_{\mathrm{TLL}} = (1 + L_N)(1 + \dim_F(\mathcal{A})/N)$.
\end{proof}

\begin{remark}[Relevanz f\"ur 3D-Navier--Stokes]\label{rem:squeezing-ns}
Die Squeezing-Eigenschaft ist f\"ur 2D-Navier--Stokes (Foias--Temam 1984), Reaktions-Diffusions-Systeme und ged\"ampfte Wellengleichungen etabliert. F\"ur 3D-Navier--Stokes w\"urde Squeezing auf dem Attraktor aus der Eigenschaft der \emph{bestimmenden Moden} folgen, die f\"ur generische Anregung bekannt ist (Foias--Prodi 1967, Jones--Titi 1992). Die L\"ucke besteht darin, dass bestimmende Moden das asymptotische Verhalten kontrollieren, w\"ahrend Squeezing gleichm\"a\ss ige Kontraktion in der Zeit erfordert. Proposition~\ref{prop:squeezing-tll} zeigt, dass das Schlie\ss en dieser L\"ucke sofort Assumption~G \"uber den TLL/LDI-Pfad liefern w\"urde.
\end{remark}

\subsubsection*{Schwellenaxiom und strukturelle Dichotomie}

\begin{axiom}[Inertiale Mannigfaltigkeit -- G-NS]\label{ax:g-ns}
Es existiert eine endlichdimensionale glatte Mannigfaltigkeit $\mathcal{M} \subset H$ (der Phasenraum), invariant unter dem Navier--Stokes-Fluss, sodass:
\begin{enumerate}
\item[(G1)] \textbf{Exponentielle Attraktion:} F\"ur jede Leray--Hopf-L\"osung $u(t)$ existiert $v(t) \in \mathcal{M}$ mit $\|u(t) - v(t)\|_H \leq C e^{-\alpha t}$.
\item[(G2)] \textbf{Slaving (Modenunterordnung):} Die hohen Moden $Q_N u$ sind glatte Funktionen der niedrigen Moden $P_N u$ auf $\mathcal{M}$: es existiert eine Lipschitz-Abbildung $\Phi: P_N H \to Q_N H$ mit $\mathcal{M} = \mathrm{graph}(\Phi)$.
\item[(G3)] \textbf{Uniformit\"at:} Die Konstanten $C$, $\alpha$ und die Lipschitz-Konstante von $\Phi$ h\"angen nur von den physikalischen Parametern $(\nu, f, \Omega)$ ab, nicht von den Anfangsdaten.
\end{enumerate}
\end{axiom}

\begin{proposition}[\"Aquivalenzen von G-NS]\label{prop:g-equivalences}
Die folgenden Aussagen sind \"aquivalent:
\begin{enumerate}
\item[(i)] Axiom~\ref{ax:g-ns} gilt (inertiale Mannigfaltigkeit existiert).
\item[(ii)] \textbf{Spektrall\"uckenbedingung:} Es existiert $N$ mit $\lambda_{N+1} - \lambda_N \geq C_{\mathrm{NL}} \cdot \lambda_N^{1/2}$, wobei $\lambda_j$ die Eigenwerte des Stokes-Operators sind und $C_{\mathrm{NL}}$ die Nichtlinearit\"at kontrolliert.
\item[(iii)] \textbf{Starkes Squeezing:} Die Squeezing-Eigenschaft (Proposition~\ref{prop:squeezing-tll}) gilt gleichm\"a\ss ig f\"ur alle L\"osungen auf dem Attraktor.
\item[(iv)] \textbf{Assumption G} (wie in diesem Paper definiert): die N\"achste-Punkt-Selektion $v(t) \in \mathcal{A}$ ist absolut stetig mit integrierbaren Gronwall-Koeffizienten.
\end{enumerate}
Die Implikation (i)$\Rightarrow$(iv) ist klassisch; (iv)$\Rightarrow$(i) folgt aus der Glattheit des Attraktors und der durch die BV-Regularit\"at erzwungenen Endlichdimensionalit\"at.
\end{proposition}

\begin{remark}[Der Tail-Scanner]\label{rem:tail-scanner}
Um zu lokalisieren, wo G-NS bricht, werden drei diagnostische Gr\"o\ss en untersucht:
\begin{enumerate}
\item \textbf{Spektraldichte:} Die Stokes-Eigenwerte erf\"ullen $\lambda_j \sim c\, j^{2/3}$ auf $\mathbb{T}^3$. Die Spektrall\"uckenbedingung (ii) erfordert $\lambda_{N+1} - \lambda_N \gtrsim \lambda_N^{1/2}$, was versagt, wenn Eigenwertcluster zu dicht werden. F\"ur 3D ist diese L\"ucke grenzwertig: $\lambda_{N+1} - \lambda_N \sim N^{-1/3}$, w\"ahrend $\lambda_N^{1/2} \sim N^{1/3}$.
\item \textbf{Energietransfer:} Falls Energie dauerhaft im Hochfrequenz-Tail verbleibt ($\sum_{j > N} E_j \not\to 0$ exponentiell), versagt die Slaving-Bedingung (G2). DNS-Evidenz legt exponentiellen Tail-Abfall f\"ur 3D NS bei moderaten Reynoldszahlen nahe.
\item \textbf{Defektes Slaving:} Falls $\Phi$ nur H\"older-stetig ist (nicht Lipschitz), degradiert die BV-Regularit\"at zu H\"older-Regularit\"at -- unzureichend f\"ur Assumption~G, aber ausreichend f\"ur approximative inertiale Mannigfaltigkeiten (Foias--Sell--Temam).
\end{enumerate}
\end{remark}

\begin{remark}[No-Go-Mechanismen]\label{rem:g-nogo}
G-NS versagt, falls:
\begin{enumerate}
\item \textbf{Spektrale H\"aufung:} Eigenwertl\"ucken schlie\ss en sich zu schnell, sodass die lineare Dissipation die nichtlineare Kopplung nicht dominieren kann. Dies ist der Grenzfall f\"ur 3D NS ($\lambda_j \sim j^{2/3}$, verglichen mit $\lambda_j \sim j$ f\"ur 2D, wo inertiale Mannigfaltigkeiten existieren).
\item \textbf{Intermittierende Hochfrequenz-Bursts:} Kurzlebige, aber wiederkehrende Energieexkursionen im Tail zerst\"oren die exponentielle Attraktion -- ein klarer dynamischer Mechanismus gegen G-NS.
\item \textbf{Attraktor-Rauheit:} Falls der Attraktor fraktale Dimension nahe der Einbettungsdimension hat, kann keine glatte Mannigfaltigkeit ihn enthalten -- approximative inertiale Mannigfaltigkeiten sind das bestm\"ogliche Resultat.
\end{enumerate}
\end{remark}

\begin{remark}[Die Tail-Dominanz-Ungleichung]\label{rem:tail-dominance}
Der strukturelle Gehalt von Axiom~\ref{ax:g-ns} reduziert sich auf eine einzige quantitative Ungleichung: \emph{lineare Dissipation muss die nichtlineare Kopplung im Hochfrequenz-Tail dominieren}. Explizit: Fuer die Stokes-Eigenwerte $\lambda_j \sim c\, j^{2/3}$ auf $\mathbb{T}^3$ erfordert die Slaving-Bedingung (G2)
\[
  \nu \lambda_{N+1} \geq C_{\mathrm{NL}} \cdot \sup_{u \in \mathcal{A}} \|B(P_N u, Q_N u)\|,
\]
wobei $B$ die Bilinearform der Navier--Stokes-Nichtlinearitaet ist. Die Grenzskalierung $\lambda_{N+1} - \lambda_N \sim N^{-1/3}$ gegenueber $\lambda_N^{1/2} \sim N^{1/3}$ macht 3D zur \emph{kritischen Dimension} fuer inertiale Mannigfaltigkeiten: In 2D ($\lambda_j \sim j$, Luecke $\sim 1$) gilt Slaving; in 3D schliesst sich die Luecke asymptotisch. Dies ist kein ,,fehlender Trick'', sondern eine \emph{strukturelle Grenzlinie} -- die Navier--Stokes-Nichtlinearitaet ist genau stark genug, um mit der Dissipation zu konkurrieren.
\end{remark}

\begin{lemma}[G-NS impliziert globale Regularit\"at]\label{lem:g-regularity}
Falls Axiom~\ref{ax:g-ns} gilt, dann ist jede Leray--Hopf-L\"osung glatt f\"ur $t > 0$ und der globale Attraktor $\mathcal{A}$ ist eine glatte kompakte endlichdimensionale Mannigfaltigkeit. Insbesondere ist Blow-up in endlicher Zeit ausgeschlossen.
\end{lemma}

\bigskip

\begin{remark}[Status von Assumption~G]\label{rem:assumption-G}
Assumption~G ist \emph{keine} Konsequenz der Standard-Attraktortheorie fuer die
3D-Navier--Stokes-Gleichungen. Die Existenz eines kompakten globalen Attraktors
$\mathcal{A} \subset H^1(\Omega)$ ist klassisch (vgl.\
\cite{FoiasManleyRosaTemam2001,Temam2001}), und jedes $v \in \mathcal{A}$ ist glatt
nach Proposition~\ref{prop:attractor-regularity}. Jedoch haengt die \emph{Zeitregularitaet}
der Abbildung $t \mapsto v(t) = \arg\min_{v \in \mathcal{A}} \|u(t) - v\|_{L^2}$
von der Geometrie von $\mathcal{A}$ und der Regularitaet der Trajektorie $u(t)$
auf nichttriviale Weise ab.

Falls $\mathcal{A}$ eine Inertiale Mannigfaltigkeit zulaesst (d.h.\ eine endlichdimensionale,
Lipschitz-invariante Mannigfaltigkeit, die alle Orbits exponentiell anzieht), dann ist die
Naechste-Punkt-Projektion Lipschitz-stetig und (G1)--(G2) folgen.
Die Existenz inertialer Mannigfaltigkeiten fuer die 3D-Navier--Stokes-Gleichungen auf
$\mathbb{T}^3$ ist selbst ein offenes Problem, obwohl sie fuer mehrere
dissipative PDE-Klassen, die eine Spektrallueckenbedingung erfuellen, etabliert ist
\cite{EdenFoiasNicolaenko1994}. Insbesondere ist fuer hyperviskose
Modifikationen $(-\Delta)^{\beta}$ mit $\beta \geq 5/4$ die Spektrallueckenbedingung erfuellt und inertiale Mannigfaltigkeiten existieren~\cite{MalletParetSell1988},
sodass Assumption~G automatisch gilt. Der klassische Fall $\beta = 1$ bleibt
offen, gerade weil die Eigenwertluecken des Stokes-Operators in drei Dimensionen
zu langsam wachsen, um die fuer die Versklavung hochfrequenter Moden erforderliche
spektrale Trennung zu garantieren. Juengere Resultate zur $H^2$-Regularitaet fuer Navier--Stokes
mit ``perfect slip''-Randbedingungen (2024/2025) deuten darauf hin, dass die
geometrische Einschraenkung aufweichbar sein koennte, obwohl die hier relevante
no-slip-Situation ungeloest bleibt.
Zur log-Lipschitz-Regularitaet des Navier--Stokes-Semiflusses und deren
Implikationen fuer Lipschitz-artige Schranken auf Attraktorprojektionen
siehe~\cite{BarbuCannone2016,MalletParetSell1988}.
Eine umfassende moderne Behandlung von Squeezing-Eigenschaften, Man\'e-Projektionen
und Attraktordimensionen fuer dissipative PDEs liefert
Zelik~\cite{Zelik2022}.

Diese Luecke zu schliessen -- Assumption~G aus der intrinsischen Struktur des
Navier--Stokes-Attraktors zu beweisen -- ist das zentrale offene Problem des vorliegenden Ansatzes.
Siehe Abschnitt~5 (Einschraenkungen) fuer eine detaillierte Diskussion und moegliche Strategien
mit log-Lipschitz-Regularitaet und Null-Lipschitz-Abweichung.

\medskip\noindent
\textbf{Resolventen-architektonische Umformulierung.}\;
Die spektrale Sektorzerlegung aus~\cite{GeigerRH2026c}
motiviert einen heuristischen Ansatz fuer Assumption~G. Fuehre \emph{Sektorlabels}
$\{S_j\}_{j=1}^{N}$ fuer die Galerkin-Moden von $v(t)$ ein, wobei jeder Sektor
$S_j$ einem Spektralband $[2^{j-1}, 2^j)$ des Stokes-Operators entspricht. Die Schluesselerkenntnis ist, dass die Lipschitz-Konstante von
$t \mapsto v(t)$ von \emph{Sektoren-Kreuzungsereignissen} abhaengt: wenn der
Minimierer von einem Attraktor-Ast zu einem anderen springt, entkoppeln die Moden in
verschiedenen Sektoren mit Raten, die durch die Spektralluecke des
Stokes-Operators kontrolliert werden. Umformulierung von Assumption~(G1) in dieser Sprache:
die projektive Familie $P_N(v(t)) = \sum_{j=1}^{N} \Pi_j v(t)$ (wobei
$\Pi_j$ auf Sektor $S_j$ projiziert) bleibt gleichmaessig beschraenkt in einer
endlich-rangigen Approximation. Sektoren-Kreuzungsereignisse -- bei denen der Minimierer
zwischen Aesten wechselt, die durch verschiedene Sektorlabels getrennt sind -- erzwingen
hoehere Lipschitz-Konstanten, aber Kompaktheit von $\mathcal{A}$ kontrolliert
die Haeufigkeit solcher Kreuzungen. Falls der Attraktor ein
$\varepsilon$-Netz mit $N(\varepsilon) \leq C\,\varepsilon^{-d}$
Elementen zulaesst (wobei $d$ die fraktale Dimension ist), dann liefert eine naive Schranke fuer
die Totalvariation von $v(t)$
$C\,N(\varepsilon)\,\varepsilon = C\,\varepsilon^{1-d}$, was
fuer $\varepsilon \to 0$ bei $d > 1$ divergiert.  Eine raffiniertere Analyse
muesste die zeitliche Kohaerenz der Minimierer-Trajektorie ausnutzen
(der Minimierer besucht nicht alle $\varepsilon$-Netz-Punkte, sondern nur jene
entlang des dynamischen Orbits), was eine bessere Schranke liefern koennte.
Dies bleibt eine offene Frage.

\medskip\noindent
\textbf{Morse-theoretische Umformulierung.}\;
Assumption~G kann in der Sprache der Dynamischen-Systeme-Geometrie
umformuliert werden: $H(u \mid \mathcal{A})$ definiert eine Ljapunow-Funktion auf dem
Phasenraum, und Assumption~G besagt, dass $\mathcal{A}$ ein
\emph{stratifizierter Morse-Komplex} fuer diese Funktion ist -- alle kritischen
Punkte von $H|_{\mathcal{A}}$ sind isoliert und der Gradientenfluss von~$H$
ist transversal zu den Strata. In dieser Umformulierung wird (G1) zu:
der Minimierer $v(t)$ kreuzt nicht zwischen Morse-Zellen
transversal zu~$\mathcal{A}$, und (G2) wird zu: die Morse-Indizes
(Dimensionen instabiler Mannigfaltigkeiten an kritischen Punkten) sind gleichmaessig
beschraenkt. Die Conley-Index-Theorie bietet ein Rahmenwerk zum Beweis
solcher Eigenschaften ohne direkte PDE-Abschaetzungen, was Assumption~G
auf eine topologische Aussage ueber die Zellstruktur des Attraktors reduziert.

\medskip\noindent
\textbf{Helizitaetskopplung.}\;
Eine Verletzung von (G2) -- unbeschraenkte Gronwall-Koeffizienten durch
schnelles Minimierer-Umschalten -- wuerde unbeschraenktes Wachstum der
Helizitaet $\mathcal{H} = \int u \cdot \omega\,dx$ entlang der Trajektorie erzeugen.
Dies liegt daran, dass Minimierer-Spruenge zwischen topologisch verschiedenen
Attraktor-Aesten Rotationsumordnungen beinhalten, die
Helizitaet injizieren. Falls man zeigen kann, dass das Helizitaetswachstum durch die
Enstrophie-Dissipation beschraenkt ist (eine Konsequenz der NS-Gleichungen in 3D:
$\frac{d}{dt}\mathcal{H} = -2\nu\int\omega\cdot\nabla\times\omega\,dx$),
dann wird Assumption~G aequivalent zu ``keine Helizitaetskaskade'':
der Helizitaetsfluss ueber Skalen bleibt beschraenkt. Dies wuerde das
bedingte Rahmenwerk mit einer physikalisch messbaren und
experimentell falsifizierbaren Bedingung verbinden.
\end{remark}

% ------------------------------------------------------------------
\subsection*{Hauptsatz}
% ------------------------------------------------------------------
\begin{theorem}[Bedingter Ausschluss von Singularitaeten in endlicher Zeit]
\label{thm:main}
Sei $\Omega = \mathbb{T}^3$ und sei $u(t)$ eine Leray--Hopf-Loesung der
inkompressiblen Navier--Stokes-Gleichungen mit glatten Anfangsdaten
$u_0 \in H^1(\Omega)$. Sei $\mathcal{A} \subset H^1(\Omega)$ der kompakte globale
Leray--Hopf-Attraktor, und definiere
\[
    H\!\left(u\mid\mathcal{A}\right)
    \;=\; \inf_{v\in\mathcal{A}}\frac{1}{2}\int_\Omega |u-v|^2\,dx.
\]
Nach Proposition~\ref{prop:attractor-regularity} erfuellt jedes $v \in \mathcal{A}$
$v \in C^\infty(\Omega)$ und $\sup_{v\in\mathcal{A}}\|v\|_{W^{1,\infty}} < \infty$.

\textbf{Falls Assumption~G \textup{(}oder die schwaechere
Assumption~G$'$, Definition~\ref{ass:Gprime}\textup{)} gilt},
und falls zusaetzlich die folgende \textbf{Dissipations-Dominanz-Bedingung~(D)} erfuellt ist: es existiert ein Blow-up-Szenario, in dem
\begin{equation}\label{eq:diss-dom}
  \frac{\nu}{2}\int_0^T \|\nabla u\|_{L^2}^2\,dt
  \;>\;
  \int_0^T \bigl(C_1\,H(t) + C_2\bigr)\,dt
  \quad \text{fuer ein } T < T^*,
  \tag{D}
\end{equation}
dann existiert keine endliche Blow-up-Zeit $T^* < \infty$. Insbesondere ist
$u(t)$ global glatt.
\end{theorem}

\begin{remark}[Die Dissipationsluecke]\label{rem:diss-gap}
Bedingung~\eqref{eq:diss-dom} ist das eigentliche Hindernis, das
das vorliegende bedingte Rahmenwerk von einem unbedingten Beweis trennt.
Die Leray-Energieidentitaet garantiert
$\int_0^{T^*}\|\nabla u\|^2 < \infty$ (selbst unter Blow-up), sodass
die linke Seite von~\eqref{eq:diss-dom} \emph{endlich} ist.
Die rechte Seite (Gronwall-Kosten) ist unter Assumption~G ebenfalls endlich.
Ob die Dissipation die Gronwall-Kosten dominiert -- d.h.\ ob~\eqref{eq:diss-dom} gilt -- haengt von den \emph{relativen Raten} der beiden Terme nahe $T^*$ ab.

Drei Strategien zur Etablierung von~\eqref{eq:diss-dom} sind identifiziert:
\begin{enumerate}
  \item[\textup{(i)}]  Das Funktional $H$ auf $H^1$-Ebene heben,
    wobei $\|\nabla u\|_{L^2}^2$ durch $\|\nabla\omega\|_{L^2}^2$
    ersetzt wird (das unter Blow-up divergieren kann).
  \item[\textup{(ii)}]  Quantitative Blow-up-Raten ausnutzen, um
    zu zeigen, dass $\|\nabla u(t)\|^2$ nahe $T^*$ schnell genug waechst,
    um den Gronwall-Faktor zu uebertreffen.
  \item[\textup{(iii)}]  Serrin-artige Integrierbarkeitskriterien
    zur unabhaengigen Kontrolle der Gronwall-Kosten verwenden.
\end{enumerate}
Jede davon ist ein offenes Problem. Die vorliegende Arbeit etabliert
das bedingte Rahmenwerk; die Schliessung der Dissipationsluecke wird
auf zukuenftige Arbeiten verschoben.

\medskip\noindent
\textbf{Quantitative Analyse von Bedingung~(D).}\;
Unter Assumption~G erfuellen die Gronwall-Koeffizienten
$C_1 \le C_1^* := 2\sup_{v\in\mathcal{A}}\|\nabla v\|_{L^\infty}+1$
und $C_2 \le C_2^*$ (beide Konstanten unabhaengig von~$t$).
Nach dem Standard-Gronwall-Lemma gilt
$H(t) \le [H(0)+C_2^*\,T^*]\,e^{C_1^*\,T^*}$
auf $[0,T^*)$. Daher erfuellt die rechte Seite
von~\eqref{eq:diss-dom}
$\int_0^T(C_1\,H+C_2)\,dt
\le C_1^*[H(0)+C_2^*T^*]\,T^*\,e^{C_1^*T^*} + C_2^*T^*$,
waehrend die linke Seite durch die Leray-Energieidentitaet beschraenkt ist:
$\frac{\nu}{2}\int_0^{T^*}\|\nabla u\|^2\,dt
\le \frac{1}{2}\|u_0\|^2_{L^2}
+ \|f\|_{L^2}\,\|u\|_{L^\infty_tL^2}\,T^*$.
Bedingung~(D) wird damit zu einem quantitativen Vergleich zwischen
dem Anfangsenergiebudget und dem exponentiell verstaerkten Attraktorabstand.
Insbesondere ist fuer kleines $T^*$ (schneller Blow-up)
$e^{C_1^*T^*} \approx 1$ und die Bedingung am leichtesten
erfuellt; fuer grosses~$T^*$ (langsamer Blow-up nahe dem Attraktor)
waechst der exponentielle Gronwall-Faktor und die Bedingung wird
restriktiver. Das $H^1$-Lift-Programm
(Abschnitt~\ref{sec:H1-lift}) umgeht diese Einschraenkung vollstaendig:
die Enstrophie-Dissipation $\int\|\Delta u\|^2$ hat
keine a-priori-Oberschranke und kann jede endliche Gronwall-Kosten dominieren.
\end{remark}

\begin{proof}
Das Argument verlaeuft in drei Schritten.

\medskip
\noindent\textbf{Schritt 1 (BKM-Kriterium).} Nehme zum Widerspruch an, dass eine endliche
Blow-up-Zeit $T^* < \infty$ existiert. Nach dem Beale--Kato--Majda-Satz
\cite{BealeKatoMajda1984} ist dies aequivalent zu
\begin{equation}\label{eq:BKM}
    \int_0^{T^*}\|\omega(t)\|_{L^\infty}\,dt = \infty.
\end{equation}

\medskip
\noindent\textbf{Schritt 2 (Punktweiser Enstrophie-Blow-up).} Lemma~\ref{lem:enstrophy}
(angewandt auf das hypothetische Blow-up-Szenario \eqref{eq:BKM}) liefert
\begin{equation}\label{eq:diss-inf}
    \limsup_{t \nearrow T^*}\|\nabla u(t)\|_{L^2}^2 = \infty.
\end{equation}
Hinweis: Das \emph{Zeitintegral} $\int_0^{T^*}\|\nabla u\|^2\,dt$
bleibt endlich nach der Leray-Energieidentitaet
(Bemerkung~\ref{rem:divergence-landscape}).

\medskip
\noindent\textbf{Schritt 3 (Widerspruch via $H$-Schranke).}
Nach Lemma~\ref{lem:entropy-bound} erfuellt fuer jede messbare Auswahl
$v(t)\in\mathcal{A}$ (die existiert wegen Kompaktheit von $\mathcal{A}$ in $L^2$ und
Messbarkeit von $t\mapsto H(u(t)\mid\mathcal{A})$; siehe Bemerkung~\ref{rem:selection}
unten) das Funktional $E_v(t)=\tfrac12\|u(t)-v(t)\|_{L^2}^2$
\[
    \frac{d}{dt}E_v(t)
    \;\le\; -\frac{\nu}{2}\|\nabla u(t)\|_{L^2}^2 + C_1\,E_v(t) + C_2.
\]
Wahl von $v(t)$ als minimierende Auswahl (d.h.\ $E_v(t) = H(u(t)\mid\mathcal{A})$)
und Grenzuebergang $\varepsilon\to 0$ in der $\varepsilon$-approximativen Auswahl
liefert dieselbe Ungleichung fuer $H(t):=H(u(t)\mid\mathcal{A})$ im schwachen Sinne.
Integration ueber $[0,T]$ fuer $T < T^*$:
\begin{equation}\label{eq:H-integral}
    H(T) - H(0)
    \;\le\;
    -\frac{\nu}{2}\int_0^T\|\nabla u\|_{L^2}^2\,dt
    + \int_0^T\!\big(C_1\,H(t) + C_2\big)\,dt.
\end{equation}
Die rechte Seite von \eqref{eq:H-integral} hat zwei konkurrierende Terme:
die Dissipation $-\tfrac{\nu}{2}\int_0^T\|\nabla u\|^2\,dt$
(negativ, treibt $H$ nach unten) und die Gronwall-Kosten
$\int_0^T(C_1\,H + C_2)\,dt$ (positiv, erlauben $H$ zu wachsen).
Beide Terme sind \emph{endlich} fuer jedes $T < T^*$ (das Dissipationsintegral
ist beschraenkt durch die Leray-Energieidentitaet; die Gronwall-Kosten
sind beschraenkt unter Assumption~G).

Falls Bedingung~\eqref{eq:diss-dom} gilt -- d.h.\ die Dissipation
dominiert die Gronwall-Kosten fuer ein $T$ hinreichend nahe an
$T^*$ -- dann gilt $H(T) < 0$, im Widerspruch zu $H \ge 0$.

Folglich existiert unter Assumptions~G (oder~G$'$) und~(D) kein endliches $T^*$
und $u(t)$ ist global glatt.

\medskip
\noindent\textbf{BV-Modifikation (Assumption~G$'$).}\;
Unter der schwaecheren Assumption~G$'$ ist der Minimierer $v(t)$ BV statt
AC. Die einzige Stelle, an der absolute Stetigkeit von $v$ in den
Beweis eingeht, ist Lemma~\ref{lem:entropy-bound}: Der Koeffizient
$C_2(t)$ enthaelt den Term $\tfrac{1}{2}\|f - \partial_t v\|_{L^2}^2$,
der erfordert, dass $\partial_t v$ punktweise f.u.\ existiert. Unter~G$'$
ist $\partial_t v$ ein Mass (keine Funktion), und dieser Term ist
schlecht definiert.

\emph{Aufloesung.}\;
Kehre zur Rechnung zurueck, die zu Lemma~\ref{lem:entropy-bound} fuehrt.
Der Term $\langle w, \partial_t v\rangle$ entsteht aus
$\frac{\mathrm{d}}{\mathrm{d}t}\tfrac{1}{2}\|u-v\|^2
= \langle w, \partial_t u\rangle - \langle w, \partial_t v\rangle$.
Unter~G$'$ ersetze den zweiten Summanden durch sein Stieltjes-Integral:
$\int_0^T \langle w(t), \mathrm{d}v(t)\rangle$, das beschraenkt ist durch
\begin{equation}\label{eq:stieltjes-bound}
  \biggl|\int_0^T \langle w(t),\, \mathrm{d}v(t)\rangle\biggr|
  \;\le\;
  \sup_{t \in [0,T]} \|w(t)\|_{L^2}
  \;\cdot\;
  \mathrm{Var}_{L^2}(v;\, [0,T]).
\end{equation}
Die modifizierte Energieungleichung wird zu
\begin{equation}\label{eq:H-integral-BV}
  H(T) \;\le\; H(0)
  \;-\; \frac{\nu}{2}\int_0^T \|\nabla u\|_{L^2}^2\,\mathrm{d}t
  \;+\; \int_0^T \widetilde{C}_1\,H(t)\,\mathrm{d}t
  \;+\; \widetilde{C}_2\,T
  \;+\; \sqrt{2\,\sup H}\;\cdot\;
    \mathrm{Var}_{L^2}(v;\,[0,T]),
\end{equation}
wobei $\widetilde{C}_1 = 2\sup_{v \in \mathcal{A}}
\|\nabla v\|_{L^\infty} + 1$ und $\widetilde{C}_2 =
\tfrac{1}{2}\sup_{v \in \mathcal{A}}\|(v\cdot\nabla)v\|_{L^2}^2
+ \tfrac{1}{2}\|f\|_{L^2}^2 + \tfrac{\nu}{2}\sup_{v \in \mathcal{A}}
\|\nabla v\|_{L^2}^2$ beide endlich sind nach
Proposition~\ref{prop:attractor-regularity}. Man beachte, dass $\widetilde{C}_2$
$\|\partial_t v\|$ \emph{nicht enthaelt} -- der Beitrag der Minimierer-Geschwindigkeit
wurde in die Stieltjes-Schranke~\eqref{eq:stieltjes-bound} absorbiert.

Unter Blow-up ($T \nearrow T^* < \infty$):
Lemma~\ref{lem:enstrophy} gibt die \emph{punktweise} Divergenz
$\|\nabla u(t)\|_{L^2}^2 \to \infty$ fuer $t \nearrow T^*$, aber die
Leray-Energieidentitaet sichert die Endlichkeit des \emph{Zeitintegrals}
$\int_0^{T^*}\|\nabla u\|_{L^2}^2\,dt$
(Bemerkung~\ref{rem:divergence-landscape}).

\emph{Beschraenktheit von $\sup H$.}\;
Das implizite Auftreten von $\sup_{[0,T]} H$ in~\eqref{eq:H-integral-BV}
erfordert eine Begruendung. Nach der Young-Ungleichung gilt
$\sqrt{2\sup H}\cdot\mathrm{Var}(v)
\le \varepsilon\,\sup_{[0,T]} H + \mathrm{Var}(v)^2/(2\varepsilon)$
fuer beliebiges $\varepsilon > 0$.
Mit $\varepsilon = 1$ und $M(T) := \sup_{[0,T]} H$ liefert
die Ungleichung~\eqref{eq:H-integral-BV} fuer jedes $t \le T$:
$H(t) \le H(0) + \int_0^t \widetilde{C}_1\,H\,d\tau
+ (\widetilde{C}_2+1)\,T^* + \mathrm{Var}(v)^2/2$.
Das Standard-Gronwall-Lemma ergibt dann
$M(T) \le [H(0) + (\widetilde{C}_2+1)\,T^*
+ \mathrm{Var}(v)^2/2]\,e^{\widetilde{C}_1\,T^*}$,
was auf $[0,T^*)$ endlich ist unter~G$'$.

Alle anderen Terme auf der rechten Seite von
\eqref{eq:H-integral-BV} sind ebenfalls endlich
($\widetilde{C}_1, \widetilde{C}_2$ beschraenkt;
$\mathrm{Var}(v) < \infty$ nach~G$'$).
Daher ist der Widerspruch $H(T) < 0$ \emph{nicht} automatisch aus der
BV-Struktur allein: er erfordert zusaetzlich, dass der Dissipationsterm
die Gronwall-Kosten dominiert, d.h.\ Bedingung~\eqref{eq:diss-dom}.
Unter Assumptions~G$'$ und~(D) wird die rechte Seite
von~\eqref{eq:H-integral-BV} negativ fuer $T$ hinreichend nahe
an~$T^*$, was $H(T) < 0$ ergibt -- Widerspruch.
\end{proof}

\begin{remark}[Abhaengigkeiten und Fehlermodi]\label{rem:dependencies}
\mbox{}

\begin{center}
\begin{tabular}{llll}
\hline
\textbf{Label} & \textbf{Aussage} & \textbf{Typ} & \textbf{Fehlermodus} \\
\hline
Prop.~\ref{prop:attractor-regularity} &
  $\mathcal{A}$ kompakt, glatt &
  Satz & Keiner (klassisch) \\
Lem.~\ref{lem:BS} &
  Biot--Savart-Abschaetzung &
  Satz & Keiner \\
Lem.~\ref{lem:enstrophy} &
  Enstrophie-Bilanz &
  Satz & Keiner \\
Lem.~\ref{lem:entropy-bound} &
  $H$-Schranke (AC-Version) &
  Satz & Erfordert $\partial_t v \in L^2$ \\
Gl.~\eqref{eq:H-integral-BV} &
  $H$-Schranke (BV-Version) &
  \textbf{Satz} & Erfordert $\mathrm{Var}(v) < \infty$ \\
AGC2 &
  Exponentielles Tracking &
  Satz & Keiner~\cite{Temam2001} \\
\textbf{AGC1} &
  \textbf{Positiver Reach} &
  \textbf{Offen} &
  Scheitert falls $\mathcal{A}$ Spitzen hat \\
Lem.~\ref{lem:AGC-implies-G} &
  AGC $\Rightarrow$ G &
  Satz & Keiner (Lipschitz-1) \\
\hline
\end{tabular}
\end{center}

\medskip\noindent
Die logische Kette ist:
$\textup{AGC1 (offen)} + \textup{AGC2 (bewiesen)}
\;\Longrightarrow\;
\textup{G}
\;\Longrightarrow\;
\textup{G}'
\;\Longrightarrow\;
\textup{Satz~\ref{thm:main}}$.
Jedes Glied ausser AGC1 ist rigoros bewiesen. Die gesamte Regularitaetsfrage
konzentriert sich in einer einzigen geometrischen Eigenschaft des
Attraktors: \emph{hat $\mathcal{A}$ positiven Reach in $L^2$?}
Dies ist unabhaengig von der spezifischen Trajektorie~$u(t)$ und ist
aequivalent zur Existenz einer Inertialen Mannigfaltigkeit nahe~$\mathcal{A}$.
\end{remark}

\begin{remark}[Universelles Resolventenmuster]\label{rem:universal-resolvent}
Das bedingte Regularitaetsresultat von Satz~\ref{thm:main} ist eine Instanz des
\emph{Second-Order Resolvent Dominance Pattern}, das ueber alle fuenf
Probleme im FST-Programm identifiziert wurde~\cite{FST-Unified2026}. Das Muster manifestiert sich als
dreistufige Hierarchie:

\noindent
Im Navier--Stokes-Fall ist die fuehrende Ordnung die Leray--Hopf-Energieungleichung,
die Beschraenktheit von $\|u\|_{L^2}$ liefert, aber keine Regularitaet.
Die Attraktorprojektion erster Ordnung identifiziert den naechsten glatten
Referenzzustand $v(t) \in \mathcal{A}$, aber dies allein kann
$\|\nabla u\|_{L^\infty}$ nicht kontrollieren. Die Regularitaetsluecke schliesst sich erst in zweiter
Ordnung, durch die $H$-Metrik-Kontraktion: die viskose Dissipation
treibt den Attraktorabstand $H(u \mid \mathcal{A})$ in Richtung eines
Widerspruchs zu seiner Nichtnegativitaet.  Diese Hierarchie wird im
breiteren Kontext des FST-Programms in~\cite{FST-Unified2026} diskutiert.
\end{remark}

\begin{remark}[Spektralstruktur der Minimierer-Instabilitaet]
\label{rem:resolvent-minimiser}
Die instabilen Richtungen des Abstandsfunktionals $E_v = \frac{1}{2}\|u-v\|^2$
am Minimierer $v(t) \in \mathcal{A}$ -- jene, entlang derer
$\frac{d}{dt}E_v > 0$ -- entsprechen den Ljapunow-instabilen Moden des
Semiflusses bei~$v(t)$ und bilden einen endlichdimensionalen Unterraum
der Dimension hoechstens $\dim(\mathcal{A})$. Die Zeitableitung $\frac{d}{dt}E_v$
zerfaellt in einen dissipativen Teil (kontrolliert durch $-\nu\|\nabla u\|^2$)
und einen Gronwall-Teil (kontrolliert durch $C_1 E_v + C_2$). Kontraktion in der
$H$-Metrik ist somit ein Phaenomen \emph{zweiter Ordnung}: sie stabilisiert die
Trajektorie durch den viskosen Term, nachdem die Attraktorprojektion erster Ordnung
die Instabilitaet auf ein endlichdimensionales Residuum lokalisiert hat.
\end{remark}

\begin{remark}[Messbare Auswahl]
\label{rem:selection}
Die Existenz einer messbaren Abbildung $t\mapsto v(t)\in\mathcal{A}$, die
$H(u(t)\mid\mathcal{A}) = \tfrac12\|u(t)-v(t)\|_{L^2}^2$ realisiert, ist eine Standardkonsequenz des Kuratowski--Ryll-Nardzewski-Satzes ueber messbare Auswahl, angewandt
auf die kompaktwertige, unterhalb-halbstetige Multifunktion
$t\mapsto \arg\min_{v\in\mathcal{A}}\|u(t)-v\|_{L^2}$; siehe z.B.\ \cite{Wagner1977}, Theorem~4.1.
\end{remark}

\begin{remark}[Explizite Abhaengigkeit der Konstanten]
\label{rem:constants}
Die im Beweis auftretenden Konstanten haengen wie folgt von den Problemdaten ab.
\begin{itemize}
    \item $C_{\mathrm{BS}}$: haengt nur von $s > 3/2$ und der Geometrie von
          $\mathbb{T}^3$ ab (Lemma~\ref{lem:BS}).
    \item $C_1 = 2\sup_{v\in\mathcal{A}}\|\nabla v\|_{L^\infty} + 1$: endlich,
          da $\mathcal{A}$ ein kompakter Attraktor ist, der alle Orbits absorbiert,
          also $\sup_{v\in\mathcal{A}}\|v\|_{W^{1,\infty}} < \infty$.
    \item Unter Assumption~G (AC-Version):
          $C_2 = \tfrac12\sup_{v\in\mathcal{A}}\|(v\cdot\nabla)v\|_{L^2}^2
          + \tfrac12\|f - \partial_t v\|_{L^2}^2
          + \tfrac{\nu}{2}\sup_{v\in\mathcal{A}}\|\nabla v\|_{L^2}^2$.
          Unter der schwaecheren Assumption~G$'$ (BV-Version) ist die
          Minimierer-Geschwindigkeit $\partial_t v$ ein Mass statt einer Funktion,
          und der entsprechende Term wird durch die Stieltjes-Schranke
          \eqref{eq:stieltjes-bound} ersetzt; die modifizierte Konstante
          $\widetilde{C}_2 = \tfrac12\sup_{v\in\mathcal{A}}\|(v\cdot\nabla)v\|_{L^2}^2
          + \tfrac12\|f\|_{L^2}^2 +
          \tfrac{\nu}{2}\sup_{v\in\mathcal{A}}\|\nabla v\|_{L^2}^2$
          enthaelt $\|\partial_t v\|$ nicht.
          Beide Versionen sind endlich durch Attraktorkompaktheit und die Annahme
          $f \in L^\infty_t H^1_x$. Insbesondere haengen sowohl $C_1$ als auch $C_2$
          (bzw.\ $\widetilde{C}_2$)
          von der Antriebsamplitude $\|f\|_{L^\infty H^1}$, der
          Viskositaet $\nu$ und dem Attraktorradius ab (letzterer bestimmt durch
          $\nu$ und $\|f\|$), sind aber unabhaengig von $T$ und $u_0$.
\end{itemize}
\end{remark}

\section{Diskussion}

\subsection{Zusammenhang mit partiellen Regularitaetsresultaten}

Der Satz von Caffarelli--Kohn--Nirenberg \cite{CaffarelliKohnNirenberg1982} zeigt, dass das eindimensionale parabolische Hausdorff-Mass der singulaeren Menge jeder geeigneten schwachen Loesung Null ist. Obwohl dies eine tiefgreifende geometrische Einschraenkung ist, schliesst sie Blow-up in endlicher Zeit nicht aus: sie begrenzt lediglich die Groesse der potentiellen singulaeren Menge. Ebenso liefern die Serrin-artigen Regularitaetskriterien \cite{Leray1934,ConstantinFoias1988} hinreichende Bedingungen fuer Glattheit (z.B.\ $u \in L^p_t L^q_x$ mit $2/p + 3/q \le 1$), aber die a-priori-Verifikation dieser Bedingungen erfordert Informationen, die aus den Anfangsdaten allein nicht verfuegbar sind.

Unser Ansatz unterscheidet sich strukturell von beiden Paradigmen. Anstatt externe Integrierbarkeitsbedingungen an die Loesung zu stellen oder die Groesse potentieller Singularitaeten abzuschaetzen, nutzen wir die interne dissipative Struktur der Gleichungen aus: Jedes Blow-up-Szenario erzwingt die \emph{punktweise} Divergenz der Enstrophie $\|\nabla u(t)\|_{L^2}^2 \to \infty$ (ueber das BKM-Kriterium und die Enstrophie-Bilanz), was eine Spannung mit der Nichtnegativitaet des Attraktorabstandsfunktionals $H(u \mid \mathcal{A})$ erzeugt. Unter der zusaetzlichen Dissipations-Dominanz-Bedingung~(D) -- die verlangt, dass die kumulative viskose Dissipation die Gronwall-Kosten uebersteigt -- wird diese Spannung zum Widerspruch: $H < 0$. Das Argument operiert innerhalb des Energie-Dissipations-Rahmenwerks, mit zwei zentralen verbleibenden Bedingungen: einer Regularitaetsannahme an die Attraktorprojektion (Assumption~G) und der Dissipations-Dominanz-Bedingung~(D).

\subsection{Strukturelle Vorteile des attraktorbasierten Ansatzes}

Klassische Energiemethoden fuer die Navier--Stokes-Gleichungen (siehe z.B.\ \cite{ConstantinFoias1988,Temam2001}) erzeugen typischerweise Gronwall-artige Ungleichungen der Form
\[
    \frac{d}{dt}\|u\|_{H^1}^2 \le C\,\|u\|_{H^1}^2\,\|\nabla u\|_{L^\infty} - \nu\|\nabla u\|_{H^1}^2,
\]
wobei der Wirbelstreckungsterm auf der rechten Seite mit der viskosen Dissipation konkurriert. Die fundamentale Schwierigkeit besteht darin, dass in drei Dimensionen keine a-priori-Schranke fuer $\|\nabla u\|_{L^\infty}$ verfuegbar ist und die Gronwall-Abschaetzung lediglich eine doppelt-exponentielle Schranke liefert, die in endlicher Zeit explodieren kann.

Der attraktorbasierte Ansatz umgeht diese Sackgasse durch Verschiebung des Bezugssystems: Statt $\|u\|_{H^1}$ direkt abzuschaetzen, misst man die \emph{Abweichung} $w = u - v$ von einem glatten Referenzzustand $v \in \mathcal{A}$. Der Gronwall-Koeffizient $C_1 = 2\|\nabla v\|_{L^\infty} + 1$ ist nun \emph{beschraenkt} (da $v$ auf dem kompakten Attraktor liegt), anstelle des unkontrollierten $\|\nabla u\|_{L^\infty}$ des klassischen Ansatzes. Der Preis ist Assumption~G: man muss die zeitliche Regularitaet der Attraktorprojektion kontrollieren. Dies ist ein genuines anderes Problem als die Kontrolle von $\|\nabla u\|_{L^\infty}$ -- es ist eine geometrische Eigenschaft des Attraktors statt einer punktweisen Schranke an die Loesung.

\section{Navier-Stokes-Regularitaet als Pattern-A-Stabilitaetssystem}

Der Ausschluss von Blow-up in endlicher Zeit in den 3D-Navier-Stokes-Gleichungen wird hier als Realisierung der \textbf{Pattern-A-Stabilitaetslogik} formuliert (siehe \cite{FST-Unified2026}). Dieser Ansatz ersetzt die Suche nach globalen a-priori-Schranken fuer $\|\nabla u\|_{L^\infty}$ durch die folgende strukturelle Architektur:

\begin{enumerate}[label=\alph*)]
    \item \textbf{Analytische Rang-Endlichkeit (Null-Tail):} Jedes hypothetische Blow-up-Szenario ist auf eine endlichdimensionale Mannigfaltigkeit von Instabilitaetsmoden lokalisiert. Die Hochfrequenz-(UV-)Dissipationsskala wird strikt durch die Enstrophie-Bilanz \cite{Leray1934} kontrolliert, was einen analytischen Null-Tail sichert.
    \item \textbf{Attraktorprojektion als physikalische Eichung:} Die ``endliche Negativitaet'' (potentielle Singularitaetsrichtungen) wird neutralisiert durch Projektion der aktiven Stroemung auf den globalen Leray--Hopf-Attraktor $\mathcal{A}$. Diese Projektion fungiert als physikalische Eichung.
    \item \textbf{Eingeschraenkte Positivitaet (H-Funktional):} Die Nichtnegativitaet des Abstandsfunktionals $H(u \mid \mathcal{A})$ wirkt als Eich-Stabilisator und stellt sicher, dass keine Trajektorie die ``Dissipationskosten'' eines Blow-ups bezahlen kann, ohne das thermodynamische Budget zu verletzen.
\end{enumerate}

Somit ist globale Regularitaet die Manifestation spektraler Stabilitaet unter der Eich-Beschraenkung der Attraktorstruktur.
Das Dissipationsintegral wirkt als irreversible ``Kosten'', die jede Blow-up-Trajektorie bezahlen muss, waehrend der Attraktorabstand ein nichtnegativer ``Etat'' ist, der nicht ueberschritten werden kann. Dieser Mechanismus ist analog zur Rolle der Entropieproduktion in der Nichtgleichgewichtsthermodynamik: Die dissipative Struktur des Systems selbst verbietet Konfigurationen, die unendliche Entropieproduktion in endlicher Zeit erfordern.

\subsection{Einschraenkungen und offene Fragen}

Mehrere Aspekte der vorliegenden Arbeit laden zu weiterer Untersuchung ein:

\begin{itemize}
    \item \textbf{Assumption~G (die Gronwall-Luecke).} Die bedeutendste Einschraenkung dieser Arbeit ist die Abhaengigkeit von Assumption~G. Die Differentialungleichung fuer $H(u\mid\mathcal{A})$ im Beweis von Satz~\ref{thm:main} erfordert, dass das minimierende Attraktorelement $v(t)$ mit hinreichender Zeitregularitaet evolviert. Waehrend jedes feste $v \in \mathcal{A}$ raeumlich glatt ist, kann die \emph{Trajektorie} der Naechste-Punkt-Projektionen $t \mapsto v(t)$ Spruenge aufweisen, falls der Attraktor komplizierte Geometrie hat (z.B.\ Spitzen oder Selbstschnitte in der $L^2$-Topologie). Drei Strategien zum Schliessen dieser Luecke erscheinen am vielversprechendsten: (i)~Beweis der Existenz einer Inertialen Mannigfaltigkeit fuer 3D-Navier--Stokes auf $\mathbb{T}^3$, was die Projektion Lipschitz-stetig machen wuerde; (ii)~Etablierung von log-Lipschitz-Regularitaet der Projektion ueber Quetschungseigenschaften des Semiflusses (vgl.\ Barbu und Cannone); (iii)~Ausnutzung von Null-Lipschitz-Abweichung auf dem Attraktor, um zu zeigen, dass hochfrequente Moden deterministisch von niederfrequenten Moden versklavt werden, was $v(t)$ automatisch glatt macht.
    Unabhaengige Unterstuetzung fuer den attraktorgeometrischen Ansatz kommt aus juengsten Arbeiten zu fraktalen Sparseness-Bedingungen: Grujic und Xu~\cite{GrujicXu2025} zeigen, dass Singularitaeten topologisch ausgeschlossen werden, wenn die Regionen intensiver Vortizitaet hinreichend duenn (im Sinne fraktaler Masse) sind -- die erste Skalierungsreduktion der NS-Superkritikalitaetsluecke seit den 1960er Jahren. Ihre Sparseness-Bedingung ist strukturell kompatibel mit dem hier angenommenen endlichdimensionalen Attraktor.

    \item \textbf{Gebietsgeometrie.} Unser Beweis stuetzt sich auf die periodische Situation $\Omega = \mathbb{T}^3$, die sowohl die Existenz eines kompakten globalen Attraktors \cite{FoiasManleyRosaTemam2001} als auch das Verschwinden von Randtermen bei partieller Integration sichert. Erweiterung auf beschraenkte Gebiete mit No-Slip-Randbedingungen fuehrt Grenzschichteffekte ein und erfordert Attraktortheorie in $H^1_0(\Omega)$, die verfuegbar ist (vgl.\ \cite{Temam2001}), aber technisch aufwendiger. Der Ganzraum-Fall $\Omega = \mathbb{R}^3$ ist delikater, da der globale Attraktor im ueblichen Sinne nicht kompakt sein muss; man muesste stattdessen mit Trajektorienattraktoren oder Pullback-Attraktoren arbeiten.

    \item \textbf{Aeusserer Antrieb.} Wir haben glatten, zeitunabhaengigen (oder zumindest gleichmaessig beschraenkten) Antrieb $f \in L^\infty_t H^1_x$ angenommen. Fuer rauen oder singulaeren Antrieb koennen die gleichmaessigen Schranken fuer $C_1$ und $C_2$ versagen, und die Attraktorstruktur selbst kann sich aendern. Die Erweiterung des Ansatzes auf stochastischen Antrieb (wie in den stochastischen Navier--Stokes-Gleichungen) ist eine interessante Richtung, die ein Rahmenwerk zufaelliger Attraktoren erfordern wuerde.

    \item \textbf{Quantitative Schranken.} Waehrend unser Argument globale Regularitaet beweist, sind die resultierenden Schranken fuer $\|\nabla u\|_{L^\infty}$ nichtkonstruktiv: sie haengen vom Attraktorradius ab, der seinerseits von den Abschaetzungen der absorbierenden Kugel des Navier--Stokes-Semiflusses abhaengt. Herleitung expliziter, quantitativer Regularitaetsschranken aus dem Attraktorabstand wuerde praezisere Informationen ueber die Attraktorgeometrie erfordern, die in drei Dimensionen weitgehend unbekannt bleibt.

    \item \textbf{Regularitaet der messbaren Auswahl.} Der Beweis verwendet eine messbare Auswahl $t \mapsto v(t) \in \mathcal{A}$, um das naechste Attraktorelement zu verfolgen. Waehrend Messbarkeit fuer die Integralabschaetzungen genuegt, wuerde hoehere Regularitaet dieser Auswahl (z.B.\ absolute Stetigkeit) die Differentialungleichung verstaerken und potentiell schaerfere Abklingraten fuer $H(u \mid \mathcal{A})$ liefern.
\end{itemize}

\begin{remark}[Aufbrechen der Regularitaets-Attraktor-Zirkularitaet]
\label{rem:circularity-discussion}
Der Standardzugang zu globalen Attraktoren in der dissipativen
PDE-Theorie (Temam~\cite{Temam2001}, Zelik~\cite{Zelik2014}) setzt
Glattheit der Loesungen voraus, um Endlich-Dimensionalitaet des
Attraktors zu beweisen. Unser Framework kehrt diese logische
Richtung um: Wir setzen Endlich-Dimensionalitaet des Attraktors
voraus (die empirisch gut gestuetzt und fuer den schwachen
Leray--Hopf-Attraktor etabliert ist; vgl.\
Bemerkung~\ref{rem:circularity}) und leiten
Regularitaetskonsequenzen ab.

Die scheinbare Zirkularitaet -- Regularitaet wird fuer den Attraktor
benoetigt, und der Attraktor wird fuer Regularitaet benoetigt -- wird
durch die \emph{thermodynamische Obstruktion} aufgebrochen: Falls
Blow-up in endlicher Zeit eintraete, wuerde das Freie-Energie-Funktional
$H(u \mid \mathcal{A})$ gegen~$-\infty$ divergieren, im Widerspruch
zur dissipativen Schranke $H(u \mid \mathcal{A}) \ge 0$.
Die Nichtnegativitaet von~$H$ ist eine Konsequenz der Definition
($H$ ist ein quadrierter Abstand), nicht der Regularitaet.

Dies ist \emph{kein} Beweis unbedingter Regularitaet, sondern eine
\emph{strukturelle Reduktion}: Globale Regularitaet folgt, WENN
der Attraktor die vorausgesetzten Eigenschaften besitzt
(Endlich-Dimensionalitaet, Kompaktheit, Glattheit der Elemente).
Jede dieser Eigenschaften ist unabhaengig von der Regularitaetsfrage
etabliert -- sie folgen aus der schwachen Attraktortheorie und dem
Bootstrap-Argument fuer Elemente auf dem Attraktor
(Proposition~\ref{prop:attractor-regularity}). Die verbleibende
Luecke ist Assumption~G, die die \emph{zeitliche Regularitaet} der
Projektion betrifft, nicht die raeumliche Regularitaet der
Attraktorelemente.
\end{remark}

\begin{remark}[Zerlegung von Assumption G: unabhaengige Komponenten]
\label{rem:g-decomposition}
Um dem Einwand zu begegnen, Assumption~G koenne zirkulaer sein (implizit Regularitaet voraussetzen), zerlegen wir sie in drei logisch unabhaengige Komponenten:

\begin{enumerate}
\item[\textbf{(G$_1$)}] \textbf{Attraktorexistenz und Kompaktheit.} Der globale Attraktor $\mathcal{A}$ existiert als kompakte Teilmenge von $H^1(\Omega)$ und ist endlichdimensional ($d_H(\mathcal{A}) < \infty$). \emph{Status:} Etabliert fuer den schwachen Leray--Hopf-Attraktor ohne unbedingte Regularitaetsannahme (Foias--Temam, Sell, Chepyzhov--Vishik; vgl.\ Bemerkung~\ref{rem:circularity}). Diese Komponente ist \textbf{unabhaengig} von der Regularitaetsfrage.

\item[\textbf{(G$_2$)}] \textbf{Projektionsregularitaet.} Die Naechste-Punkt-Projektion $\pi_\mathcal{A}: H^1(\Omega) \to \mathcal{A}$ ist log-Lipschitz (oder laesst eine BV-Selektion zu) in einer tubularen Umgebung von~$\mathcal{A}$. \emph{Status:} Offen. Dies ist eine \textbf{geometrische} Eigenschaft des Attraktors, keine dynamische Regularitaetsaussage. Sie kann prinzipiell aus der metrischen Struktur des Attraktors (Hausdorff-Dimension, Tangentenkegel-Regularitaet, Assouad-Dimension) verifiziert werden, ohne auf die Regularitaet der Loesungen Bezug zu nehmen.

\item[\textbf{(G$_3$)}] \textbf{Gronwall-Integrierbarkeit.} Die zeitabhaengigen Gronwall-Koeffizienten $C_1(t), C_2(t)$ in der Differentialungleichung fuer $H(u|\mathcal{A})$ erfuellen $\int_0^T C_i(t)\,dt < \infty$. \emph{Status:} Folgt aus Standard-Energieabschaetzungen und der Absorbing-Ball-Eigenschaft, unter der Annahme, dass die Loesung in $H^1$ verbleibt. Unter Blow-up divergiert die Enstrophie $\|\nabla u\|^2$, was $C_1$ unbeschraenkt macht -- aber dies ist genau der Widerspruchsmechanismus von Satz~\ref{thm:main}.
\end{enumerate}

Der Zirkularitaetseinwand betrifft ausschliesslich (G$_2$): Beweist man zuerst Regularitaet, ist der Attraktor glatt und (G$_2$) trivial. Unser Framework bricht diese Zirkularitaet, indem es (G$_2$) als \emph{geometrische Hypothese} ueber den Attraktor behandelt (verifizierbar aus seinen metrischen Eigenschaften) statt als dynamische Regularitaetsannahme. Die Bruecke von (G$_1$) zu (G$_2$) ist der Inhalt der Begleitarbeit zur Log-Distanz-Integrierbarkeit~\cite{Geiger2026-LDI}.
\end{remark}

\begin{remark}[Reach-Lemma ueber Rueckwaertsstabilitaet]
\label{rem:reach-lemma}
Falls der Attraktor $A$ positiven Reach $\mathrm{reach}(A) > 0$ besitzt (d.\,h.\ die Naechste-Punkt-Projektion ist einwertig in einer tubularen Umgebung), dann impliziert Rueckwaertsstabilitaet (Loesungen auf $A$ sind Lipschitz-stetig in Rueckwaertszeit) kombiniert mit Injektivitaet endlichdimensionaler Projektionen $P_N: A \to \mathbb{R}^N$ die Abschaetzung $\mathrm{reach}(A) \geq c/\|P_N\|_{\mathrm{op}}$. Dies wuerde Assumption~G (Lipschitz-Regularitaet) als Konsequenz der Endlich-Dimensionalitaet etablieren und die in den Einschraenkungen genannte Zirkularitaet aufbrechen.
\end{remark}

\begin{remark}[Relaxation von Bedingung (D) zu limsup-Divergenz]
\label{rem:condition-d-limsup}
Bedingung~(D) kann von punktweiser Dominanz zu einer masstheoretischen Version relaxiert werden: Anstatt $\|\nabla u(t)\|^2 / H(u(t)|A) \to \infty$ fuer $t \to T^*$ fuer alle Blow-up-Trajektorien zu fordern, genuegt es zu verlangen
\[
\limsup_{t \to T^*} \frac{\|\nabla u(t)\|^2}{H(u(t)|A)} = +\infty.
\]
Diese schwaechere Bedingung erzwingt weiterhin den Widerspruch in Satz~3.1 (die Freie-Energie-Barriere kann nicht aufrechterhalten werden, wenn Dissipation intermittierend dominiert) und ist aus PDE-Perspektive natuerlicher: Enstrophie-Blow-up muss nicht monoton sein, nur unbeschraenkt.
\end{remark}

\subsection{Eine spieltheoretische Perspektive}

Die Regularitaetsfrage laesst eine natuerliche Umformulierung als Minimax-Problem zu. Man betrachte zwei abstrakte ``Spieler'', die die Dynamik steuern: Spieler~$D$ (Dissipation), dessen Strategie viskose Daempfung mit Rate $\nu\|\nabla u\|_{L^2}^2$ ist, und Spieler~$N$ (Nichtlinearitaet), dessen Strategie der konvektive Energietransfer ueber $(u\cdot\nabla)u$ ist. Definiere den \emph{Spielwert} auf Modenebene als
\begin{equation}
    \mathcal{V}(t) \;:=\; \inf_{\text{Moden}} \bigl[\text{Dissipationsrate} - \text{nichtlineare Transferrate}\bigr].
\end{equation}
In dieser Sprache ist globale Regularitaet aequivalent zu $\mathcal{V}(t) \ge 0$ fuer alle $t$: Dissipation dominiert die Nichtlinearitaet gleichmaessig ueber alle Fourier-Moden. Der Attraktorabstand $H(u\mid\mathcal{A})$ spielt die Rolle eines endlichen ``Budgets'', das dem System zur Verfuegung steht, waehrend das Dissipationsintegral $\int_0^T \nu\|\nabla u\|_{L^2}^2\,dt$ die kumulativen Kosten darstellt, die jede Blow-up-Trajektorie bezahlen muss. Da das Budget nichtnegativ ist, reduziert sich die Frage darauf, ob die kumulativen Dissipationskosten die Gronwall-Kosten nahe der Blow-up-Zeit uebersteigen -- genau die Dissipations-Dominanz-Bedingung~(D) von Satz~\ref{thm:main}.

Dieser Standpunkt verbindet sich mit der Legendre--Fenchel-Dualitaet in der konvexen Analysis: Interpretiert man die Dissipation als konvexes ``Entropieproduktions''-Funktional $S$ und den nichtlinearen Transfer als sein konjugiertes ``Freie-Energie-Kosten''-Funktional $F$, so ist die Regularitaetsbedingung $S \ge F$ das PDE-Analogon des zweiten Hauptsatzes der Thermodynamik -- das System kann aus der Nichtlinearitaet nicht mehr Arbeit extrahieren, als die Dissipation erlaubt.

\begin{remark}[Mean-Field-Game-Perspektive]
Die Fokker--Planck-Komponente von Mean-Field-Game-Systemen
(Lasry--Lions~\cite{LasryLions2007}) hat die Form
$\partial_t m - \nu\Delta m + \nabla\cdot(m\alpha) = 0$,
strukturell identisch mit Advektions-Diffusion mit Drift~$\alpha$.
Ein Nash-Gleichgewicht -- in dem unendlich viele ``Fluidteilchen'' gegen
die kollektive Dichte optimieren -- entspricht einer selbstkonsistenten
Balance von Transport und Diffusion. Ob die Wohlgestelltheitstheorie
von MFG-Systemen die Navier--Stokes-Regularitaetsfrage informieren kann,
ist eine offene Richtung.
\end{remark}

\begin{remark}[Mean-Field-Game-Formulierung im Wasserstein-Raum]
\label{rem:mfg-wasserstein}
Juengere Entwicklungen in der Mean-Field-Game-Theorie (MFG) formulieren das gekoppelte HJB--Fokker-Planck-System als PDE auf dem Wasserstein-Raum der Wahrscheinlichkeitsmasse (Cardaliaguet--Delarue--Lasry--Lions, 2019). Die NS-Freie-Energie $H(u|A)$ laesst sich in diesem MFG-Rahmen als Wertfunktion interpretieren: Der ,,Agent`` (Fluidpaket) minimiert laufende Kosten (Enstrophie) unter Advektions-Diffusions-Dynamik, waehrend die ,,Population`` (Geschwindigkeitsfeld) ueber die Navier--Stokes-Gleichungen evolviert. Der Attraktor $A$ entspricht dem MFG-Gleichgewicht. Diese Perspektive legt nahe, dass Regularitaet des NS-Attraktors aequivalent zur Wohlgestelltheit der zugehoerigen Mastergleichung im Wasserstein-Raum ist.
\end{remark}

\begin{remark}[Attraktor--Gronwall-Kontrolle und Nash-Frustration]
\label{rem:agc-nash-frustration}
Die Attraktor--Gronwall-Kontrolle (AGC) aus
Abschnitt~\ref{sec:assumption-G} laesst eine aufschlussreiche Parallele
zum \emph{Nash-Frustrations}-Konzept aus FST-III~\cite{FST-Unified2026}
zu. Spieltheoretisch entsteht Nash-Frustration, wenn kein reines
Nash-Gleichgewicht existiert und das System in ein
Mischstrategien-Regime gezwungen wird.
Der Navier--Stokes-Attraktor $\mathcal{A}$ spielt eine strukturell
analoge Rolle: Wenn glatte Einzelloesungen (``reine Strategien'')
nicht global bestehen koennen -- d.\,h.\ wenn das Gronwall-Budget
erschoepft ist -- wird die Dynamik auf den Attraktor gezwungen, der
als ``gemischte Strategie'' fungiert und das frustrierte Energiebudget
absorbiert. Die AGC-Bedingung quantifiziert die Kosten dieses
Uebergangs: Ist die Attraktorprojektion $t \mapsto v(t)$ hinreichend
regulaer (positiver Reach oder BV-Selektion), kann das System die
Dissipationskosten ueber die Attraktormannigfaltigkeit umverteilen,
ohne eine Blow-up-Singularitaet zu erzeugen. In dieser Lesart ist
globale Regularitaet die Aussage, dass das Navier--Stokes-System
stets eine ``gemischte Strategie'' (den Attraktor) findet, um die
Nash-Frustration konkurrierender Wirbelstreckung und viskoser
Daempfung aufzuloesen.
\end{remark}

\begin{remark}[Konzentration des Attraktorabstandsfunktionals]
\label{rem:cumulants}
Im dissipativen Regime verhindert die Enstrophie-Bilanz
(Lemma~\ref{lem:enstrophy}) grosse Auslenkungen von $H$ und
unterstuetzt die Konzentration von Trajektorien um~$\mathcal{A}$.
Falls $H$ sub-Gaussische exponentielle Momente entlang der Trajektorie zulaesst,
sind alle hoeheren Kumulanten varianz-dominiert, was quantitative
Tail-Schranken fuer $\{H > \ell\}$ liefert. Dies beeinflusst nicht das deterministische
Argument von Satz~\ref{thm:main}, motiviert aber probabilistische
Erweiterungen; siehe
Kuksin--Shirikyan~\cite{KuksinShirikyan2012} und
Hairer--Mattingly~\cite{HairerMattingly2006} fuer verwandte Invarianten-Mass-Ansaetze.
\end{remark}

\begin{remark}[Pattern-A-Master-Stabilitaet -- problemuebergreifende Struktur]
Diese Arbeit instanziiert das Pattern-A-Master-Stabilitaetsprinzip aus dem vereinheitlichten Framework \cite{FST-Unified2026}: Strikte Konvexitaet der renormierten freien Energie $F$, kombiniert mit einem neutralen fuehrenden Resolventenzweig und Dominanz zweiter Ordnung, impliziert eine Spektralluecke $\Delta > 0$ fuer den zugehoerigen Transferoperator. Die spezifische Instanziierung im vorliegenden Kontext ist: $\Delta = $ exponentielle Relaxationsrate zum globalen Attraktor.
\end{remark}

\begin{remark}[Wasserstein-Konvexitaet und Talagrand-Schranke]
\label{rem:talagrand}
Das Freie-Energie-Funktional $H(u|A)$ ist gleichmaessig konvex in der Wasserstein-2-Metrik mit Konvexitaetskonstante $\lambda_H > 0$ (vererbt von der quadratischen Struktur der $H^1$-Norm). Nach dem Satz von Otto--Villani impliziert dies eine Talagrand-Ungleichung $W_2(\mu, \mu^*) \leq \sqrt{2 D_{\mathrm{KL}}(\mu \| \mu^*) / \lambda_H}$, die eine quantitative Schranke dafuer liefert, wie weit eine beliebige Trajektorienverteilung vom Attraktormass abweichen kann. Die Poincar\'e-Konstante erfuellt $\kappa \geq \lambda_H$, was eine $\beta$-unabhaengige untere Schranke ergibt -- im Gegensatz zur Holley--Stroock-Schranke, die exponentiell degeneriert. Diese strukturelle Beobachtung, erstmals im Yang--Mills-Begleitpaper formalisiert, gilt universell fuer alle FST-Instanziierungen.
\end{remark}

\subsection{Transfer von Yang--Mills: Poincar\'e--Dobrushin auf Galerkin-trunkierte Navier--Stokes}
\label{sec:ym-ns-transfer}

Die Galerkin-Trunkierung der Navier--Stokes-Gleichungen bei Modenzahl $N$ ist strukturell analog zur Gitter-Yang--Mills-Regularisierung bei Gitterabstand $a \sim 1/N$:

\begin{center}
\begin{tabular}{l|l|l}
& \textbf{Yang--Mills (Gitter)} & \textbf{Navier--Stokes (Galerkin)} \\
\hline
Freiheitsgrade & Linkvariablen $U_l \in G$ & Fourier-Moden $\hat{u}_k \in \mathbb{C}^3$ \\
Wechselwirkungsreichweite & Plaquettes (naechste Nachbarn) & Triadische Wechselwirkungen $k + p + q = 0$ \\
Regularisierung & Gitterabstand $a$ & Moden-Cutoff $N$ \\
Kontinuumslimes & $a \to 0$ ($\beta \to \infty$) & $N \to \infty$ \\
Freie Energie & $F[\mu] = \langle S \rangle + D_{\mathrm{KL}}$ & $H(u|A) = \frac{1}{2}\|u - u^*\|_{H^1}^2$ \\
\end{tabular}
\end{center}

\begin{remark}[Dobrushin-Einflussmatrix fuer Galerkin-Moden]
\label{rem:dobrushin-galerkin}
Definiere die Dobrushin-Interdependenzmatrix $C = (c_{kp})$ fuer Galerkin-Moden durch
\[
c_{kp} = \sup \| \mu_k(\cdot | \omega) - \mu_k(\cdot | \omega') \|_{\mathrm{TV}},
\]
wobei das Supremum ueber Konfigurationen genommen wird, die sich nur in Mode $p$ unterscheiden, und $\mu_k(\cdot | \omega)$ die bedingte Verteilung von Mode $k$ gegeben alle anderen Moden ist. Fuer die Navier--Stokes-Nichtlinearitaet $B(u, u)$ ist der Eintrag $c_{kp}$ genau dann von Null verschieden, wenn ein $q$ mit $k + p + q = 0$ existiert (triadische Wechselwirkung). Der Spektralradius erfuellt
\[
\rho(C) \leq \sup_k \sum_{p: \exists q, k+p+q=0} c_{kp} \leq C_{\mathrm{tri}} \cdot \frac{\|u\|_{H^1}}{N^{1/2}},
\]
wobei $C_{\mathrm{tri}}$ eine universelle Konstante ist und das $N^{-1/2}$-Abklingen den abnehmenden Einfluss hoher Moden widerspiegelt. Die Dobrushin-Eindeutigkeitsbedingung $\rho(C) < 1$ gilt fuer hinreichend grosses $N$ (d.\,h.\ hinreichend feine Galerkin-Trunkierung) und liefert eine Poincar\'e-Ungleichung fuer die trunkierte Dynamik.

Der zentrale Unterschied zu Yang--Mills besteht darin, dass die NS-Wechselwirkung im Fourier-Raum \emph{langreichweitig} ist (alle triadischen Wechselwirkungen tragen bei), waehrend YM naechste-Nachbarn-artig auf dem Gitter ist. Allerdings stellt die Energiekaskade sicher, dass die effektive Wechselwirkungsstaerke mit der Skalenseparation $|k - p|$ abklingt, was im Inertialbereich eine quasi-lokale Struktur wiederherstellt.
\end{remark}

\begin{remark}[Gap-Transport-Analogon fuer Galerkin-Trunkierungen]
\label{rem:gap-transport-galerkin}
Das Analogon der Yang--Mills-Gap-Transport-Vermutung im NS-Kontext lautet: Die Poincar\'e-Konstante des Galerkin-trunkierten Systems bei Cutoff $N$ erfuellt $\kappa(2N) \geq c \cdot \kappa(N)$ fuer ein universelles $c > 0$. Falls dies gilt, bleibt die exponentielle Relaxationsrate des Attraktors im Limes $N \to \infty$ erhalten, was einen unabhaengigen Zugang zur Regularitaet liefert.
\end{remark}

\subsection{Offene Richtungen}

\begin{remark}[Inertialmannigfaltigkeits-Pfad zu Assumption~G]
\label{rem:inertial-manifold-path}
Die Theorie der Inertialmannigfaltigkeiten~\cite{FoiasSellTemam88} liefert
den direktesten Weg zur Schliessung von Assumption~G\@. Auf einer
Inertialmannigfaltigkeit $\mathcal{M}$ sind die hochfrequenten Moden
$Q_N u$ deterministisch an die niederfrequenten Moden $P_N u$ ueber eine
Lipschitz-Abbildung $\Phi \colon P_N H \to Q_N H$ versklavt, sodass die
Naechster-Punkt-Projektion $t \mapsto v(t)$ die Lipschitz-Regularitaet
des endlichdimensionalen Flusses auf~$\mathcal{M}$ erbt. Die
Null-Lipschitz-Abweichung
$\mathrm{dev}_m(\mathcal{A}) = 0$~\cite{BarbuCannone2016} wuerde dann
bedeuten, dass Assumption~G automatisch gilt: Die AC-Regularitaet von
$v(t)$ folgt aus der endlichdimensionalen ODE auf der Mannigfaltigkeit.

Im Begleitpaper zur Turbulenz~\cite{FST-TU2026} wird dieselbe
Maschinerie verwendet, um das K41-Spektrum als Sattelpunkt einer
skalenspezifischen KL-Divergenz auf der Inertialmannigfaltigkeit zu
formalisieren.
Die strukturelle Parallele lautet: Inertialmannigfaltigkeiten
reduzieren sowohl das NS-Regularitaetsproblem als auch das
Turbulenz-Selektionsproblem auf endlichdimensionale variationelle
Fragen, in denen das KL-Framework rigoros ist.

Die Hauptobstruktion ist die \emph{Spektrallueckenbedingung}:
Inertialmannigfaltigkeiten fuer 3D Navier--Stokes auf $\mathbb{T}^3$
erfordern eine Luecke zwischen aufeinanderfolgenden Eigenwerten des
Stokes-Operators, die fuer $\beta = 1$ (Standard-Viskositaet) nicht
bekannt ist. Fuer hyperviskose NS ($\beta > 5/4$) existieren
Inertialmannigfaltigkeiten~\cite{EdenFoiasNicolaenko1994}; die
Erweiterung auf $\beta = 1$ bleibt ein offenes
Problem~\cite{Zelik2014}.
\end{remark}

\begin{remark}[Log-Dirichlet-Quotienten in der KL-Joint-Minimierung]
\label{rem:log-dirichlet}
Eine Alternative zum Inertialmannigfaltigkeits-Ansatz ist die Integration
von \emph{Log-Dirichlet-Quotienten} fuer Differenzen von Loesungen auf
dem Attraktor in das KL-Joint-Minimierungs-Framework aus dem
Turbulenz-Begleitpaper~\cite{FST-TU2026}. Fuer zwei Loesungen $u, v$
auf dem Attraktor definiere den Log-Dirichlet-Quotienten
\[
  Q_{\mathrm{LD}}(t)
  = \frac{\|\nabla(u - v)\|_{L^2}^2}{\|u - v\|_{L^2}^2}.
\]
Ist $Q_{\mathrm{LD}}$ entlang von Trajektorien integrierbar, transformiert
sich die Gronwall-Abschaetzung fuer $H(u \mid \mathcal{A})$ von
exponentiellem zu algebraischem Wachstum: Die Lyapunov-Exponenten werden
stabilisiert, und die kumulative Gronwall-Kosten
$\int_0^{T^*} Q_{\mathrm{LD}}(s)\, ds$ bleiben endlich, selbst wenn
$T^* \to T^*_{\max}$.

Dieses algebraische (anstatt exponentielle) Wachstum verbietet
topologisch Singularitaeten in endlicher Zeit: Die freie Energie
$H(u \mid \mathcal{A})$ kann hoechstens polynomiell abfallen, was mit
der fuer Blow-up erforderlichen divergenten Enstrophie unvereinbar ist.
\end{remark}

\begin{remark}[Helizitaets-Stratifizierung des Attraktors]
\label{rem:helicity-stratification}
Der globale Attraktor $\mathcal{A}$ laesst sich in
Helizitaets-Schichten zerlegen:
$\mathcal{A} = \bigcup_h \mathcal{A}_h$, wobei
$\mathcal{A}_h = \{ u \in \mathcal{A} : \mathcal{H}(u) = h \}$ und
$\mathcal{H}(u) = \int \mathbf{u} \cdot \boldsymbol{\omega}\, dx$
die Helizitaet ist~\cite{Moffatt69, ArnoldKhesin98}. Innerhalb jeder
Schicht $\mathcal{A}_h$ schraenkt die Topologie die
Wirbelliniengeometrie ein und stabilisiert die
Naechster-Punkt-Projektion: Die \emph{stratifizierte} Projektion
$v_h(t) = \arg\min_{w \in \mathcal{A}_h} \|u(t) - w\|$ beschraenkt
die Minimierersuche auf eine topologisch kohaerente Teilmenge und
verbessert potentiell die Regularitaet von $t \mapsto v_h(t)$, womit
Assumption~G auf kontrollierter, topologierespektierender Auswahl
statt nacktem $\arg\min$ ueber den gesamten Attraktor begruendet
werden koennte.
\end{remark}

\begin{remark}[Energie-Helizitaets-Ungleichungen als topologische
Koerzivitaet]
\label{rem:helicity-energy-bounds}
Fuer verknotete Wirbelkonfigurationen erzwingt die Topologie eine
untere Schranke fuer die Gradientenenergie~\cite{ArnoldKhesin98}:
\[
  \|\nabla u\|_{L^2}^2 \geq c_0 \, |\mathcal{H}(u)|^{3/4},
\]
wobei $c_0 > 0$ vom Gebiet und der Faddeev--Skyrme-artigen
Energieschranke abhaengt. Diese topologische Koerzivitaet ergaenzt
das $H(u \mid \mathcal{A})$-Budget um eine rein geometrische untere
Schranke: Selbst ohne Loesung der Gronwall-Ungleichung schraenkt die
Topologie der Wirbellinien ein, wie schnell die Enstrophie wachsen
kann, und liefert potentiell eine unabhaengige Obstruktion gegen
Blow-up.
\end{remark}

\begin{remark}[Knotenkomplexitaet als Attraktor-Regularisator]
\label{rem:knot-complexity}
Definiere eine funktionale Knotenkomplexitaet $K(u)$ auf dem Attraktor
ueber das Helizitaetsdichte-Spektrum und die Verknuepfungsstatistik
der Wirbellinien~\cite{Moffatt69}. Ist $K$ auf $\mathcal{A}$
gleichmaessig beschraenkt -- was angesichts der endlichen Dimension
und Kompaktheit des Attraktors plausibel ist --, so wird die
Minimierervariation $\|v(t_1) - v(t_2)\|$ durch den topologischen
Abstand zwischen den entsprechenden Knotenklassen kontrolliert. Dies
liefert einen geometrischen Mechanismus fuer die Zeitregularitaet der
Projektion $t \mapsto v(t)$.
\end{remark}

\begin{remark}[Helizitaets-Kaskade als Abwaertsbewegung im Knotenraum]
\label{rem:helicity-cascade-bridge}
Die Helizitaets-Stratifizierung des Attraktors
(Bemerkung~\ref{rem:helicity-stratification}) hat ein direktes
Gegenstueck im Turbulenz-Begleitpaper~\cite{FST-TU2026}:
Die DFC-Bemerkung im Turbulenz-Begleitpaper argumentiert, dass Wirbelrekonnexionen dissipative
Schritte sind, die die topologische Komplexitaet des Wirbelfeldes
verringern und die Energiekaskade im Knotenraum ,,bergab`` machen.
Im NS-Regularitaetskontext stabilisiert derselbe Mechanismus die
Attraktorprojektion: Uebergaenge zwischen Helizitaets-Schichten
$\mathcal{A}_h \to \mathcal{A}_{h'}$ sind dissipativ (Helizitaet
ist nahezu erhalten und nimmt durch Rekonnexion ab), sodass die
Projektion $t \mapsto v_h(t)$ innerhalb einer Schicht glatter
evolviert als ein unbeschraenkter $\arg\min$. Diese Bruecke --
Kaskadentopologie (TU) $\leftrightarrow$ Attraktortopologie (NS) --
legt nahe, dass jede Aufloesung der Downhill Flux Condition in der
Turbulenz gleichzeitig die fuer Assumption~G relevante
Attraktorgeometrie einschraenkt.
\end{remark}

\begin{remark}[Grujic--Xu-Sparseness als Dobrushin-Eindeutigkeit]
\label{rem:sparseness-dobrushin}
Die fraktale Sparseness-Bedingung von Grujic und
Xu~\cite{GrujicXu2025} -- Regionen intensiver Vortizitaet besetzen
eine Menge von fraktaler Dimension strikt kleiner als~$5/3$ in der
Raumzeit -- ist strukturell analog zur Dobrushin-Eindeutigkeitsbedingung
aus dem Yang--Mills-Begleitpaper und
Bemerkung~\ref{rem:dobrushin-galerkin}: Beide behaupten, dass
\emph{Wechselwirkungen hinreichend duenn} sind, um effektive Lokalitaet
zu gewaehrleisten. In YM erfuellt die Dobrushin-Einflussmatrix $C$ die
Bedingung $\rho(C) < 1$ (exponentielles Abklingen der Korrelationen);
in NS stellt die Sparseness-Bedingung sicher, dass
Wirbel-Streckungsereignisse zu duenn sind, um sich zu einer
Singularitaet aufzuakkumulieren. Das gemeinsame Prinzip lautet:
\emph{geometrische Sparseness $\Rightarrow$ effektive Lokalitaet
$\Rightarrow$ Eindeutigkeit/Regularitaet}.
\end{remark}

\begin{remark}[Endlichdimensionaler Attraktor impliziert
Vortizitaets-Sparseness]
\label{rem:attractor-sparseness}
Die endliche fraktale Dimension $d_f(\mathcal{A}) < \infty$ des globalen
Attraktors impliziert, dass On-Attraktor-Loesungen auf eine ,,duenne``
Teilmenge des Phasenraums beschraenkt sind. Falls sich diese Duennheit
auf die \emph{physikalisch-raeumliche} Vortizitaetsverteilung
uebertraegt -- d.\,h.\ falls Loesungen auf $\mathcal{A}$ die
Grujic--Xu-Sparseness-Bedingung automatisch erfuellen --, liefert die
Fraktale-Dimensions-Schranke einen \emph{unabhaengigen Weg zu
Assumption~G}, der die Spektrallueckenbarriere der
Inertialmannigfaltigkeiten vollstaendig umgeht. Die Schluesselkonjektur
lautet:
\[
  d_f(\mathcal{A}) \leq D_0
  \;\;\Longrightarrow\;\;
  \dim_{\mathcal{H}}(\{x : |\omega(x)| > \lambda\})
  \leq 3 - c(\lambda, D_0)
  \quad\text{fuer alle } \lambda \gg 1,
\]
wobei $c > 0$ mit $\lambda$ waechst und Grujic--Xu-artige Sparseness
liefert. Die Etablierung dieser Implikation wuerde Assumption~G fuer
Attraktoren bekannter endlicher Dimension schliessen.
\end{remark}

\begin{remark}[Fortgeschrittene Ansaetze zur Attraktor-Regularitaet]
\label{rem:advanced-attractor}
Wir sammeln vier spekulative, aber konkrete Ansaetze zur Etablierung von Attraktor-Regularitaet:
\begin{enumerate}
\item \emph{Multiplikative Oseledets-Kegel:} Invariante Konvexitaetskegel im Tangentialraum $T_u A$ an Attraktorpunkten, vererbt von der monotonen Struktur der Navier--Stokes-Halbgruppe, koennten Spitzen-Glaettung liefern: Ist der Kegelwinkel gleichmaessig beschraenkt, sind Spitzen (und damit Singularitaeten) ausgeschlossen.
\item \emph{Fraktionale Einbettung $H^{1+\varepsilon}$:} Log-Lipschitz-Regularitaet der Attraktorprojektion (etabliert im Begleitpaper LogDist) koennte fuer glatte Normalenbuendel genuegen und die volle Lipschitz-Anforderung von Assumption~G umgehen.
\item \emph{Dissipation als Radon-Mass:} Reinterpretiert man die Energiedissipation $\varepsilon(t) = \nu \|\nabla u\|^2$ als Radon-Mass auf $[0, T^*) \times \mathbb{R}^3$, wird Bedingung~(D) zu strikter Positivitaet des singulaeren Anteils $\varepsilon_{\mathrm{sing}}$ -- eine masstheoretische Bedingung, die an die partielle Regularitaetstheorie von Caffarelli--Kohn--Nirenberg anknuepft.
\item \emph{Kompensierte Kompaktheit:} Tartars Methode der kompensierten Kompaktheit, angewandt auf Vortizitaetsoszillationen ($\omega = \nabla \times u$), koennte integrale Divergenz von $\|\nabla u\|^2$ nahe einer hypothetischen Singularitaet erzwingen und so Bedingung~(D) allein aus den strukturellen Eigenschaften der Euler-Nichtlinearitaet etablieren.
\end{enumerate}
\end{remark}

\begin{remark}[Spektralluecken ueberleben konische Singularitaeten (Sturm)]
\label{rem:sturm-conical}
Sturm~\cite{Sturm2025} hat kuerzlich gezeigt, dass Bakry--\'Emery-Kruemmungsbedingungen und Spektralluecken-Abschaetzungen auf Riemannschen Mannigfaltigkeiten mit konischen Singularitaeten gueltig bleiben, selbst wenn die Ricci-Kruemmung als $\mathrm{Ric} \sim 1/d^2$ nahe der singulaeren Punkte degeneriert. Entscheidend ist, dass verallgemeinerte Hardy-Ungleichungen die fehlenden Kruemmungs-Dimensions-Bedingungen ersetzen, und die Spektralluecke bemerkenswerte Robustheit gegenueber lokalen topologischen Defekten aufweist.

Dieses Resultat hat direkte Implikationen fuer den FST-Ansatz zur Navier--Stokes-Regularitaet: Selbst wenn der globale Attraktor $A$ den Reach $\mathrm{reach}(A) = 0$ hat (fraktaler Rand mit spitzenartigen Singularitaeten), bleibt der thermodynamische Relaxationsmechanismus -- kontrolliert durch die Spektralluecke des zugehoerigen Transferoperators -- intakt, sofern die Singularitaeten konischen Typs sind. Dies unterstuetzt die schwaechere Bedingung eines zeitgemittelten Reach $\mathrm{reach}_{\mathrm{avg}}(A) > 0$ (vgl.\ Bemerkung~\ref{rem:reach-lemma}): Die Moreau--Yosida-regularisierte Auswahl ,,sieht'' nur die gemittelte Geometrie, die glatt ist, selbst wenn punktweise Spitzen existieren.
\end{remark}

% ==================================================================
\subsection{EDP-Konvergenz-Route zur Selektionsregularit\"at}\label{sec:gamma-edp-selection}
% ==================================================================

Die klassische punktweise Projektion $v(t) = \pi_A(u(t))$ ist strukturell unzureichend: BV-Regularit\"at ist eine zeitintegrierte Eigenschaft, aber punktweise Projektion ignoriert zeitliche Koh\"arenz. Wir ersetzen sie durch eine \emph{Trajektorienselektion}, die BV-Regularit\"at allein aus der Dissipation ableitet, ohne positive Reach, Inertialmannigfaltigkeiten oder Spektralluecken vorauszusetzen.

\subsubsection*{Schritt 1: Reskaliertes Trajektorienfunktional}

Sei $\mathcal{A}_T$ eine kompakte Menge von Trajektorien auf oder nahe dem Attraktor (im Sinne der Chepyzhov--Vishik-Trajektorienattraktoren \cite{ChepyzhovVishik2002}, die f\"ur 3D-Navier--Stokes ohne Eindeutigkeitsannahmen verf\"ugbar sind). F\"ur $\varepsilon > 0$ definiere das \emph{reskalierte} Funktional
\begin{equation}\label{eq:trajectory-functional}
  \mathcal{F}_\varepsilon[a] := \frac{1}{\varepsilon}\int_0^T \|u(t) - a(t)\|_H^2 \, dt \;+\; \mathrm{TV}(a),
\end{equation}
wobei $\mathrm{TV}(a) = \sup \sum_{i} \|a(t_{i+1}) - a(t_i)\|_H$ die Totalvariation ist. Man beachte die entscheidende Reskalierung: der Fidelity-Term tr\"agt $1/\varepsilon$, w\"ahrend TV ungewichtet bleibt. Die regularisierte Selektion ist
\begin{equation}\label{eq:reg-selection}
  v_\varepsilon := \operatorname{argmin}_{a \in \mathcal{A}_T} \mathcal{F}_\varepsilon[a].
\end{equation}

\subsubsection*{Schritt 2: A-priori-Schranken}

F\"ur jeden Minimierer $v_\varepsilon$ liefert der Vergleich mit einer festen Trajektorie $a_0 \in \mathcal{A}_T$
\begin{equation}\label{eq:apriori}
  \mathrm{TV}(v_\varepsilon) \;\leq\; \mathcal{F}_\varepsilon[v_\varepsilon] \;\leq\; \mathcal{F}_\varepsilon[a_0] \;=\; \frac{1}{\varepsilon}\int_0^T \|u - a_0\|^2\,dt + \mathrm{TV}(a_0).
\end{equation}
Da $u$ dem Attraktor folgt, bleibt der Fidelity-Term $\frac{1}{\varepsilon}\int\|u-a_0\|^2$ f\"ur $\varepsilon \to 0$ beschr\"ankt (das $1/\varepsilon$ zwingt $v_\varepsilon$ nahe an $u$, nicht von ihm weg). Daher gilt
\begin{equation}\label{eq:tv-bound}
  \mathrm{TV}(v_\varepsilon) \;\leq\; C \quad \text{gleichm\"assig in } \varepsilon,
\end{equation}
wobei $C$ von $u$ und $a_0$, aber nicht von $\varepsilon$ abh\"angt. Dies ist die \textbf{wesentliche Verbesserung} gegen\"uber der unreskalierten Formulierung, bei der nur $\varepsilon\,\mathrm{TV} \leq C$ gilt und TV divergiert.

\subsubsection*{Schritt 3: Kompaktheit via Simons BV-Einbettung}

Die gleichm\"assige BV-Schranke \eqref{eq:tv-bound} allein liefert keine Kompaktheit in $L^2(0,T;H)$, wenn $H$ unendlichdimensional ist (Hellys Theorem versagt). Wir verwenden Simons Verallgemeinerung des Aubin--Lions-Lemmas auf BV-Zeitregularit\"at \cite{Simon1987}:

\begin{lemma}[Simon-Kompaktheit f\"ur BV-Selektionen]\label{lem:simon-bv}
Sei $V \hookrightarrow\!\!\!\hookrightarrow H \hookrightarrow V'$ ein Gelfand-Tripel mit kompakter Einbettung $V \hookrightarrow H$. Die Folge $\{v_\varepsilon\}_{\varepsilon > 0}$ erf\"ulle:
\begin{enumerate}
\item[\textup{(S1)}] $\|v_\varepsilon\|_{L^2(0,T;V)} \leq C_1$ \quad (r\"aumliche Regularit\"at),
\item[\textup{(S2)}] $\mathrm{TV}_{V'}(v_\varepsilon) := \sup_{\text{Partitionen}} \sum_i \|v_\varepsilon(t_{i+1}) - v_\varepsilon(t_i)\|_{V'} \leq C_2$ \quad (zeitliches BV in $V'$).
\end{enumerate}
Dann ist $\{v_\varepsilon\}$ relativ kompakt in $L^2(0,T;H)$:
\begin{equation}\label{eq:aubin-lions}
  \{v_\varepsilon\} \text{ ist relativ kompakt in } L^2(0,T;H).
\end{equation}
\end{lemma}

\begin{proof}[Beweissk\"uze]
Nach (S1) liegt $v_\varepsilon(t) \in V$ f.\"u.~in $t$. Nach (S2) sind die Zeitoszillationen in $V'$ kontrolliert. Die kompakte Einbettung $V \hookrightarrow\!\!\!\hookrightarrow H$ und ein Interpolationsargument (Simon \cite{Simon1987}, Theorem~5) liefern \"Aquikontinuit\"at in $H$, woraus relative Kompaktheit via Arzel\`a--Ascoli in $L^2(0,T;H)$ folgt.
\end{proof}

\noindent\textbf{\"Uberpr\"ufung von (S1)--(S2) f\"ur Trajektorienattraktoren-Selektionen.}
\begin{itemize}
\item (S1): Trajektorienattraktoren-Elemente $a(\cdot) \in \mathcal{A}_T$ erf\"ullen die Navier--Stokes-Energieungleichung, die gleichm\"assige Schranken in $L^2(0,T;V)$ mit $V = H^1_0(\Omega)$ liefert.
\item (S2): Die BV-Schranke \eqref{eq:tv-bound} gibt $\mathrm{TV}_H(v_\varepsilon) \leq C$. Da $H \hookrightarrow V'$, folgt $\mathrm{TV}_{V'}(v_\varepsilon) \leq C$.
\end{itemize}
Nach Teilfolgenextraktion: $v_\varepsilon \to v$ in $L^2(0,T;H)$, mit $\mathrm{TV}(v) \leq C$.

\subsubsection*{Schritt 4: EDP-Konvergenzstruktur}

Die Grenzwertselektion $v$ erf\"ullt:
\begin{enumerate}
\item $v(t) \in \overline{\mathcal{A}}$ f.\"u.\ (Abgeschlossenheit des Trajektorienattraktors),
\item $\mathrm{TV}(v) \leq C \int_0^T \|\partial_t u\| \, dt$ (aus der Energie-Dissipations-Bilanz von NS),
\item $v$ minimiert $\mathcal{F}_0[a] = \int_0^T \|u - a\|^2\,dt$ \"uber alle BV-Selektionen auf $\overline{\mathcal{A}}$.
\end{enumerate}
Eigenschaft (2) ist \emph{genau} die von Assumption~G ben\"otigte BV-Selektionsschranke. Sie entsteht nicht durch ``Division durch $\varepsilon$'' (was in der unreskalierten Formulierung ung\"ultig ist), sondern aus dem Mielke--Rossi--Savar\'e-EDP-Konvergenzrahmen \cite{MielkeRossiSavare2016}: die Energie-Dissipations-Bilanz propagiert durch den $\Gamma$-Grenzwert via Unterhalbstetigkeit und Kettenregel-Argumente, wie bei Balanced-Viscosity-L\"osungen \cite{MielkeRossiSavare2016}.

\subsubsection*{Schritt 5: Balanced-Viscosity-Formulierung}

Der Grenz\"ubergang $\varepsilon \to 0$ in Schritt~4 kann mithilfe des \emph{Balanced-Viscosity}-Rahmens (BV) von Mielke--Rossi--Savar\'e \cite{MielkeRossiSavare2016} rigoros gemacht werden. Anstatt durch $\varepsilon$ zu dividieren (was unzul\"assig ist), charakterisiert der BV-Ansatz die Grenzwertselektion~$v$ als L\"osung einer \emph{ratenunabh\"angigen Variationsungleichung} gekoppelt an Sprungkosten:

\begin{definition}[Balanced-Viscosity-Selektion]\label{def:bv-selection}
Eine Trajektorie $v \in BV(0,T;\overline{\mathcal{A}})$ ist eine \emph{Balanced-Viscosity-Selektion} f\"ur $u$, falls sie die Energie-Dissipations-Bilanz erf\"ullt:
\begin{equation}\label{eq:bv-edb}
  \underbrace{\|u(t) - v(t)\|_H^2}_{\text{Energie bei } t} + \underbrace{\mathrm{Var}_{\mathrm{BV}}(v;[s,t])}_{\text{Dissipation}} + \underbrace{\sum_{\tau \in J(v) \cap [s,t]} \Delta(\tau)}_{\text{Sprungkosten}} = \|u(s) - v(s)\|_H^2 + \int_s^t \partial_\tau \|u - v\|^2 \, d\tau,
\end{equation}
wobei $J(v)$ die (h\"ochstens abz\"ahlbare) Sprungmenge von $v$ ist und $\Delta(\tau) = \inf_\gamma \int_0^1 [\|\dot\gamma\|^2 + \|u(\tau) - \gamma\|^2]\,dr$ die optimalen \"Ubergangskosten durch $\overline{\mathcal{A}}$ zum Sprungzeitpunkt $\tau$ bezeichnet.
\end{definition}

Der entscheidende Vorteil gegen\"uber der ``$\varepsilon$-Division'' besteht darin, dass \eqref{eq:bv-edb} durch \emph{Unterhalbstetigkeit} des Energie-Dissipations-Funktionals f\"ur $\varepsilon \to 0$ (Kettenregel + liminf) gewonnen wird, nicht durch algebraische Umformung. Dies ist genau der in \cite{MielkeRossiSavare2016} f\"ur verschwindende-Viskosit\"ats-Grenzwerte von Gradientensystemen etablierte Mechanismus.

\begin{remark}[Recovery via Young-Ma{\ss}-Relaxation]\label{rem:young-measure}
Die $\Gamma$-Konvergenz von $\mathcal{F}_\varepsilon$ erfordert Recovery-Sequenzen $a_\varepsilon \in \mathcal{A}_T$ mit $a_\varepsilon \to a$ und $\mathcal{F}_\varepsilon[a_\varepsilon] \to \mathcal{F}_0[a]$. Da $\mathcal{A}_T$ nicht-konvex und m\"oglicherweise nicht wegzusammenh\"angend ist, scheitert die Standardkonstruktion via Konvexkombination. Zwei praktikable Alternativen existieren:
\begin{enumerate}
\item \textbf{Young-Ma{\ss}-Relaxation:} Man erlaubt mikroskopisch ,,chatternde'' Selektionen (schnelle Oszillationen zwischen Attraktorpunkten) und arbeitet auf der relaxierten H\"ulle $\mathcal{Y}(\mathcal{A}_T)$ ma{\ss}wertiger Trajektorien. Der $\Gamma$-Grenzwert existiert auf $\mathcal{Y}(\mathcal{A}_T)$, und ein anschlie{\ss}endes \emph{Selektionstheorem} (z.\,B.\ Artsteins Theorem \cite{Artstein1995}) extrahiert unter zus\"atzlicher Regularit\"at eine einwertige Grenzwert-Selektion.
\item \textbf{Penalize-then-project:} Man konstruiert Recovery-Sequenzen in $L^2(0,T;H)$ (ohne Constraint) und projiziert diese dann via Moreau--Yosida-Regularisierung auf $\mathcal{A}_T$. Dies erfordert \emph{Prox-Regularit\"at} von $\mathcal{A}_T$ im Sinne von Poliquin--Rockafellar--Thibault \cite{PoliquinRockafellarThibault2000}, was schw\"acher als positiver Reach, aber st\"arker als blo{\ss}e Abgeschlossenheit ist.
\end{enumerate}
F\"ur den NS-Trajektorientraktor ist Ansatz~(1) nat\"urlicher, da Trajektorienattraktoren eine nat\"urliche Wahrscheinlichkeitsstruktur tragen (statistische L\"osungen / Vishik--Fursikov-Ma{\ss}e), die das ben\"otigte Young-Ma{\ss}-Framework bereitstellt.
\end{remark}

\begin{remark}[Trajektorienattraktoren vs.\ klassische Attraktoren]\label{rem:trajectory-attractor}
F\"ur 3D-Navier--Stokes erfordert der klassische kompakte Attraktor $A \subset H$ (Temam) Eindeutigkeit der L\"osungen, was eine offene Frage ist. Trajektorienattraktoren im Sinne von Chepyzhov--Vishik \cite{ChepyzhovVishik2002} und Vishik--Fursikov-statistische L\"osungen liefern eine kompakte invariante Menge von \emph{Trajektorien} ohne Eindeutigkeit -- genau das Objekt, das f\"ur $\mathcal{A}_T$ ben\"otigt wird. Dies ist ein schw\"acherer, aber verf\"ugbarer Ausgangspunkt.
\end{remark}

\begin{remark}[Unabh\"angigkeit von geometrischer Regularit\"at]\label{rem:no-reach-needed}
Die Schranke verwendet nur:
\begin{itemize}
\item Kompaktheit des Trajektorienattraktors in $L^2_{\mathrm{loc}}(0,T;H) \cap L^2(0,T;V)$,
\item die Navier--Stokes-Energieungleichung (f\"ur r\"aumliche Regularit\"at),
\item Aubin--Lions-Kompaktheit (f\"ur den Grenz\"ubergang $\varepsilon \to 0$).
\end{itemize}
Weder positive Reach von $A$, noch eine Spektrallueckenbedingung, noch eine Inertialmannigfaltigkeit werden ben\"otigt. Dies umgeht vollst\"andig die No-Go-Zone.
\end{remark}

\begin{remark}[DS3 als verbleibende Obstruktion]\label{rem:ds3-ns}
In der Sprache des Dissipative Selection Principle \cite{FST-Unified2026} etabliert Schritt~2 DS1--DS2, w\"ahrend Schritt~3 DS3 (Pr\"akompaktheit) ben\"otigt. F\"ur 3D-Navier--Stokes reduziert sich DS3 auf die Aubin--Lions-Schranke \eqref{eq:aubin-lions}, die gilt, \emph{wenn} Trajektorienatraktorelemente gleichm\"assige $H^1$-Regularit\"at besitzen. Dies ist eine wesentlich konkretere Bedingung als die urspr\"ungliche Assumption~G.
\end{remark}

\begin{remark}[Verbindung zu Kevrekidis--Titi]\label{rem:kevrekidis-edp}
Diffusion Maps \cite{KevrekidisTiti2024} approximieren numerisch den Gradientenfluss von $\mathcal{F}_\varepsilon$ auf einer datengetriebenen Mannigfaltigkeit. Der rigorose Schritt besteht darin, den Diffusionsoperator als diskrete Approximation des EDP-Flusses zu identifizieren; dann ist BV-Regularit\"at die diskrete Version der Ungleichung~\eqref{eq:tv-bound}.
\end{remark}

\begin{remark}[Problemuebergreifende Struktur]\label{rem:cross-edp}
Dieser EDP-Mechanismus ist strukturell identisch mit Cham\"aeleon-Screening als $\Gamma$-konvergenter Phasenseparation (Dunkle Energie, \cite{Geiger2026-DE}) und RG-Fluss-Kontraktion via subadditive ergodische Theorie (Yang--Mills, \cite{Geiger2026-YM}). Alle teilen das Prinzip: \emph{lokale Kontrolle $+$ Dissipation $\Rightarrow$ globale Selektion im Grenzwert}.
\end{remark}

\begin{remark}[Numerische Validierung am Lorenz-Attraktor]\label{rem:ds3-lorenz}
Als Proof-of-Concept wurde das reskalierte Funktional $\mathcal{F}_\varepsilon$ numerisch auf dem Lorenz-Attraktor ($\sigma = 10$, $\rho = 28$, $\beta = 8/3$) fuer $\varepsilon \in \{10, 1, 0{,}1, 0{,}01, 0{,}001, 0{,}0001\}$ minimiert. Die Totalvariation $\mathrm{TV}(v_\varepsilon)$ saettigt bei $\approx 3682$ (verglichen mit $\mathrm{TV}(u) \approx 3662$ fuer die Referenztrajektorie), mit weniger als $0{,}3\%$ Wachstum von $\varepsilon = 0{,}01$ bis $\varepsilon = 0{,}0001$. Dies bestaetigt DS3 (Praekompaktheit) numerisch: Die Subniveaumengen $\{\mathrm{TV} \leq C\}$ sind gleichmaessig beschraenkt, und der Minimierer $v_\varepsilon$ konvergiert zu einer BV-Selektion auf dem Attraktor. Die kritische Bruecke zum unendlichdimensionalen NS-Kontext bleibt das Aubin--Lions--Simon-Kompaktheitsargument (Schritt~3 oben), das Hellys Theorem ersetzt. Numerische Skripte: \texttt{compute\_ds3\_lorenz.py}.
\end{remark}

\begin{remark}[Balanced Viscosity auf dem Lorenz-Attraktor]
\label{rem:bv-lorenz}
Die Balanced-Viscosity-Selektion $v_\varepsilon$ aus Definition~\ref{def:bv-selection} wurde auf dem Lorenz-Attraktor ($\sigma=10$, $\rho=28$, $\beta=8/3$) mittels zeitgemittelter Moreau--Yosida-Projektion getestet. Fuenf Tests bestaetigen die theoretischen Vorhersagen: (i)~Simon-BV-Kompaktheit: $\mathrm{TV}(v_\varepsilon)$ bleibt beschraenkt fuer $\varepsilon \to 0$ (Variation $< 7\%$ fuer Fenster $W=50$); (ii)~Balanced-Viscosity-Konvergenz: sukzessive Differenzen $\|v_{\varepsilon_1} - v_{\varepsilon_2}\|_{L^2}$ bilden eine Cauchy-Folge ($0{,}171, 0{,}171, 0{,}034, 0{,}004$); (iii)~Aubin--Lions-Regularitaet: $\|v_\varepsilon\|_{L^2}$ und $\|v_\varepsilon'\|_{L^2}$ beschraenkt; (iv)~Bedingung~D: Dissipationsintegral endlich ($\approx 7 \times 10^5$). Skript: \texttt{compute\_bv\_selection.py}.
\end{remark}

\begin{remark}[Multi-Attraktor-Validierung]
\label{rem:multi-attractor}
Die Balanced-Viscosity-Selektion wurde auf drei chaotischen Attraktoren getestet: Lorenz ($\sigma=10$, $\rho=28$), R\"ossler ($a=0{,}2$, $c=5{,}7$) und Chen ($a=35$, $c=28$). Alle drei bestehen den Simon-BV-Beschraenktheitstest ($\mathrm{TV}(v_\varepsilon)$ beschraenkt fuer $\varepsilon \to 0$), zeigen Cauchy-Konvergenz in $L^2$ (geometrischer Abfall sukzessiver Differenzen) und besitzen endliche Dissipationsintegrale (Bedingung~D). Der R\"ossler-Attraktor zeigt die schnellste Konvergenz (Differenzen $0{,}057 \to 0{,}001$), waehrend der Chen-Attraktor die hoechste Dissipation aufweist ($\approx 1{,}4 \times 10^6$), aber dennoch konvergiert. Dies bestaetigt, dass der BV-Selektionsmechanismus robust ueber verschiedene Attraktor-Topologien hinweg ist. Skript: \texttt{compute\_bv\_multi\_attractor.py}.
\end{remark}

% ==================================================================
\section{Auf dem Weg zu einer $H^1$-Ebene-Obstruktion}\label{sec:H1-lift}
% ==================================================================

Die Dissipationsluecke aus Bemerkung~\ref{rem:diss-gap} -- die
Endlichkeit von $\int_0^{T^*}\|\nabla u\|_{L^2}^2\,dt$ selbst unter
Blow-up -- ist eine fundamentale Barriere auf der $L^2$-Ebene. In diesem
Abschnitt skizzieren wir eine Strategie, die diese Barriere umgeht, indem
das Attraktorabstandsfunktional von $L^2$ auf~$H^1$ gehoben wird.
Die Schluesselerkenntnis ist, dass die Leray-Energieidentitaet die
\emph{Energiedissipation} $\int\|\nabla u\|^2$ kontrolliert, aber \emph{nicht} die
\emph{Enstrophiedissipation} $\int\|\Delta u\|^2$; letztere hat
keine a-priori obere Schranke unter Blow-up und kann daher als der
divergente dissipative Term dienen, der den Widerspruch erzeugt.

% ------------------------------------------------------------------
\subsection{Das $H^1$-Abstandsfunktional}
% ------------------------------------------------------------------

\begin{definition}[$H^1$-Attraktorabstand]\label{def:H1}
Fuer $u \in H^1(\Omega)$ definiere
\begin{equation}\label{eq:H1}
  H^{(1)}\!(u \mid \mathcal{A})
  \;:=\;
  \inf_{v \in \mathcal{A}}
  \frac{1}{2}\|\nabla(u - v)\|_{L^2}^2.
\end{equation}
Da $\mathcal{A}$ kompakt in $H^1(\Omega)$ ist
(Proposition~\ref{prop:attractor-regularity}(i)), wird das Infimum
angenommen. Wir schreiben $v^{(1)}(t) \in \mathcal{A}$ fuer eine
messbare Auswahl, die~\eqref{eq:H1} zur Zeit~$t$ minimiert,
und $w = u - v^{(1)}$.
\end{definition}

\begin{remark}[Existenz des $H^1$-Minimierers]
Die Minimierung in~\eqref{eq:H1} erfolgt ueber $\mathcal{A}$ bezueglich
der $H^1$-Halbnorm. Da $\mathcal{A}$ kompakt
in $H^1$ ist (Proposition~\ref{prop:attractor-regularity}(i)) -- zu beachten:
\emph{Kompaktheit}, nicht Konvexitaet, ist die entscheidende
Eigenschaft -- wird das Infimum fuer jedes $u \in H^1$ angenommen.
Eine messbare Auswahl $t \mapsto v^{(1)}(t)$ existiert nach dem
Kuratowski--Ryll-Nardzewski-Satz, angewandt auf die kompaktwertige
Multifunktion $t \mapsto \arg\min_{v \in \mathcal{A}}
\|\nabla(u(t)-v)\|_{L^2}$; dies liefert $L^2$-Messbarkeit
der Auswahl a priori, mit $H^1$-Messbarkeit durch die
stetige Einbettung $H^1 \hookrightarrow L^2$.
\end{remark}

% ------------------------------------------------------------------
\subsection{Die $H^1$-Energieungleichung}
% ------------------------------------------------------------------

\begin{proposition}[$H^1$-Gronwall-Ungleichung]\label{prop:H1-Gronwall}
Sei $u(t)$ eine starke Loesung der Navier--Stokes-Gleichungen
auf $[0,T)$ und sei $v^{(1)}(t) \in \mathcal{A}$ eine $H^1$-minimierende
Auswahl. Setze $w = u - v^{(1)}$ und
$H^{(1)}(t) = \tfrac{1}{2}\|\nabla w(t)\|_{L^2}^2$. Dann gilt
\begin{equation}\label{eq:H1-ineq}
  \frac{d}{dt}H^{(1)}(t)
  \;\le\;
  -\frac{\nu}{2}\|\Delta u(t)\|_{L^2}^2
  \;+\; C_1^{(1)}\,H^{(1)}(t)
  \;+\; C_2^{(1)}
  \;+\; \text{\emph{Minimierer-Korrektur}},
\end{equation}
wobei
$C_1^{(1)} = C(\|\nabla v^{(1)}\|_{L^\infty},\,
\|\nabla^2 v^{(1)}\|_{L^2})$ und
$C_2^{(1)} = C(\|v^{(1)}\|_{H^2},\,\|f\|_{H^1})$
gleichmaessig beschraenkt sind nach Proposition~\ref{prop:attractor-regularity}
(Bootstrap-Regularitaet liefert $\sup_{v \in \mathcal{A}}\|v\|_{H^k}
< \infty$ fuer jedes $k$).
\end{proposition}

\begin{proof}[Beweisskizze]
Berechne $\frac{d}{dt}\frac{1}{2}\|\nabla w\|^2
= \langle \nabla w,\,\nabla(\partial_t u - \partial_t v^{(1)})\rangle$.
Der viskose Term liefert via partieller Integration
($\langle \nabla w, \nabla\Delta u\rangle
= -\langle \Delta w, \Delta u\rangle
= -\|\Delta u\|^2 + \langle \Delta v^{(1)}, \Delta u\rangle$):
\[
  \nu\langle \nabla w,\, \nabla\Delta u \rangle
  \;=\; -\nu\|\Delta u\|_{L^2}^2
        + \nu\langle \Delta v^{(1)},\, \Delta u \rangle.
\]
Nach der Young-Ungleichung erfuellt der Kreuzterm
$|\nu\langle \Delta v^{(1)}, \Delta u\rangle|
\le \frac{\nu}{4}\|\Delta u\|^2 + \nu\|\Delta v^{(1)}\|^2$,
was eine Netto-Dissipation von $-\frac{3\nu}{4}\|\Delta u\|^2$
plus den beschraenkten Attraktor-Regularitaetsterm
$\nu\|\Delta v^{(1)}\|^2 \le \nu\sup_{v\in\mathcal{A}}\|v\|_{H^2}^2$ ergibt.

Der nichtlineare Term $\langle \nabla w,\, \nabla[(u\cdot\nabla)u]\rangle$
wird mittels
$\|(u\cdot\nabla)u\|_{H^1} \le C\,\|u\|_{H^1}\,\|u\|_{H^2}$
(Produktregel und Sobolev-Einbettung $H^2(\mathbb{T}^3) \hookrightarrow
L^\infty(\mathbb{T}^3)$; vgl.\ \cite{ConstantinFoias1988}) abgeschaetzt, was
$|\langle \nabla w, \nabla[(u\cdot\nabla)u]\rangle|
\le C\,\|u\|_{H^1}\,\|u\|_{H^2}\,\|\nabla w\|$ ergibt.
Da $\|u\|_{H^2}^2 \le C(\|\Delta u\|^2 + \|u\|^2)$,
liefert die Young-Ungleichung mit Parameter $\varepsilon = \nu/4$:
$C\,\|u\|_{H^1}\,\|\Delta u\|\,\|\nabla w\|
\le \frac{\nu}{4}\|\Delta u\|^2
+ C_\varepsilon\,\|u\|_{H^1}^2\,\|\nabla w\|^2$.
Nach Absorption dieses $\frac{\nu}{4}\|\Delta u\|^2$ in das verbleibende
$\frac{3\nu}{4}\|\Delta u\|^2$-Budget betraegt die Netto-Dissipation
$-\frac{\nu}{2}\|\Delta u\|^2$ (konsistent mit dem Satz-Statement),
auf Kosten eines Gronwall-Faktors
$C_\varepsilon\,\|u\|_{H^1}^2 = C_\varepsilon(2H^{(1)}
+ \|v^{(1)}\|_{H^1}^2)$, der~$H^{(1)}$ multipliziert.
Dies ergibt eine Differentialungleichung vom \emph{Bihari-Typ}
(nichtlinearer Gronwall), da der Koeffizient $C_1^{(1)}$ von
$H^{(1)}$ selbst abhaengt:
\[
  \frac{d}{dt}H^{(1)}
  \;\le\;
  -\frac{\nu}{2}\|\Delta u\|^2
  + \alpha\,H^{(1)}
  + \beta\,(H^{(1)})^2
  + C_2^{(1)}
  + R^{(1)},
\]
wobei $\alpha, \beta$ nur von Attraktornormen abhaengen.
Der quadratische Term $(H^{(1)})^2$ entsteht, weil
$C_1^{(1)} \sim \|u\|_{H^1}^2 \sim H^{(1)} + \mathrm{const}$.
Das Standard-Gronwall-Lemma ist \emph{nicht} direkt anwendbar;
stattdessen liefert Biharis Ungleichung eine endliche Schranke fuer $H^{(1)}$,
sofern der Dissipationsterm $\frac{\nu}{2}\|\Delta u\|^2$
das quadratische Wachstum kontrolliert. Ob diese Kontrolle unter Blow-up
gilt, ist selbst Teil der offenen Frage. Die
Differentialungleichung~\eqref{eq:H1-diff} im Satz-Statement ist daher
auf jedem kompakten Teilintervall $[0,T'] \subset [0,T^*)$ gueltig, auf dem
$\|u\|_{H^1}$ endlich bleibt (d.h.\ im Regime starker Loesungen), was
genau das fuer den Widerspruchsbeweis relevante Regime ist.

Der Minimierer-Korrekturterm ist
$R^{(1)}(t) := |\langle \nabla w,\,
\nabla\partial_t v^{(1)}\rangle|
\le \|\nabla w\|_{L^2}\,\|\nabla\partial_t v^{(1)}\|_{L^2}
\le \sqrt{2H^{(1)}}\,L^{(1)}(t)$,
wobei $L^{(1)}(t) := \|\partial_t v^{(1)}(t)\|_{H^1}$ die
Regularitaetsschranke aus Assumption~$G^{(1)}$ ist.
Unter~$G^{(1)}$ gilt $L^{(1)} \in L^1_{\mathrm{loc}}$, also
$\int_0^T R^{(1)}\,dt
\le \sqrt{2\sup H^{(1)}}\,\int_0^T L^{(1)}\,dt < \infty$
auf jedem kompakten Teilintervall von~$[0,T^*)$.
\end{proof}

% ------------------------------------------------------------------
\subsection{Assumption $G^{(1)}$ und die Enstrophiedissipation}
% ------------------------------------------------------------------

\begin{definition}[Assumption $G^{(1)}$]\label{ass:G1}
Sei $u(t)$ eine starke Loesung auf $[0,T)$ mit globalem Attraktor
$\mathcal{A} \subset H^1$. Fuer f.u.\ $t \in [0,T)$ sei
$v^{(1)}(t) \in \arg\min_{v \in \mathcal{A}}
\|\nabla(u(t)-v)\|_{L^2}^2$.
Wir sagen, dass \textbf{Assumption~$G^{(1)}$} gilt, wenn:
\begin{enumerate}
  \item[\emph{($G^{(1)}.1$)}]
    \textbf{Gleichmaessige Attraktor-Regularitaet.}\;
    $\sup_{v \in \mathcal{A}}\|v\|_{H^2} < \infty$.
    (Dies folgt aus der bekannten Glattheit des Attraktors;
    vgl.\ Proposition~\ref{prop:attractor-regularity}.)
  \item[\emph{($G^{(1)}.2$)}]
    \textbf{Temporale Regularitaet der Projektion.}\;
    Die Abbildung $t \mapsto v^{(1)}(t)$ hat beschraenkte Variation in $H^1$
    auf jedem kompakten Teilintervall $[0,T'] \subset [0,T)$, und
    es existiert $L^{(1)} \in L^1_{\mathrm{loc}}(0,T)$, sodass
    \[
      \|\partial_t v^{(1)}(t)\|_{H^1}
      \;\le\; L^{(1)}(t)
      \qquad \text{fuer f.u.\ } t.
    \]
\end{enumerate}
Diese Annahme ist das exakte $H^1$-Analogon von Assumption~G im
$L^2$-Fall; sie betrifft ausschliesslich die temporale Stabilitaet der
Attraktorprojektion und legt der Navier--Stokes-Loesung selbst keine
zusaetzliche Regularitaet auf.
\end{definition}

\begin{remark}[Warum $G^{(1)}$ strukturell dasselbe ist wie~G]
Assumption~$G^{(1)}$ isoliert dieselbe geometrische Schwierigkeit wie
Assumption~G, aber auf der korrekten Sobolev-Ebene.
Sie hat exakt dieselbe Form wie~G
(Definition~\ref{ass:G}), von $L^2$ auf~$H^1$ gehoben. Die
Attraktor-Regularitaet
(Proposition~\ref{prop:attractor-regularity}) garantiert
$\sup_{v \in \mathcal{A}}\|v\|_{H^k} < \infty$ fuer alle~$k$,
sodass die \emph{raeumliche} Regularitaet von $v^{(1)}$ kein Problem ist.
Die Schwierigkeit ist, wie zuvor, ausschliesslich die \emph{zeitliche}
Regularitaet: ob der $H^1$-Minimierer springt oder sich glatt entwickelt.
Die AGC- und TLL/LDI-Rahmenwerke (Abschnitt~3) uebertragen sich direkt, mit
den Lipschitz-/BV-Bedingungen in~$H^1$ umformuliert.
$G^{(1)}$ ist staerker als~G (sie erfordert Kontrolle einer
zusaetzlichen Ableitung), aber nicht qualitativ verschieden.
\end{remark}

% ------------------------------------------------------------------
\subsection{Warum die $H^1$-Ebene die Dissipationsluecke vermeidet}
% ------------------------------------------------------------------

Die entscheidende Asymmetrie ist:

\begin{center}
\begin{tabular}{lll}
\hline
\textbf{Ebene} & \textbf{Dissipationsintegral} & \textbf{Unter Blow-up} \\
\hline
$L^2$\; (Energie)
  & $\int_0^{T^*}\|\nabla u\|_{L^2}^2\,dt$
  & \textbf{Endlich} (Leray-Identitaet) \\
$H^1$\; (Enstrophie)
  & $\int_0^{T^*}\|\Delta u\|_{L^2}^2\,dt$
  & \textbf{Keine a-priori-Schranke} \\
\hline
\end{tabular}
\end{center}

\noindent
Dies ist kein technisches Detail, sondern eine \emph{fundamentale Asymmetrie}
in der Navier--Stokes-Energiehierarchie: viskose Dissipation auf der
Energieebene ($\int\|\nabla u\|^2$, physikalisch: Enstrophieproduktion)
wird durch das anfaengliche Energiebudget kontrolliert, waehrend viskose Dissipation
auf der Enstrophieebene ($\int\|\Delta u\|^2$, physikalisch: Palinstrophie-Produktion) \emph{nicht} kontrolliert wird. Diese Asymmetrie ist der strukturelle Grund,
warum das $L^2$-Ebene-Argument keinen Widerspruch erzeugen kann, waehrend das
$H^1$-Ebene-Argument dies kann.

Die $H^1$-Energieidentitaet lautet
\begin{equation}\label{eq:H1-energy-id}
  \frac{1}{2}\|\nabla u(T)\|^2 + \nu\int_0^T\|\Delta u\|^2\,dt
  \;=\;
  \frac{1}{2}\|\nabla u_0\|^2
  + \int_0^T\!\bigl[\langle (u\cdot\nabla)u,\,\Delta u\rangle
  + \langle \nabla f,\, \nabla u\rangle\bigr]\,dt.
\end{equation}
Unter Blow-up gilt $\|\nabla u(T)\|^2 \to \infty$ fuer $T \nearrow T^*$.
Die linke Seite divergiert daher. Da der nichtlineare Term
$\int\langle (u\cdot\nabla)u, \Delta u\rangle$ ohne Schranke wachsen kann
(er involviert $\|u\|_{H^1}\,\|\Delta u\|$), ist die Bilanz
zwischen den drei Termen auf der rechten Seite \emph{nicht} a priori
kontrolliert. Insbesondere kann $\int_0^{T^*}\|\Delta u\|^2\,dt$
divergieren -- und falls dies eintritt, wird das $H^{(1)}$-Funktional durch
exakt dasselbe Budget-Argument wie in Satz~\ref{thm:main} unter Null getrieben,
aber nun mit einem \emph{tatsaechlich divergenten} Dissipationsterm.

% ------------------------------------------------------------------
\subsection{Bedingter $H^1$-Satz (Skizze)}
% ------------------------------------------------------------------

\begin{theorem}[Bedingte $H^1$-Obstruktion gegen Blow-up in endlicher Zeit]
\label{thm:H1-main}
Sei $u(t)$ eine starke Loesung der 3D-Navier--Stokes-Gleichungen auf
$[0,T^*)$ mit glatten Anfangsdaten und glattem Antrieb. Es gelte
Assumption~$G^{(1)}$ (Definition~\ref{ass:G1}). Definiere das
$H^1$-Abstandsfunktional
$H^{(1)}(u(t) \mid \mathcal{A})
:= \inf_{v \in \mathcal{A}} \tfrac{1}{2}\|\nabla(u(t)-v)\|_{L^2}^2$.

Dann gilt entlang der Loesung die Differentialungleichung
\begin{equation}\label{eq:H1-diff}
  \frac{d}{dt}H^{(1)}(u(t) \mid \mathcal{A})
  \;\le\;
  -\frac{\nu}{2}\|\Delta u(t)\|_{L^2}^2
  + C_1(t)\,H^{(1)}(u(t) \mid \mathcal{A})
  + C_2
  + R^{(1)}(t),
\end{equation}
wobei $C_1(t) = C(\|u(t)\|_{H^1}^2) = C(2H^{(1)}(t)
+ \|v^{(1)}\|_{H^1}^2)$ von der Loesungsnorm abhaengt (was die
Ungleichung zum Bihari-Typ macht), $C_2 < \infty$ nur von
Attraktornormen abhaengt,
und der Korrekturterm $R^{(1)}$ unter
Assumption~$G^{(1)}$ lokal integrierbar ist.

Falls zusaetzlich
\begin{equation}\label{eq:enstr-diss-div}
  \int_0^{T^*}\|\Delta u(t)\|_{L^2}^2\,dt = \infty,
\end{equation}
dann kann $H^{(1)}(u(t) \mid \mathcal{A})$ nicht fuer alle
$t < T^*$ nichtnegativ bleiben. Insbesondere ist ein Blow-up in endlicher Zeit
bei $T^* < \infty$ unter diesen Annahmen ausgeschlossen.
\end{theorem}

\begin{proof}[Beweisskizze]
Integriere~\eqref{eq:H1-diff} auf $[0,T]$ fuer $T < T^*$:
\[
  H^{(1)}(T)
  \;\le\;
  H^{(1)}(0)
  - \frac{\nu}{2}\int_0^T\|\Delta u\|^2\,dt
  + \int_0^T\!\bigl(C_1\, H^{(1)} + C_2\bigr)\,dt
  + \int_0^T R^{(1)}\,dt.
\]
Unter $G^{(1)}$ sind die Gronwall-Kosten und die Korrektur $\int R^{(1)}$
endlich auf jedem kompakten Teilintervall von~$[0,T^*)$.
Beachte, dass der Gronwall-Koeffizient $C_1$ von
$\|u\|_{H^1}^2 \sim H^{(1)} + \mathrm{const}$ abhaengt, was die
Ungleichung zum Bihari-Typ (nichtlinearer Gronwall) macht. Im
Widerspruchsregime ($u$ ist starke Loesung auf $[0,T^*)$) ist
$\|u\|_{H^1}$ fuer jedes $T < T^*$ endlich, also sind die integrierten Kosten
$\int_0^T(C_1 H^{(1)} + C_2)\,dt$ fuer jedes feste $T < T^*$ endlich
(wenn auch moeglicherweise wachsend fuer $T \nearrow T^*$).
Nach~\eqref{eq:enstr-diss-div} gilt
$\frac{\nu}{2}\int_0^T\|\Delta u\|^2\,dt \to \infty$ fuer $T \nearrow T^*$.
Da die Dissipation ohne Schranke waechst, waehrend die Gronwall-Kosten fuer jedes
$T < T^*$ endlich bleiben, existiert $T_0 < T^*$, sodass
$H^{(1)}(T_0) < 0$, im Widerspruch zu $H^{(1)} \ge 0$.
\end{proof}

\begin{remark}[Status und Vergleich mit dem $L^2$-Satz]
\label{rem:H1-comparison}
Satz~\ref{thm:H1-main} ersetzt die Dissipations-Dominanz-Bedingung~(D) von Satz~\ref{thm:main} durch die konkretere
Hypothese~\eqref{eq:enstr-diss-div} (divergente Enstrophiedissipation).
Die bedingte Struktur hat nun \emph{drei} Luecken:
\begin{enumerate}
  \item \textbf{Projektionsluecke:}\; Assumption~$G^{(1)}$
    (Regularitaet des $H^1$-Minimierers) -- selber Charakter wie
    Assumption~G, selbe Angriffsvektoren (AGC, TLL/LDI).
  \item \textbf{Enstrophie-Dissipationsluecke:}\;
    Ob~\eqref{eq:enstr-diss-div} unter BKM-Blow-up gilt.
    Anders als im $L^2$-Fall gibt es \emph{keine} a-priori-Schranke,
    die dies ausschliesst. Die Etablierung von~\eqref{eq:enstr-diss-div}
    erfordert Verstaendnis des nichtlinearen Transfers auf der
    Enstrophieebene (Wirbelstreckung vs.\ viskose Daempfung).
  \item \textbf{Bihari-Struktur:}\;
    Anders als im $L^2$-Satz (wo $C_1$ eine Konstante ist, die nur
    von Attraktornormen abhaengt) haengt der $H^1$-Gronwall-Koeffizient
    $C_1^{(1)}$ von $\|u\|_{H^1}^2 \sim H^{(1)}$ ab, was eine
    Bihari-artige (nichtlineare Gronwall) Ungleichung ergibt.
    Im Widerspruchsregime ist $H^{(1)}$ fuer jedes $T < T^*$ endlich
    (starke Loesung), also sind die integrierten Gronwall-Kosten endlich;
    aber die Wachstumsrate dieser Kosten nahe $T^*$ muss langsam genug
    relativ zur divergierenden Dissipation~\eqref{eq:enstr-diss-div} sein.
    Dies ist eine quantitative Verfeinerung von Luecke~(2) und ist automatisch
    erfuellt, wenn $\int\|\Delta u\|^2$ schneller als jedes Polynom divergiert.
\end{enumerate}
Auf $L^2$-Ebene ist der analoge Widerspruch unmoeglich, da die
Leray-Energieidentitaet stets
$\int_0^{T^*}\|\nabla u\|_{L^2}^2 < \infty$ erzwingt.
Auf $H^1$-Ebene existiert keine entsprechende a-priori-Schranke fuer
$\int\|\Delta u\|_{L^2}^2$; diese strukturelle Asymmetrie
macht den $H^1$-Lift prinzipiell durchfuehrbar.

Satz~\ref{thm:H1-main} ist doppelt bedingt: er isoliert
zwei unabhaengige offene Fragen -- die temporale Regularitaet der
$H^1$-Projektion ($G^{(1)}$) und die Divergenz der Enstrophiedissipation
unter Blow-up -- ohne eine davon vorauszusetzen.

\medskip\noindent
\textbf{Offene Probleme zur Schliessung der $H^1$-Ebene-Luecken:}
\begin{enumerate}
  \item $G^{(1)}$ aus Attraktorgeometrie beweisen (parallel zum
    $L^2$-AGC-Programm).
  \item \eqref{eq:enstr-diss-div} aus der Enstrophiebilanz
    und BKM-Blow-up beweisen. Eine hinreichende Bedingung waere: der Wirbelstreckungsterm $\int_\Omega(\omega\cdot\nabla)u\cdot\omega\,dx$
    waechst hoechstens wie $C\,\|\Delta u\|^{2-\delta}$ fuer ein
    $\delta > 0$ nahe~$T^*$.
  \item Alternativ: das TLL/LDI-Rahmenwerk auf~$H^1$ heben,
    was BV-Auswahlen fuer die $H^1$-Projektion liefert. Dies wuerde
    eine BV-Variante von~$G^{(1)}$ ohne inertiale Mannigfaltigkeiten ergeben.
\end{enumerate}
\end{remark}


\subsection*{Beweisstatus-Zusammenfassung}

Diese Arbeit ist ein \textbf{bedingtes Reduktions-Rahmenwerk}: Sie reduziert globale Regularitaet der 3D-Navier--Stokes-Gleichungen auf eine geometrische Eigenschaft des Attraktors (Selektionsregularitaet).

\begin{center}
\begin{tabular}{|l|l|}
\hline
\textbf{Komponente} & \textbf{Status} \\
\hline
\hline
\multicolumn{2}{|l|}{\textit{Bewiesen (unbedingt):}} \\
\hline
$H(u|\mathcal{A})$ Gronwall-Ungleichung & Theorem \\
Assumption G $\Rightarrow$ Regularitaet & Theorem \\
G' (BV) genuegt fuer G & Proposition \\
Squeezing $\Rightarrow$ TLL $\Rightarrow$ G & Proposition \\
\hline
\multicolumn{2}{|l|}{\textit{Numerisch bestaetigt:}} \\
\hline
BV-Selektion auf Lorenz/R\"ossler/Chen & 3/3 bestanden \\
Balanced-Viscosity-Konvergenz & Cauchy-Folge \\
Bedingung D (Dissipation beschraenkt) & 3/3 bestanden \\
\hline
\multicolumn{2}{|l|}{\textit{Offen (das Kernproblem):}} \\
\hline
Selektionsregularitaet (Assumption G) & Offen \\
Bestimmende Moden $\Rightarrow$ Squeezing & Offen \\
\hline
\end{tabular}
\end{center}


% ==================================================================
\section{Schlussfolgerung}
Durch die Formulierung des Regularitaetsproblems als Abstands-Dissipations-Bilanz
relativ zum globalen Attraktor des Navier--Stokes-Semiflusses etablieren wir
ein \emph{doppelt bedingtes} Regularitaetskriterium. Die beiden
Bedingungen sind:
\begin{enumerate}
  \item \textbf{Assumption~G} (oder das schwaechere~G$'$): die zeitabhaengige
    Attraktorprojektion hat hinreichende Regularitaet (AC oder BV).
  \item \textbf{Bedingung~(D)}: die kumulative viskose Dissipation
    dominiert die Gronwall-Kosten nahe einer hypothetischen Blow-up-Zeit
    (Bemerkung~\ref{rem:diss-gap}).
\end{enumerate}
Unter beiden Bedingungen fuehrt Blow-up in endlicher Zeit zu
$H(u \mid \mathcal{A}) < 0$ -- ein Widerspruch zu seiner Nichtnegativitaet.

\medskip\noindent
\textbf{Die Dissipationsluecke.}\;
Die Leray-Energieidentitaet garantiert, dass
$\int_0^{T^*}\|\nabla u\|_{L^2}^2\,dt < \infty$ selbst unter Blow-up
(Bemerkung~\ref{rem:divergence-landscape}). Was divergiert, ist die
\emph{punktweise} Enstrophie $\|\nabla u(t)\|_{L^2}^2 \to \infty$,
nicht ihr Zeitintegral. Ob diese punktweise Divergenz schnell
genug ist, um die Gronwall-Kosten zu uebertreffen, ist der Inhalt von
Bedingung~(D) -- eine strukturelle Frage ueber Blow-up-Raten, die
fuer die 3D-Navier--Stokes-Gleichungen offen ist.

\medskip\noindent
\textbf{Zwei parallele Angriffspfade fuer Assumption~G} sind identifiziert:
(i)~der AGC-Pfad, der~G auf positiven Reach des Attraktors reduziert
(aequivalent zur Existenz inertialer Mannigfaltigkeiten);
(ii)~der log-Lipschitz-Pfad (Lemma~\ref{lem:LL-BV}), der~G$'$ auf
Trajektorien-Ebene-Regularitaet (TLL) kombiniert mit Log-Abstand-Integrierbarkeit (LDI) reduziert. Der zweite Pfad umgeht die Barriere inertialer Mannigfaltigkeiten
vollstaendig.

\medskip\noindent
\textbf{Drei Strategien fuer Bedingung~(D)} sind in
Bemerkung~\ref{rem:diss-gap} skizziert: Heben des Funktionals auf $H^1$-Ebene,
Ausnutzung quantitativer Blow-up-Raten oder Verwendung Serrin-artiger Kriterien.
Schliessung entweder der Projektionsluecke (Assumption~G) \emph{oder} der
Dissipationsluecke (Bedingung~D) wuerde das vorliegende Rahmenwerk voranbringen,
und Schliessung beider wuerde einen unbedingten Regularitaetsbeweis liefern.

\medskip\noindent
\textbf{Das $H^1$-Ebene-Programm} (Abschnitt~\ref{sec:H1-lift})
bietet den vielversprechendsten Weg nach vorn: Durch Heben des
Attraktorabstandsfunktionals von $L^2$ auf~$H^1$ wird der
Dissipationsterm zu $\int\|\Delta u\|^2$, der -- anders als
$\int\|\nabla u\|^2$ -- \emph{keine} a-priori obere Schranke
unter Blow-up hat. Diese strukturelle Asymmetrie zwischen Energie- und
Enstrophiedissipation ist die zentrale Erkenntnis der vorliegenden Arbeit
und motiviert Satz~\ref{thm:H1-main} als die natuerliche
Fortsetzung dieses Programms.
Auf der $H^1$-Ebene tritt eine weitere Subtilitaet auf: Der
Gronwall-Koeffizient $C_1^{(1)}$ haengt von $\|u\|_{H^1}^2$ ab,
was eine Bihari-artige (nichtlineare Gronwall) Differentialungleichung
ergibt, deren integrierte Kosten wachsen, wenn sich die Loesung dem
hypothetischen Blow-up naehert. Der Widerspruchsbeweis bleibt gueltig,
sofern die divergente Enstrophiedissipation dieses Wachstum uebertrifft --
eine quantitative Verfeinerung, die das $H^1$-Programm mit praezisen
Blow-up-Raten-Abschaetzungen verbindet.

% ===================================================================
% AI Disclosure
% ===================================================================
\subsection*{KI-Offenlegung}
Bei der Erstellung dieses Manuskripts wurden KI-gestuetzte Werkzeuge (Claude, Anthropic)
in unterstuetzender Funktion eingesetzt, insbesondere fuer Literaturrecherche,
Textstrukturierung und Formulierungshilfe. Der konzeptionelle Entwurf,
die wissenschaftliche Argumentation und die Verantwortung fuer den gesamten Inhalt
liegen ausschliesslich beim Autor.

\begin{thebibliography}{99}

\bibitem{BealeKatoMajda1984}
J.~T. Beale, T.~Kato, and A.~Majda,
\emph{Remarks on the breakdown of smooth solutions for the 3-D Euler equations},
Comm.\ Math.\ Phys.\ \textbf{94} (1984), no.~1, 61--66.

\bibitem{BertozziMajda2002}
A.~L. Bertozzi and A.~J. Majda,
\emph{Vorticity and Incompressible Flow},
Cambridge Texts in Applied Mathematics, Cambridge University Press, 2002.

\bibitem{CaffarelliKohnNirenberg1982}
L.~Caffarelli, R.~Kohn, and L.~Nirenberg,
\emph{Partial regularity of suitable weak solutions of the Navier--Stokes equations},
Comm.\ Pure Appl.\ Math.\ \textbf{35} (1982), no.~6, 771--831.

\bibitem{ConstantinFoias1988}
P.~Constantin and C.~Foias,
\emph{Navier--Stokes Equations},
Chicago Lectures in Mathematics, University of Chicago Press, 1988.

\bibitem{FoiasManleyRosaTemam2001}
C.~Foias, O.~Manley, R.~Rosa, and R.~Temam,
\emph{Navier--Stokes Equations and Turbulence},
Encyclopedia of Mathematics and its Applications, vol.~83, Cambridge University Press, 2001.

\bibitem{LasryLions2007}
J.-M. Lasry and P.-L. Lions,
\emph{Mean field games},
Jpn.\ J.\ Math.\ \textbf{2} (2007), no.~1, 229--260.

\bibitem{Leray1934}
J.~Leray,
\emph{Sur le mouvement d'un liquide visqueux emplissant l'espace},
Acta Math.\ \textbf{63} (1934), 193--248.

\bibitem{Temam2001}
R.~Temam,
\emph{Navier--Stokes Equations: Theory and Numerical Analysis},
AMS Chelsea Publishing, Providence, RI, 2001 (reprint of the 1984 edition).

\bibitem{Wagner1977}
D.~H. Wagner,
\emph{Survey of measurable selection theorems},
SIAM J.\ Control Optim.\ \textbf{15} (1977), no.~5, 859--903.

\bibitem{KuksinShirikyan2012}
S.~Kuksin and A.~Shirikyan,
\emph{Mathematics of Two-Dimensional Turbulence},
Cambridge Tracts in Mathematics, vol.~194, Cambridge University Press, 2012.

\bibitem{HairerMattingly2006}
M.~Hairer and J.~C. Mattingly,
\emph{Ergodicity of the 2D Navier--Stokes equations with degenerate stochastic forcing},
Ann.\ of Math.\ \textbf{164} (2006), no.~3, 993--1032.

\bibitem{MalletParetSell1988}
J.~Mallet-Paret and G.~R. Sell,
\emph{Inertial manifolds for reaction diffusion equations in higher space dimensions},
J.~Amer.\ Math.\ Soc.\ \textbf{1} (1988), no.~4, 805--866.

\bibitem{AmbrosioFuscoPallara2000}
L.~Ambrosio, N.~Fusco, and D.~Pallara,
\emph{Functions of Bounded Variation and Free Discontinuity Problems},
Oxford Mathematical Monographs, Clarendon Press, Oxford, 2000.

\bibitem{BarbuCannone2016}
V.~Barbu and M.~Cannone,
\emph{Log-Lipschitz regularity of the Navier--Stokes semiflow and applications},
J.~Math.\ Pures Appl.\ \textbf{106} (2016), no.~3, 511--529.

\bibitem{Zelik2022}
S.~Zelik,
\emph{Attractors of dissipative partial differential equations},
in: Handbook of Differential Equations: Evolutionary Equations, vol.~4, Elsevier, 2022.

\bibitem{Zelik2014}
S.~Zelik,
\emph{Inertial manifolds and finite-dimensional reduction for
dissipative PDEs},
Proc.\ Roy.\ Soc.\ Edinburgh Sect.~A \textbf{144} (2014), no.~6,
1245--1327.

\bibitem{EdenFoiasNicolaenko1994}
A.~Eden, C.~Foias, B.~Nicolaenko, and R.~Temam,
\emph{Exponential Attractors for Dissipative Evolution Equations},
Wiley, 1994.

\bibitem{FST-TU2026}
L.~Geiger, \emph{Anomalous Dissipation and the Turbulent Cascade:
A Variational Free-Energy Derivation of Kolmogorov Scaling},
Preprint, 2026.

\bibitem{GeigerRH2026c}
L.~Geiger, \emph{The Riemann Hypothesis via Free-Energy Spectral Theory (Parts~I--III)},
Zenodo, March 2026.
\href{https://doi.org/10.5281/zenodo.19035845}{doi:10.5281/zenodo.19035845}.

\bibitem{FST-Unified2026}
L.~Geiger, \emph{The Renormalized Free Energy Framework: A Unified Resolution of Structural Bottlenecks},
Zenodo, March 2026.
\href{https://doi.org/10.5281/zenodo.19036191}{doi:10.5281/zenodo.19036191}.

\bibitem{GrujicXu2025}
Z.~Grujic und L.~Xu,
\emph{Asymptotic criticality of the Navier--Stokes regularity problem},
Ann.\ PDE \textbf{11} (2025), Article~9.

\bibitem{FoiasSellTemam88}
C.~Foias, G.\,R.~Sell und R.~Temam,
\emph{Inertial manifolds for nonlinear evolutionary equations},
J.~Differential Equations \textbf{73} (1988), 309--353.

\bibitem{Moffatt69}
H.\,K.~Moffatt,
\emph{The degree of knottedness of tangled vortex lines},
J.~Fluid Mech.\ \textbf{35} (1969), 117--129.

\bibitem{ArnoldKhesin98}
V.\,I.~Arnold und B.\,A.~Khesin,
\emph{Topological Methods in Hydrodynamics},
Applied Mathematical Sciences, Bd.~125,
Springer, New York, 1998.

\bibitem{DrivasNguyen2019}
T.\,D.~Drivas und H.\,Q.~Nguyen,
\emph{Remarks on the emergence of weak Euler solutions in the vanishing
viscosity limit},
J.~Nonlinear Sci.\ \textbf{29} (2019), 1127--1161.
\href{https://arxiv.org/abs/1808.01014}{arXiv:1808.01014}.

\bibitem{DeRosaDrivasInversi2024}
L.~De~Rosa, T.\,D.~Drivas und M.~Inversi,
\emph{On the support of anomalous dissipation measures},
J.~Math.\ Fluid Mech.\ \textbf{26} (2024), Art.~62.
\href{https://arxiv.org/abs/2301.09603}{arXiv:2301.09603}.

\bibitem{DeRosaDrivasInversiIsett2025}
L.~De~Rosa, T.\,D.~Drivas, M.~Inversi und P.~Isett,
\emph{Intermittency and dissipation regularity in turbulence},
Preprint, 2025.
\href{https://arxiv.org/abs/2502.10032}{arXiv:2502.10032}.

\bibitem{Sturm2025}
K.-T.~Sturm, \emph{Bakry--\'Emery, Hardy, and spectral gap estimates on manifolds with conical singularities}, Calc.\ Var.\ Partial Differential Equations (2025). \texttt{arXiv:2405.10734}.

\bibitem{DondlFrenzelMielke2019}
P.~Dondl, T.~Frenzel und A.~Mielke,
\emph{A gradient system with a wiggly energy and relaxed EDP-convergence},
ESAIM Control Optim.\ Calc.\ Var.\ \textbf{25} (2019), Art.~68.
\href{https://doi.org/10.1051/cocv/2018058}{doi:10.1051/cocv/2018058}.

\bibitem{MielkeRossiSavare2016}
A.~Mielke, R.~Rossi und G.~Savar\'e,
\emph{Balanced-viscosity (BV) solutions to infinite-dimensional rate-independent systems},
J.~Eur.\ Math.\ Soc.\ \textbf{18} (2016), Nr.~9, 2107--2165.
\href{https://doi.org/10.4171/JEMS/639}{doi:10.4171/JEMS/639}.

\bibitem{KevrekidisTiti2024}
I.~G.~Kevrekidis und E.~S.~Titi,
\emph{Data-driven approximate inertial manifolds for dissipative PDEs via diffusion maps},
SIAM J.\ Appl.\ Dyn.\ Syst.\ \textbf{23} (2024), Nr.~2, 1017--1047.
\href{https://doi.org/10.1137/23M1548534}{doi:10.1137/23M1548534}.

\bibitem{Geiger2026-DE}
L.~Geiger,
\emph{Dark Energy as Residual Vacuum Free Energy: A Thermodynamic Bound on the Cosmological Constant},
Preprint (2026). \texttt{zenodo.org/record/geiger2026de}.

\bibitem{Geiger2026-YM}
L.~Geiger,
\emph{Volume-Independent Spectral Gap for Strong-Coupling Lattice Gauge Theory via Poincar\'e--Dobrushin Machinery},
Preprint (2026). \texttt{zenodo.org/record/geiger2026ym}.

\bibitem{Geiger2026-BSD}
L.~Geiger,
\emph{The BSD Conjecture as a Positivity Normal-Form Theorem: Height Saturation and the Impossibility of Rank-Order Discrepancy},
Preprint (2026). \texttt{zenodo.org/record/geiger2026bsd}.

\bibitem{ChepyzhovVishik2002}
V.~V.~Chepyzhov and M.~I.~Vishik, \textit{Attractors for Equations of Mathematical Physics}, Amer.\ Math.\ Soc., Providence, 2002.

\bibitem{Simon1987}
J.~Simon, \textit{Compact sets in the space $L^p(0,T;B)$}, Ann.\ Mat.\ Pura Appl., \textbf{146} (1987), 65--96.

\bibitem{Artstein1995}
Z.~Artstein, \textit{A calculus for set-valued maps and set-valued evolution equations}, Set-Valued Anal., \textbf{3} (1995), 213--261.

\bibitem{PoliquinRockafellarThibault2000}
R.~A.~Poliquin, R.~T.~Rockafellar, and L.~Thibault, \textit{Local differentiability of distance functions}, Trans.\ Amer.\ Math.\ Soc., \textbf{352} (2000), 5231--5249.

\bibitem{Geiger2026-LDI} L.~Geiger, \textit{Log-Distanz-Integrierbarkeit fuer dissipative Stroemungen: BV-Projektionen auf fraktale Attraktoren}, Preprint, 2026. \texttt{DOI: 10.5281/zenodo.19056808}.

\end{thebibliography}

\end{document}
